%---------------------------------------------------------------------------------------------------------------------------- 
\graphicspath{{Parti/P03-Meccanica_SR/C05-Problemi_CinStat_SCR/Figure/}}
\chapter{Congruenza, equilibrio e lavoro virtuale nei sistemi di corpi rigidi}\label{cap:SCR}
%----------------------------------------------------------------------------

\begin{sommario}
	Si estende ai sistemi di corpi rigidi piani la duplice lettura,
	vettoriale ed energetica, già introdotta per il singolo corpo
	(\CAPref{cap:cinematica_infinitesima}, \CAPref{cap:statica_CR}).
	Si definisce la \emph{congruenza} di una configurazione del sistema
	come compatibilità simultanea con la rigidità di ciascun corpo e con
	i vincoli, esterni e interni, che li collegano; si mostra che
	l'\emph{equilibrio} del sistema — l'equilibrio di ciascun corpo sotto
	le forze attive e le reazioni trasmesse dai corpi vicini — condivide
	con la congruenza la medesima geometria di vincolo, letta in due
	direzioni opposte. Si deriva su questa base il dualismo fra la matrice dei vincoli cinematici e quella
	di equilibrio; se ne ricava poi, come conseguenza e non come ipotesi,
	l'identità bilineare del lavoro virtuale fra stati indipendenti. Da
	questa identità discendono due famiglie speculari di risultati: le
	condizioni di Compatibilità statica e Compatibilità cinematica, che
	stabiliscono quando i due problemi ammettono soluzione, e il teorema
	degli spostamenti virtuali (principio di Müller-Breslau) e il teorema
	delle forze virtuali, che ne calcolano
	esplicitamente singole componenti senza risolvere il sistema completo.
	Si introducono quindi i quattro sottospazi fondamentali di $\bm A$ —
	meccanismi, forze equilibrabili, stati di autoequilibrio, cedimenti
	ammissibili — e se ne ricava, come conseguenza diretta del teorema
	rango-nullità, la classificazione labile/isostatico/iperstatico. Il
	capitolo si chiude con una serie di esempi applicativi, che
	ripercorrono dall'inizio alla fine, su un'unica struttura, tutti gli
	strumenti introdotti.
\end{sommario}


%----------------------------------------------------------------------------
\section{Perché una teoria dei sistemi}\label{sec:motivazione_SCR}
%----------------------------------------------------------------------------

La teoria costruita nei due capitoli precedenti è completa per un corpo
rigido isolato: la sua configurazione è descritta da tre soli parametri,
$\bigl(u(O),\,v(O),\,\vartheta\bigr)$, e il suo equilibrio da due sole
condizioni vettoriali, $\rb=\bm{0}$ e $\bm{\upmu}(O)=\bm{0}$, equivalenti
--- lo si è dimostrato --- all'annullarsi del lavoro virtuale per ogni
moto rigido. Nulla di tutto questo cambia se, invece di un corpo isolato,
se ne considerano molti, purché restino indipendenti l'uno dall'altro: si
applica la teoria del corpo singolo a ciascuno, separatamente, e non c'è
altro da dire.

Le cose cambiano quando i corpi vengono \emph{collegati fra loro}. Si
consideri il sistema di \FIGref{fig:gerber_motiv}: due travi rigide
$\mathcal{C}_1$ e $\mathcal{C}_2$, la prima incernierata al suolo in $A$
e appoggiata su un carrello verticale in $B$, la seconda appoggiata su un
carrello verticale in $D$; le due travi sono collegate fra loro da una
cerniera in $C$. Si tratta della cosiddetta \emph{trave Gerber}, che
riprenderemo per l'intero capitolo. In questo contesto la parola ``trave'' denota semplicemente un corpo rigido di forma allungata: la scelta è solo per comodità di rappresentazione grafica 
e per anticipare  esempi strutturali rilevanti; nulla nella trattazione presente dipende dalla forma del corpo;  le proprietà specifiche della trave come elemento strutturale saranno introdotte nel \CAPref{cap:Trave_Rigida}.

%----------------------------------------------------------------------------------------------------------------------
\begin{figure}[htbp]
	\centering
	\begin{tikzpicture}[>=Stealth, line width=0.8pt, font=\small,
		nodo/.style={circle, fill=white, draw=black,
			inner sep=0pt, minimum size=5pt, line width=0.8pt},
		cint/.style={circle, fill=white, draw=black,
			inner sep=0pt, minimum size=6pt, line width=0.8pt},
		richiam/.style={thin, gray, dashed},
		tacca/.style={thin, gray!70}]
		
		\def\a{1.0}
		\def\tw{0.22}\def\thA{0.38}\def\thR{0.32}
		\def\bw{0.16}\def\gw{0.30}\def\sr{0.08}
		
		\draw[line width=2.2pt] (0,0) -- ({3*\a},0);
		\draw[line width=2.2pt] ({3*\a},0) -- ({5*\a},0);
		
		\node[above=2pt] at (0,0)       {$A$};
		\node[above=2pt] at ({2*\a},0)  {$B$};
		\node[above=2pt] at ({3*\a},0)  {$C$};
		\node[above=2pt] at ({5*\a},0)  {$D$};
		
		\node at ({\a},  -0.32) {$\mathcal{C}_1$};
		\node at ({4*\a},-0.32) {$\mathcal{C}_2$};
		
		%% cerniera in A
		\draw (0,0) -- ++(-\tw,-\thA) -- ++(2*\tw,0) -- cycle;
		\draw (-\gw,-\thA) -- (\gw,-\thA);
		\foreach \x in {-0.24,-0.12,0,0.12,0.24}
		\draw[thin] (\x,-\thA) -- ++(-0.09,-0.11);
		\node[nodo] at (0,0) {};
		
		%% carrello intermedio in B
		\draw ({2*\a},0) -- ++(-\tw,-\thR) -- ++(2*\tw,0) -- cycle;
		\draw ({2*\a-\gw},-\thR) -- ({2*\a+\gw},-\thR);
		\draw ({2*\a-\bw},-\thR-0.12) circle (2.4pt);
		\draw ({2*\a+\bw},-\thR-0.12) circle (2.4pt);
		\draw ({2*\a-\gw},-\thR-0.24) -- ({2*\a+\gw},-\thR-0.24);
		\foreach \x in {-0.24,-0.12,0,0.12,0.24}
		\draw[thin] ({2*\a+\x},-\thR-0.24) -- ++(-0.09,-0.11);
		\fill[white] ({2*\a-\sr},0) arc(180:360:\sr) -- cycle;
		\draw ({2*\a-\sr},0) arc(180:360:\sr);
		
		%% cerniera interna in C
		\node[cint] at ({3*\a},0) {};
		
		%% carrello terminale in D
		\draw ({5*\a},0) -- ++(-\tw,-\thR) -- ++(2*\tw,0) -- cycle;
		\draw ({5*\a-\gw},-\thR) -- ({5*\a+\gw},-\thR);
		\draw ({5*\a-\bw},-\thR-0.12) circle (2.4pt);
		\draw ({5*\a+\bw},-\thR-0.12) circle (2.4pt);
		\draw ({5*\a-\gw},-\thR-0.24) -- ({5*\a+\gw},-\thR-0.24);
		\foreach \x in {-0.24,-0.12,0,0.12,0.24}
		\draw[thin] ({5*\a+\x},-\thR-0.24) -- ++(-0.09,-0.11);
		\node[nodo] at ({5*\a},0) {};
		
		%% quotatura
		\draw[richiam] (0,      -\thA-0.15) -- (0,      -1.55);
		\draw[richiam] ({2*\a}, -1.40)      -- ({2*\a}, -1.55);
		\draw[richiam] ({3*\a}, -0.10)      -- ({3*\a}, -1.55);
		\draw[richiam] ({5*\a}, -1.40)      -- ({5*\a}, -1.55);
		
		\draw[tacca] (0,     -1.55) -- (0,     -1.75);
		\draw[tacca] ({2*\a},-1.55) -- ({2*\a},-1.75);
		\draw[tacca] ({3*\a},-1.55) -- ({3*\a},-1.75);
		\draw[tacca] ({5*\a},-1.55) -- ({5*\a},-1.75);
		
		\draw[<->, thin, gray] (0,     -1.65) -- ({2*\a},-1.65);
		\node[above=0.04] at ({\a},   -1.65) {$2a$};
		\draw[<->, thin, gray] ({2*\a},-1.65) -- ({3*\a},-1.65);
		\node[above=0.04] at ({2.5*\a},-1.65) {$a$};
		\draw[<->, thin, gray] ({3*\a},-1.65) -- ({5*\a},-1.65);
		\node[above=0.04] at ({4*\a}, -1.65) {$2a$};
		
	\end{tikzpicture}
	\caption{La trave Gerber: due corpi rigidi $\mathcal{C}_1$, $\mathcal{C}_2$
		connessi da una cerniera in $C$; cerniera al suolo in $A$, carrelli
		verticali in $B$ e $D$.}
	\label{fig:gerber_motiv}
\end{figure}
%----------------------------------------------------------------------------------------------------------------------

Si dunque provi a ragionare sulla  cinematica della trave Gerber senza scrivere ancora
equazioni. Se il carrello in $B$ venisse rimosso, il sistema potrebbe
muoversi: $\mathcal{C}_1$ ruoterebbe rigidamente attorno ad $A$,
trascinando con sé il punto $C$ che, come punto di $\mathcal{C}_1$, è
vincolato dalla rigidità di $\mathcal{C}_1$ a seguirlo secondo quanto prescritto
dalla cinematica del corpo rigido; ma la
cerniera in $C$ non impone nulla sulla \emph{rotazione} di $\mathcal{C}_2$,
che è quindi libero di ruotare per conto proprio,
purché il vincolo prescritto dal carrello in $D$ resti soddisfatto. 

Il sistema, insomma, si
muoverebbe pur restando rigido pezzo per pezzo. Che un sistema vincolato
possa muoversi non è di per sé una novità: anche un singolo corpo rigido
è labile se i vincoli esterni non bastano a immobilizzarlo. Ma in quel caso si tratta sempre di un moto
rigido dell'intero corpo. Qui è diverso: il moto appena descritto non è un moto rigido del sistema
preso come un blocco unico --- $\mathcal{C}_1$ e $\mathcal{C}_2$ ruotano
di angoli diversi, $\vartheta_1\neq\vartheta_2$ in generale --- è un moto
che i due corpi compiono l'uno rispetto all'altro, pur restando ciascuno,
separatamente, perfettamente rigido. È questa origine \emph{interna}
della labilità, resa possibile dallo svincolo in $C$ e non riconducibile
a un difetto di vincoli esterni, il fenomeno che la teoria del corpo
singolo non conosce e non può conoscere.

C'è un secondo fenomeno, duale del primo, che vale la pena solo annunciare
qui, rimandando a un esempio concreto quando classificheremo i sistemi: riguarda non la possibilità
di muoversi, ma quella delle reazioni vincolari di autoequilibrarsi. Se i
vincoli che tengono insieme un sistema sono più numerosi di quanto
l'equilibrio richieda, le equazioni di equilibrio da sole non bastano più
a escludere che le reazioni assumano valori non nulli anche in assenza di
ogni carico esterno. È il fenomeno duale della labilità: là mancavano
equazioni di vincolo e restava un moto libero; qui ne avanzano, e resta
un insieme di reazioni libere di autoequilibrarsi.

Si osservi che nulla di ciò  che abbiamo detto per la trave Gerber dipende dal fatto che i corpi siano
\emph{molti}: dipende dal fatto che sono \emph{parzialmente} collegati.
Se in $C$ ci fosse un incastro anziché una cerniera --- se cioè anche la
rotazione relativa fosse impedita --- $\mathcal{C}_1$ e $\mathcal{C}_2$
si comporterebbero, cinematicamente, come un unico corpo rigido, e la
teoria sviluppata finora basterebbe, senza bisogno di nulla di
nuovo: un vincolo interno che blocca tutti i gradi di libertà relativi
non introduce alcuna ricchezza cinematica. È precisamente la possibilità
di uno \emph{svincolo} --- una cerniera che libera la rotazione relativa,
un carrello interno che libera anche una componente di traslazione ---
a rendere necessaria una teoria dei sistemi propriamente detta: una
teoria che sappia trattare, in un unico quadro, corpi che sono a tratti
un blocco unico e a tratti liberi di muoversi l'uno rispetto all'altro.

Il compito di questo capitolo è costruire questo quadro, mantenendo per
quanto possibile lo stesso impianto concettuale già utilizzato per il
corpo singolo.


%----------------------------------------------------------------------------
\section{La congruenza del sistema}\label{sec:congruenza_SCR}
%----------------------------------------------------------------------------

Cominciamo dalla cinematica. Che cosa significa, per il sistema della
\FIGref{fig:gerber_motiv}, che una sua configurazione è ammissibile ---
cioè compatibile sia con la natura rigida di ciascuna trave sia con i
vincoli che le legano fra loro e al suolo, considerando eventuali cedimenti? La risposta si articola in due
richieste indipendenti, di natura diversa, che vale la pena tenere
distinte prima di fonderle in un unico oggetto.

La prima richiesta riguarda ciascun corpo preso singolarmente, ed è
interamente nota dal \CAPref{cap:cinematica_infinitesima}: lo spostamento
di $\mathcal{C}_1$ dev'essere quello di un corpo rigido, dunque
completamente determinato, una volta scelto un polo $O_1$ solidale a
$\mathcal{C}_1$, dai tre parametri $\bigl(u(O_1),\,v(O_1),\,\vartheta_1\bigr)$
tramite la formula fondamentale della cinematica infinitesima; lo stesso,
con un proprio polo $O_2$ e propri parametri $\bigl(u(O_2),\,v(O_2),\,
\vartheta_2\bigr)$, vale per $\mathcal{C}_2$. Nessuna relazione lega, a
questo stadio, i parametri di un corpo a quelli dell'altro.

La seconda richiesta riguarda invece i vincoli, ed è qui che conviene
ragionare direttamente sulla Gerber prima di enunciare qualunque regola
generale. Scegliamo, per fissare le idee, il polo $O_1\equiv A$ per
$\mathcal{C}_1$ e il polo $O_2\equiv C$ per $\mathcal{C}_2$ --- una scelta
comoda ma non obbligata, come già osservato per il corpo singolo: il
polo di un corpo è arbitrario, e conviene sceglierlo dove semplifica i
conti, per esempio in corrispondenza di un vincolo. Supponiamo inoltre,
per fissare le distanze, che $B$ disti $2a$ da $A$, $C$ disti $3a$ da
$A$ e $D$ disti $2a$ da $C$, tutti lungo l'asse delle travi.

\paragraph{Cerniera in $A$.} Impone che il punto $A$, che coincide
con il polo di $\mathcal{C}_1$, resti fisso: impone dunque direttamente le
componenti di traslazione del polo, senza coinvolgere la rotazione:
\begin{equation}\label{eq:vinc_A_gerber}
	u(O_1) = 0, \qquad v(O_1) = 0.
\end{equation}

\paragraph{Carrello in $B$.} Impone che la componente verticale dello
spostamento del punto $B$ --- un punto di $\mathcal{C}_1$ diverso dal
polo --- sia assegnata. Poiché $B$ non è il polo, il suo spostamento
verticale non è semplicemente $v(O_1)$: applicando a $\mathcal{C}_1$ la
formula fondamentale della cinematica infinitesima
(\EQUref{eq:FFCI_piano_deriv}) si ottiene $v(O_1)+2a\,\vartheta_1$, e
dunque
\begin{equation}\label{eq:vinc_B_gerber}
	v(O_1) + 2a\,\vartheta_1 =- 2\bar{v},
\end{equation}
dove $2\bar{v}$ è l'entità dello spostamento verticale imposto al punto $B$.
È così che, per un vincolo non applicato nel polo, la rotazione
$\vartheta_1$ entra nell'equazione insieme alle componenti di
traslazione.

\paragraph{Cerniera in $C$.} È il vincolo davvero nuovo, quello che lega
fra loro i due corpi: impone che lo spostamento del punto $C$, considerato
come punto di $\mathcal{C}_1$, coincida con lo spostamento dello stesso
punto $C$, considerato come punto di $\mathcal{C}_2$. Avendo scelto
$O_2\equiv C$, lo spostamento di $C$ come punto di $\mathcal{C}_2$ è
semplicemente $\bigl(u(O_2),\,v(O_2)\bigr)$; come punto di
$\mathcal{C}_1$, applicando ancora la formula fondamentale, è
$\bigl(u(O_1),\,v(O_1)+3a\,\vartheta_1\bigr)$. Uguagliando le due
espressioni si ottengono due equazioni che legano \emph{per la prima
	volta} i parametri di $\mathcal{C}_1$ a quelli di $\mathcal{C}_2$:
\begin{equation}\label{eq:vinc_C_gerber}
	u(O_1) = u(O_2), \qquad
	v(O_1) + 3a\,\vartheta_1 = v(O_2).
\end{equation}
Si noti bene cosa \emph{non} compare in \eqref{eq:vinc_C_gerber}:
nessuna equazione lega $\vartheta_1$ a $\vartheta_2$. È esattamente lo
svincolo rotazionale della cerniera, già discusso in astratto nel
\PARref{sec:motivazione_SCR}, a mancare qui: $\mathcal{C}_2$ resta libero
di ruotare indipendentemente da $\mathcal{C}_1$, e questa libertà è,
letteralmente, l'assenza di un'equazione.

\paragraph{Carrello in $D$.} Impone, in tutto analogamente al carrello
in $B$ ma questa volta nei parametri di $\mathcal{C}_2$,
\begin{equation}\label{eq:vinc_D_gerber}
	v(O_2) + 2a\,\vartheta_2 =- \bar{v}.
\end{equation}

Qui abbiamo supposto che il punto $D$ si sposti della quantità nota $-\bar{v}$.

Le equazioni \eqref{eq:vinc_A_gerber}--\eqref{eq:vinc_D_gerber} sono in
tutto sei relazioni scalari (due dalla cerniera in $A$, una dal carrello
in $B$, due dalla cerniera in $C$, una dal carrello in $D$) fra i sei
parametri $\bigl(u(O_1),v(O_1),\vartheta_1,u(O_2),v(O_2),\vartheta_2
\bigr)$: esattamente tanti quanti servono, in questo caso, a determinare
univocamente la configurazione. Che questo bilancio --- vincoli contro
parametri --- sia sempre risolvente, o possa fallire in un modo o
nell'altro, è precisamente il tipo di domanda a cui la teoria che stiamo
costruendo dovrà saper rispondere in generale, non caso per caso: da qui
la necessità di un formalismo.

\paragraph{Il vettore di configurazione.}
Il primo passo del formalismo è solo un atto di notazione: raccogliere in un unico vettore
colonna i parametri cinematici di tutti i corpi del sistema.


	Dato un sistema di $n_c$ corpi rigidi piani $\mathcal{C}_1,\dots,
	\mathcal{C}_{n_c}$, scelto per ciascuno un polo $O_i$ solidale, si
	chiama \emph{vettore di configurazione}\index{vettore di configurazione|textbf}
	del sistema il vettore
	\begin{equation}\label{eq:vettore_u_SCR}
		\bm{u} \coloneqq
		\bigl\{
		u_1\; v_1\; \vartheta_1 \;\;
		u_2\; v_2\; \vartheta_2 \;\;
		\cdots \;\;
		u_{n_c}\; v_{n_c}\; \vartheta_{n_c}
		\bigr\}^{\!\top}
		\in \mathbb{R}^{n}, \qquad n\coloneqq 3n_c,
	\end{equation}
	dove $\bigl(u_i,\,v_i,\,\vartheta_i\bigr)\coloneqq\bigl(u(O_i),\,
	v(O_i),\,\vartheta_i\bigr)$ sono i parametri cinematici di
	$\mathcal{C}_i$ rispetto al polo scelto.


Non c'è, in questa definizione, alcuna informazione che non fosse già
nella cinematica del corpo singolo applicata $n_c$ volte: $\bm{u}$ è
soltanto l'elenco ordinato di ciò che già sapevamo descrivere. Vediamo
subito come si costruisce, e a cosa serve, su un caso concreto.
%
%
%\subsection{Un caso esplicito: la trave di Gerber}\label{sec:gerber_congruenza}
%
%Riprendiamo il sistema di \FIGref{fig:gerber_motiv}, scegliendo --- per
%fissare le idee, non per obbligo --- il polo $O_1\equiv A$ per
%$\mathcal{C}_1$ e il polo $O_2\equiv C$ per $\mathcal{C}_2$, con $B$ a
%distanza $2a$ da $A$, $C$ a distanza $3a$ da $A$ e $D$ a distanza $2a$
%da $C$, tutti lungo l'asse delle travi. Assegniamo un cedimento
%verticale $2\bar v$ verso il basso al carrello in $B$ e un cedimento
%verticale $\bar v$ verso il basso al carrello in $D$; la cerniera in
%$A$ e la cerniera interna in $C$ restano perfette. Il vettore di
%configurazione, per $n_c=2$, $n=6$, è
%\begin{equation}\label{eq:u_gerber_cong}
%	\bm{u} = \bigl\{u_1,\,v_1,\,\vartheta_1,\,u_2,\,v_2,\,
%	\vartheta_2\bigr\}^{\!\top}\in\mathbb{R}^6,
%\end{equation}
%avendo scritto $u_1\coloneqq u(O_1)$ e così via.
%
%\paragraph{Cerniera in $A$.} Impone che il punto $A$, che coincide con
%il polo di $\mathcal{C}_1$, resti fisso:
%\begin{equation}\label{eq:vinc_A_gerber}
%	u_1 = 0, \qquad v_1 = 0.
%\end{equation}
%
%\paragraph{Carrello in $B$.} Impone che lo spostamento verticale del
%punto $B$ valga $-2\bar v$. Poiché $B$ non è il polo, il suo spostamento
%verticale non è semplicemente $v_1$: applicando a $\mathcal{C}_1$ la
%formula fondamentale della cinematica infinitesima
%(\EQUref{eq:FFCI_piano_deriv}) si ottiene $v_1+2a\,\vartheta_1$, e
%dunque
%\begin{equation}\label{eq:vinc_B_gerber}
%	v_1 + 2a\,\vartheta_1 = -2\bar v.
%\end{equation}
%Il coefficiente $2a$ --- il braccio del vincolo rispetto al polo --- è
%ciò che fa entrare $\vartheta_1$ nell'equazione.
%
%\paragraph{Cerniera in $C$.} È il vincolo che lega fra loro i due
%corpi: impone che lo spostamento del punto $C$, calcolato come punto di
%$\mathcal{C}_1$, coincida con lo spostamento dello stesso punto $C$,
%calcolato come punto di $\mathcal{C}_2$. Come punto di $\mathcal{C}_2$,
%essendo $O_2\equiv C$, lo spostamento è semplicemente
%$(u_2,\,v_2)$; come punto di $\mathcal{C}_1$, per la formula
%fondamentale, è $(u_1,\,v_1+3a\,\vartheta_1)$. Uguagliando:
%\begin{equation}\label{eq:vinc_C_gerber}
%	u_1 = u_2, \qquad v_1 + 3a\,\vartheta_1 = v_2.
%\end{equation}
%Si noti bene cosa \emph{non} compare: nessuna equazione lega
%$\vartheta_1$ a $\vartheta_2$. È lo svincolo rotazionale della cerniera
%a mancare qui: $\mathcal{C}_2$ resta libero di ruotare indipendentemente
%da $\mathcal{C}_1$, e questa libertà è, letteralmente, l'assenza di
%un'equazione (torneremo su questo in generale nel
%\PARref{sec:congruenza_generale}).
%
%\paragraph{Carrello in $D$.} Impone, in tutto analogamente al carrello
%in $B$ ma nei parametri di $\mathcal{C}_2$,
%\begin{equation}\label{eq:vinc_D_gerber}
%	v_2 + 2a\,\vartheta_2 = -\bar v.
%\end{equation}
%
%\paragraph{Matrice, vettore dei cedimenti e soluzione.}
Il sistema di sei equazioni \eqref{eq:vinc_A_gerber}--\eqref{eq:vinc_D_gerber}
può essere facilmente scritto in forma matriciale come $\bm{Au}=\bm{s}$, dove
\begin{equation}\label{eq:A_gerber_SCR}
	\bm{A} =
	\begin{bmatrix}
		1 &  0 &   0 & 0 & 0 &  0 \\
		0 &  1 &   0 & 0 & 0 &  0 \\
		0 &  1 &  2a & 0 & 0 &  0 \\
		-1 &  0 &   0 & 1 & 0 &  0 \\
		0 & -1 & -3a & 0 & 1 &  0 \\
		0 &  0 &   0 & 0 & 1 & 2a
	\end{bmatrix},
	\qquad
	\bm{s} = \bigl\{0\;\;  0 \;\; -2\bar v \;\; 0 \;\; 0 \;\; -\bar v
	\bigr\}^{\!\top}.
\end{equation}
In questo semplice caso, il sistema si risolve
per sostituzione in avanti, una riga alla volta:
\begin{align*}
	\text{riga (i):}&\quad u_1 = 0,\\
	\text{riga (ii):}&\quad v_1 = 0,\\
	\text{riga (iii):}&\quad v_1+2a\,\vartheta_1=-2\bar v
	\;\Longrightarrow\; \vartheta_1=-\frac{\bar v}{a},\\
	\text{riga (iv):}&\quad u_2=u_1=0,\\
	\text{riga (v):}&\quad v_2=v_1+3a\,\vartheta_1=-3\bar v,\\
	\text{riga (vi):}&\quad v_2+2a\,\vartheta_2=-\bar v
	\;\Longrightarrow\;
	\vartheta_2=\frac{-\bar v+3\bar v}{2a}=\frac{\bar v}{a}.
\end{align*}
La configurazione del sistema, descritta da $\bm{u}$, è dunque completamente determinata:
$\mathcal{C}_1$ ruota in senso orario attorno ad $A$ di ampiezza
$\bar v/a$; $\mathcal{C}_2$ trasla verticalmente di $-3\bar v$ e ruota
in senso antiorario di $\bar v/a$. Una verifica diretta conferma la
coerenza della soluzione con i dati: il cedimento in $D$ vale
$v_2+2a\vartheta_2=-3\bar v+2\bar v=-\bar v$, e quello in $B$
vale $v_1+2a\vartheta_1=-2\bar v$.

La \FIGref{fig:gerber_deformata} illustra la configurazione così
trovata.
%
Abbiamo così costruito e risolto, su un caso concreto, l'intero
problema cinematico del sistema: il vettore di configurazione, la
matrice dei vincoli, il vettore dei cedimenti, e la loro soluzione.
Vediamo ora come la stessa costruzione si generalizzi a un sistema
qualunque.


\begin{figure}[htbp]
	\centering
	\begin{tikzpicture}[>=Stealth, line width=0.8pt, font=\small,
		nodo/.style={circle, fill=white, draw=black,
			inner sep=0pt, minimum size=4pt, line width=0.7pt},
		nodoR/.style={circle, fill=black, draw=black,
			inner sep=0pt, minimum size=3pt}]
		
		\def\a{1.0}
		
		% --- configurazione di riferimento (tratteggiata) ---
		\draw[dashed, gray!60, line width=0.9pt] (0,0) -- ({5*\a},0);
		\node[gray!70, above=2pt] at (0,0)      {$A$};
		\node[gray!70, above=2pt] at ({2*\a},0) {$B$};
		\node[gray!70, above=2pt] at ({3*\a},0) {$C$};
		\node[gray!70, above=2pt] at ({5*\a},0) {$D$};
		\node[nodo] at (0,0) {};
		\node[nodo] at ({2*\a},0) {};
		\node[nodo] at ({3*\a},0) {};
		\node[nodo] at ({5*\a},0) {};
		
		% --- configurazione deformata ---
		\coordinate (Ad) at (0,0);
		\coordinate (Bd) at ({2*\a},-0.6);
		\coordinate (Cd) at ({3*\a},-0.9);
		\coordinate (Dd) at ({5*\a},-0.3);
		
		\draw[line width=1.6pt] (Ad) -- (Cd);
		\draw[line width=1.6pt] (Cd) -- (Dd);
		
		\node[nodoR] at (Ad) {};
		\node[nodoR, label={[yshift=-11pt]below:$B$}] at (Bd) {};
		\node[nodoR, label={[xshift=6pt,yshift=-9pt]below:$C$}] at (Cd) {};
		\node[nodoR, label={[yshift=-11pt]below:$D$}] at (Dd) {};
		
		% --- riferimenti orizzontali locali per gli angoli ---
		\draw[dashed, gray!60, thin] (Ad) -- ++(0.9,0);
		\draw[dashed, gray!60, thin] (Cd) -- ++(0.9,0);
		
		% --- angolo theta_1 in A (orario, sotto l'orizzontale) ---
		\draw[->, thin] ({0.6*\a},0) arc[start angle=0, end angle=-16.7, radius={0.6*\a}];
		\node[font=\scriptsize] at (1,-0.16) {$\vartheta_1$};
		
		% --- angolo theta_2 in C (antiorario, sopra l'orizzontale locale) ---
		\draw[->, thin] ({3*\a+0.6},-0.9) arc[start angle=0, end angle=16.7, radius=0.6];
		\node[font=\scriptsize] at ({3*\a+1},-0.76) {$\vartheta_2$};
		
		% --- cedimenti ---
		\draw[->, thick] ({2*\a},-0.05) -- ({2*\a},-0.55)
		node[right=1pt, pos=0.6, font=\scriptsize] {$2\bar v$};
		\draw[->, thick] ({5*\a},-0.05) -- ({5*\a},-0.25)
		node[right=1pt, pos=0.6, font=\scriptsize] {$\bar v$};
		
	\end{tikzpicture}
	\caption{Configurazione deformata della trave di Gerber (in grigio
		tratteggiato la configurazione di riferimento, non deformata),
		soluzione del problema cinematico del
		\PARref{sec:gerber_congruenza}; spostamenti esagerati per
		chiarezza. $\mathcal{C}_1$ ruota rigidamente attorno ad $A$
		dell'angolo $\vartheta_1=-\bar v/a$ (orario); $\mathcal{C}_2$
		ruota, indipendentemente da $\mathcal{C}_1$, dell'angolo
		$\vartheta_2=+\bar v/a$ (antiorario) attorno al punto $C$, che
		segue $\mathcal{C}_1$ senza distaccarsene. Il punto $C$ non
		presenta alcuna discontinuità di \emph{spostamento} --- le due
		travi restano a contatto --- ma la pendenza della configurazione
		deformata cambia bruscamente in $C$: è la manifestazione
		visibile dello svincolo rotazionale della cerniera interna.}
	\label{fig:gerber_deformata}
\end{figure}


\subsection{La formulazione generale}\label{sec:congruenza_generale}

Le quattro equazioni di vincolo appena scritte per la Gerber
condividono una struttura comune, che conviene ora isolare. 




I vincoli  studiati nel \CAPref{cap:vincoli_CR} impongono
relazioni \emph{lineari} tra le componenti di $\bm{u}$. Ogni
equazione scalare di vincolo può essere scritta nella forma
\begin{equation}\label{eq:vincolo_scalare}
	\bm{a}_k^{\top}\bm{u} = s_k,
	\qquad k = 1, \dots, m,
\end{equation}
dove $\bm{a}_k \in \mathbb{R}^{n}$ è il \emph{vettore di vincolo}\index{vettore di vincolo}
dell'equazione $k$-esima --- la cui trasposta $\bm{a}_k^{\top}$
sarà la $k$-esima riga della matrice dei vincoli $\bm{A}$
assemblata nel \PARref{sec:prob_cin} --- e $s_k \in \mathbb{R}$
è il cedimento (nullo per un vincolo perfetto).
Il vettore $\bm{a}_k$ ha componenti non nulle solo nei blocchi corrispondenti ai corpi
coinvolti dal vincolo: un solo blocco per i vincoli
\emph{esterni}, due blocchi per i vincoli \emph{interni}.

Si dice che la configurazione del sistema è
\emph{esternamente congruente}\index{congruenza!esterna} con i cedimenti
$\bm{s}$ se gli spostamenti dei punti di vincolo
soddisfano le equazioni di vincolo~\eqref{eq:vincolo_scalare}
per ogni $k = 1, \dots, m$.

\paragraph{Vincoli esterni.}
Un vincolo esterno impone una condizione sullo spostamento
assoluto di un punto $A \in \mathcal{C}_i$ o sulla sua
rotazione. Per un carrello di asse $\ab$, con componenti
cartesiane $\ab = a_x\,\ab_x + a_y\,\ab_y$, la condizione
cinematica è:
\begin{equation}\label{eq:vincolo_est_vett}
	\ub_i(A) \cdot \ab = s.
\end{equation}
Si dice che la configurazione del corpo $\mathcal{C}_i$ è
\emph{internamente congruente}\index{congruenza!interna} se lo spostamento di ogni
suo punto è compatibile con il vincolo di rigidità,
ovvero se può essere espresso nella forma
\begin{equation}\label{eq:rigidita_i}
	\ub_i(A) = \ub(O_i) + \bm{\upvartheta}_i \times
	\overrightarrow{O_i A},
\end{equation}
dove $\bm{\upvartheta}_i = \vartheta_i\,\ab_z$.
Questa condizione è automaticamente soddisfatta dalla
parametrizzazione tramite il vettore di configurazione
$\bm{u}$.

Sostituendo la~\eqref{eq:rigidita_i} nella~\eqref{eq:vincolo_est_vett}
e applicando la proprietà ciclica del prodotto misto:
\begin{align}\label{eq:vincolo_est_scal_cin}
	\ub_i(A) \cdot \ab
	&= \ub(O_i)\cdot\ab
	+ \bigl(\bm{\upvartheta}_i \times \overrightarrow{O_i A}
	\bigr)\cdot\ab \notag\\
	&= \ub(O_i)\cdot\ab
	+ \bm{\upvartheta}_i \cdot
	\bigl(\overrightarrow{O_i A} \times \ab\bigr) \notag\\
	&= \ub(O_i)\cdot\ab
	+ \vartheta_i\,
	\underbrace{\bigl(\overrightarrow{O_i A}
		\times \ab\bigr)\cdot\ab_z}_{
		=\,x_i(A)\,a_y - y_i(A)\,a_x}
	= s.
\end{align}
Il termine $(\overrightarrow{O_i A}\times\ab)\cdot\ab_z$
è il \emph{momento} del versore $\ab$ rispetto al polo
$O_i$: esso misura il braccio del vincolo rispetto al polo
scelto. La riga corrispondente di $\bm{A}$ è pertanto:
\begin{equation}\label{eq:ak_esterno}
	\bm{a}_k^{\top} =
	\bigl[\,\underbrace{0 \;\cdots\; 0}_{3(i-1)}
	\;\Big|\;
	a_x \;\; a_y \;\;
	x_i(A)\,a_y - y_i(A)\,a_x
	\;\Big|\;
	\underbrace{0 \;\cdots\; 0}_{3(n_c - i)}\,\bigr].
\end{equation}

\paragraph{Vincoli interni.}
Un vincolo interno collega $\mathcal{C}_i$ a $\mathcal{C}_j$
imponendo una condizione sullo spostamento relativo tra il punto $A$ visto come appartenente a $\Cc_i$ e come appartenente
a $\Cc_j$:
\begin{equation}\label{eq:vincolo_int_vett}
\big(\ub_j(A) - \ub_i(A)\big) \cdot \ab = s.
\end{equation}
Applicando la condizione di rigidità~\eqref{eq:rigidita_i}
separatamente ai due corpi e sottraendo:
\begin{multline}\label{eq:vincolo_int_scal}
\big(\ub_j(A) - \ub_i(A)\big)\cdot\ab
	= \\
	\bigl[\ub(O_j)\cdot\ab
	+ \vartheta_j\,(\overrightarrow{O_j A}\times\ab)
	\cdot\ab_z\bigr]
	- \bigl[\ub(O_i)\cdot\ab
	+ \vartheta_i\,(\overrightarrow{O_i A}\times\ab)
	\cdot\ab_z\bigr]
	= s.
\end{multline}
Il vettore di vincolo $\bm{a}_k$ ha ora supporto su due
blocchi con segni opposti:
\begin{multline}\label{eq:ak_interno}
	\bm{a}_k^{\top} =
	\bigl[\,\cdots\;\Big|\;
	\underbrace{-a_x \;\; -a_y \;\;
		-(x_i(A)\,a_y - y_i(A)\,a_x)
	}_{\text{blocco }i}
	\;\Big|\;\cdots \\[6pt]
	\hspace{6em}
	\cdots\;\Big|\;
	\underbrace{+a_x \;\; +a_y \;\;
		+(x_j(A)\,a_y - y_j(A)\,a_x)
	}_{\text{blocco }j}
	\;\Big|\;\cdots\,\bigr].
\end{multline}

\medskip
Specializzando la struttura
generale~\eqref{eq:ak_esterno}
e~\eqref{eq:ak_interno} ai cinque vincoli elementari
del \CAPref{cap:vincoli_CR}, si ottengono le righe di
$\bm{A}$ raccolte nella \TABref{tab:righe_A}.

\medskip
\begin{table}[H]
	\centering
	\resizebox{\textwidth}{!}{%
		\renewcommand{\arraystretch}{1.9}
		\small
		\begin{tabular}{>{\bfseries}lcc}
			\toprule
			Vincolo & Eq. cinematica & Riga di $\bm{A}$ relativa al corpo $i$\\
			\midrule
			Carrello (asse $\ab$)
			& $\ub_i(A)\cdot\ab = s$
			& $\bigl[a_x,\; a_y,\;
			x_i(A)\,a_y - y_i(A)\,a_x\bigr]$\\[4pt]
			Bipendolo
			& $\vartheta_i = \bar{\vartheta}$
			& $\bigl[0,\; 0,\; 1\bigr]$\\[4pt]
			Cerniera ($x$-comp.)
			& $u_i(A) = \bar{u}$
			& $\bigl[1,\; 0,\; -y_i(A)\bigr]$\\[4pt]
			Cerniera ($y$-comp.)
			& $v_i(A) = \bar{v}$
			& $\bigl[0,\; 1,\; \phantom{-}x_i(A)\bigr]$\\[4pt]
			Glifo (asse $\ab$)
			& $\ub_i(A)\cdot\ab = s$, $\;\vartheta_i = \bar{\vartheta}$
			& (due righe: come carrello + bipendolo)\\[4pt]
			Incastro
			& $\ub_i(A) = \bar{\ub}$, $\;\vartheta_i = \bar{\vartheta}$
			& (tre righe: due cerniera + bipendolo)\\
			\bottomrule
	\end{tabular}}
	\caption{Contributo di ciascun vincolo esterno alla matrice $\bm{A}$:
		blocco $i$-esimo della riga $\bm{a}_k^\top$,
		relativo alle componenti $(u_i,\, v_i,\, \vartheta_i)$
		di $\bm{u}$; le componenti nei blocchi $j \neq i$ sono nulle.
		Per un vincolo interno, lo stesso blocco compare con segno
		opposto nel blocco $i$ e con segno positivo nel blocco $j$,
		come nella~\eqref{eq:ak_interno}.}
	\label{tab:righe_A}
\end{table}


\begin{oss}
	Nella~\eqref{eq:ak_interno},
	le coordinate $(x_i(A),\,y_i(A))$ e $(x_j(A),\,y_j(A))$
	sono le coordinate del punto $A$ rispetto ai poli $O_i$ e
	$O_j$ rispettivamente, e in generale differiscono. Solo
	quando $O_i = O_j$ esse coincidono e i coefficienti di
	$\vartheta_i$ e $\vartheta_j$ nella~\eqref{eq:ak_interno}
	sono uguali e opposti.
\end{oss}
%
%\begin{oss}[Vincoli interni come continuità con svincolo]
%	I vincoli interni possono essere reinterpretati in modo
%	unitario come vincoli di \emph{continuità} tra
%	$\mathcal{C}_i$ e $\mathcal{C}_j$, eventualmente
%	parzialmente rilasciati tramite uno \emph{svincolo}.
%	La continuità piena impone l'uguaglianza di tutti e tre
%	i parametri cinematici nel punto di connessione $A$:
%	\begin{equation}\label{eq:continuita_piena}
%		\llbracket\ub\rrbracket(A) = \bm{0}, \qquad
%		\llbracket\vartheta\rrbracket(A) = 0,
%	\end{equation}
%	corrispondenti a 3 equazioni scalari in $\mathbb{R}^2$
%	(2 traslazioni + 1 rotazione). Uno \emph{svincolo}
%	rilascia una o più di queste condizioni:
%	\begin{itemize}[noitemsep]
%		\item \textbf{cerniera interna}: svincolo rotazionale
%		($\llbracket\vartheta\rrbracket$ libero) --- restano 2 equazioni
%		scalari, la~\eqref{eq:vincolo_int_scal} con
%		$\ab = \ab_x$ e $\ab = \ab_y$;
%		\item \textbf{carrello interno} (asse $\ab$): svincolo
%		traslazionale nella direzione ortogonale ad $\ab$ e
%		svincolo rotazionale --- resta 1 equazione scalare,
%		la~\eqref{eq:vincolo_int_scal} con il versore $\ab$
%		dato;
%		\item \textbf{giunzione rigida} (incastro interno): nessuno svincolo ---
%		tutte e 3 le equazioni sono attive.
%	\end{itemize}
%	Nella matrice $\bm{A}$, ogni rilascio corrisponde
%	semplicemente all'assenza di una riga: la struttura
%	algebrica~\eqref{eq:ak_interno} rimane invariata,
%	cambia solo il numero di righe contribuite dal vincolo.
%	Questa prospettiva si generalizza in modo naturale ai
%	corpi deformabili, dove gli svincoli modellano
%	discontinuità fisiche --- come si vedrà nei volumi
%	successivi --- senza richiedere modifiche alla struttura
%	formale della trattazione.
%\end{oss}
%






\medskip

%\paragraph{La matrice dei vincoli e il problema cinematico.}
Impilando le $m$ equazioni di vincolo di un sistema qualunque ---
ciascuna della forma $\bm{a}_k^\top\bm{u}=s_k$, con $\bm{a}_k^\top$
costruita come nella \eqref{eq:vincolo_est_scal_SCR} o nella
\eqref{eq:vincolo_int_scal_SCR} --- si ottiene la \emph{matrice dei
	vincoli}\index{matrice dei vincoli|textbf} $\bm{A}\in\mathbb{R}^{m\times n}$
e il \emph{vettore dei cedimenti}\index{cedimento|textbf}
$\bm{s}=(s_1,\dots,s_m)^\top$, con $s_k=0$ per ogni vincolo perfetto.

Formuliamo dunque il seguente problema.

\subsection{Problema cinematico, in generale}\index{problema cinematico|textbf}\index{problema cinematico|seealso{dualità}}
Dato un sistema di $n_c$ corpi rigidi $\mathcal{C}_i$
($i = 1, 2, \dots, n_c$) collegati fra loro e al suolo
da $m$ vincoli, e assegnati i cedimenti $s_k$ di ciascun
vincolo ($k = 1, \dots, m$), si chiede di determinare,
se esistono, le configurazioni del sistema --- descritte
dal vettore $\bm{u} \in \mathbb{R}^n$ --- che soddisfano
simultaneamente:
\begin{itemize}[noitemsep]
	\item le \emph{equazioni di rigidità} di ciascun corpo,	che impongono che ogni corpo si muova come un solido
	rigido (congruenza interna);
	\item le \emph{equazioni di vincolo}, che impongono che	gli spostamenti siano compatibili con i cedimenti
	assegnati (congruenza esterna).
\end{itemize}
Una configurazione che è simultaneamente internamente
congruente --- ovvero compatibile con il vincolo di
rigidità di ciascun corpo --- ed esternamente congruente
--- ovvero compatibile con i cedimenti assegnati ai
vincoli --- si dice \emph{cinematicamente ammissibile}\index{configurazione cinematicamente ammissibile}.

Il problema cinematico si riduce dunque a determinare se esiste e se è unica la soluzione del seguente sistema lineare
 \begin{equation}\label{eq:prob_cin_SCR}
 	\bm{A}\,\bm{u} = \bm{s}.
 \end{equation}
 Una configurazione che è simultaneamente internamente
 congruente --- ovvero compatibile con il vincolo di
 rigidità di ciascun corpo --- ed esternamente congruente
 --- ovvero compatibile con i cedimenti assegnati ai
 vincoli --- si dice \emph{cinematicamente ammissibile}\index{configurazione cinematicamente ammissibile}.




%----------------------------------------------------------------------------
\section{L'equilibrio del sistema}\label{sec:equilibrio_SCR}
%----------------------------------------------------------------------------

Passiamo alla statica. Anche qui la domanda relativa alle condizioni di equilibrio di un sistema
di corpi si lascia scomporre in due
richieste indipendenti, dello stesso tipo di quelle incontrate parlando
di congruenza --- non a caso: è precisamente questa somiglianza che
costruiremo, passo dopo passo, fino a renderla un teorema.

La prima richiesta riguarda ancora ciascun corpo preso singolarmente, ed
è interamente nota dal \CAPref{cap:statica_CR}: $\mathcal{C}_1$ dev'essere
in equilibrio sotto l'insieme di tutte le azioni che riceve --- i carichi
attivi \emph{e} le reazioni che gli vengono trasmesse dai vincoli, esterni
e interni; lo stesso, sotto le proprie azioni, per $\mathcal{C}_2$. La
seconda richiesta riguarda invece il modo in cui l'azione si trasmette
attraverso un vincolo interno: per il principio di azione e reazione,
ciò che $\mathcal{C}_1$ scambia con $\mathcal{C}_2$ nel punto di contatto
è, per $\mathcal{C}_2$, esattamente l'opposto di ciò che $\mathcal{C}_2$
scambia con $\mathcal{C}_1$ --- stessa retta di applicazione, stesso
punto, intensità uguale e verso opposto. 

%Più precisamente, possiamo formulare il problema statico come segue.
%
%\paragraph{Problema statico.}\index{problema statico|textbf}\index{problema statico|seealso{dualità}}
%Dato un sistema di $n_c$ corpi vincolati, e assegnate le forze
%attive applicate ai corpi  si
%chiede di determinare, se esistono, le reazioni vincolari
% che, insieme alle forze attive,
%soddisfano simultaneamente:
%\begin{itemize}[noitemsep]
%	\item le condizioni di  \emph{equilibrio} di ciascun
%	corpo $\mathcal{C}_i$, inteso come solido rigido
%	soggetto alle forze attive e alle reazioni dei vincoli
%	ad esso connessi;
%	\item le \emph{condizioni di reazione}, che impongono
%	che le reazioni siano orientate secondo i versori
%	di vincolo.
%\end{itemize}
%Un sistema di reazioni che è simultaneamente equilibrato internamente
%--- compatibile con l'equilibrio di ciascun corpo --- ed esternamente
%--- diretto secondo i versori di vincolo --- si dice
%\emph{staticamente ammissibile}\index{configurazione staticamente ammissibile}.

%\tR{
%
%\paragraph{Problema statico.}\index{problema statico|textbf}\index{problema statico|seealso{dualità}}
%Dato lo stesso sistema di $n_c$ corpi rigidi collegati da $m$ vincoli,
%e assegnate le forze attive --- descritte dal vettore delle forze
%generalizzate $\bm f\in\mathbb{R}^n$ --- si chiede di determinare, se
%esistono, le reazioni vincolari --- descritte dal vettore
%$\bm r\in\mathbb{R}^m$ --- che soddisfano simultaneamente:
%\begin{itemize}[noitemsep]
%	\item le \emph{equazioni di equilibrio} di ciascun corpo
%	$\mathcal{C}_i$, inteso come solido rigido soggetto ai carichi
%	attivi e alle reazioni dei vincoli ad esso connessi (equilibrio
%	interno);
%	\item le \emph{equazioni di reazione}, che impongono che ciascuna
%	reazione sia diretta secondo il versore $\ab_k$ del vincolo cui è
%	associata (equilibrio esterno).
%\end{itemize}
%Un sistema di reazioni che è simultaneamente equilibrato internamente
%--- compatibile con l'equilibrio di ciascun corpo --- ed esternamente
%--- diretto secondo i versori di vincolo --- si dice
%\emph{staticamente ammissibile}\index{configurazione staticamente ammissibile}.
%Il problema statico si riduce dunque a determinare se esiste ed è unica
%la soluzione del sistema lineare
%\begin{equation}\label{eq:prob_stat_SCR}
%	\bm{B}\,\bm{r} = -\bm{f}, \qquad \bm{B}=\bm{A}^{\!\top}.
%	\end{equation}}


Come per la congruenza, conviene
vedere subito queste due richieste all'opera su un caso esplicito, prima
di scriverle in generale, ricordando che esistoni due modi per caratterizzare l'equilibrio e dunque per risolvere il problema statico.

%----------------------------------------------------------------------------
\subsection{Un caso esplicito: la trave di Gerber, letta due volte}\label{sec:gerber_equilibrio}
%----------------------------------------------------------------------------

Riprendiamo la Gerber di \FIGref{fig:gerber_motiv}, ora sotto un carico
attivo: una forza $2F$, verticale e diretta verso il basso, applicata nel
punto medio del tratto $AB$ di $\mathcal{C}_1$; una forza $F$, anch'essa
verticale e verso il basso, applicata nel punto medio del tratto
$CD$ di $\mathcal{C}_2$ (\FIGref{fig:gerber_carichi}). È lo stesso
sistema, con la stessa geometria e la stessa cerniera interna in $C$, che
abbiamo già letto in chiave cinematica: qui lo leggiamo in chiave
statica, e fra poco energetica.

\begin{figure}[htbp]
	\centering
	\begin{tikzpicture}[>=Stealth, line width=0.8pt, font=\small,
		nodo/.style={circle, fill=white, draw=black,
			inner sep=0pt, minimum size=5pt, line width=0.8pt},
		cint/.style={circle, fill=white, draw=black,
			inner sep=0pt, minimum size=6pt, line width=0.8pt},
		richiam/.style={thin, gray, dashed},
		tacca/.style={thin, gray!70}]
		
		\def\a{1.0}
		\def\tw{0.22}\def\thA{0.38}\def\thR{0.32}
		\def\bw{0.16}\def\gw{0.30}\def\sr{0.08}
		
		\draw[line width=2.2pt] (0,0) -- ({3*\a},0);
		\draw[line width=2.2pt] ({3*\a},0) -- ({5*\a},0);
		
		\node[above=2pt] at (0,0)       {$A$};
		\node[above=2pt] at ({2*\a},0)  {$B$};
		\node[above=2pt] at ({3*\a},0)  {$C$};
		\node[above=2pt] at ({5*\a},0)  {$D$};
		
		\node at ({\a},  -0.32) {$\mathcal{C}_1$};
		\node at ({4*\a},-0.32) {$\mathcal{C}_2$};
		
		%% cerniera in A
		\draw (0,0) -- ++(-\tw,-\thA) -- ++(2*\tw,0) -- cycle;
		\draw (-\gw,-\thA) -- (\gw,-\thA);
		\foreach \x in {-0.24,-0.12,0,0.12,0.24}
		\draw[thin] (\x,-\thA) -- ++(-0.09,-0.11);
		\node[nodo] at (0,0) {};
		
		%% carrello intermedio in B
		\draw ({2*\a},0) -- ++(-\tw,-\thR) -- ++(2*\tw,0) -- cycle;
		\draw ({2*\a-\gw},-\thR) -- ({2*\a+\gw},-\thR);
		\draw ({2*\a-\bw},-\thR-0.12) circle (2.4pt);
		\draw ({2*\a+\bw},-\thR-0.12) circle (2.4pt);
		\draw ({2*\a-\gw},-\thR-0.24) -- ({2*\a+\gw},-\thR-0.24);
		\foreach \x in {-0.24,-0.12,0,0.12,0.24}
		\draw[thin] ({2*\a+\x},-\thR-0.24) -- ++(-0.09,-0.11);
		\fill[white] ({2*\a-\sr},0) arc(180:360:\sr) -- cycle;
		\draw ({2*\a-\sr},0) arc(180:360:\sr);
		
		%% cerniera interna in C
		\node[cint] at ({3*\a},0) {};
		
		%% carrello terminale in D
		\draw ({5*\a},0) -- ++(-\tw,-\thR) -- ++(2*\tw,0) -- cycle;
		\draw ({5*\a-\gw},-\thR) -- ({5*\a+\gw},-\thR);
		\draw ({5*\a-\bw},-\thR-0.12) circle (2.4pt);
		\draw ({5*\a+\bw},-\thR-0.12) circle (2.4pt);
		\draw ({5*\a-\gw},-\thR-0.24) -- ({5*\a+\gw},-\thR-0.24);
		\foreach \x in {-0.24,-0.12,0,0.12,0.24}
		\draw[thin] ({5*\a+\x},-\thR-0.24) -- ++(-0.09,-0.11);
		\node[nodo] at ({5*\a},0) {};
		
		%% carico 2F nel punto medio di AB (x=a)
		\draw[thin] (\a,0.05) -- (\a,-0.05);
		\draw[->, thick] (\a,0.85) -- (\a,0.10)
		node[right=3pt, pos=0] {$2F$};
		
		%% carico F nel punto medio di C_2 (x=4a)
		\draw[thin] ({4*\a},0.05) -- ({4*\a},-0.05);
		\draw[->, thick] ({4*\a},0.85) -- ({4*\a},0.10)
		node[right=3pt, pos=0] {$F$};
		
		%% quotatura
		\draw[richiam] (0,      -\thA-0.15) -- (0,      -1.55);
		\draw[richiam] ({2*\a}, -1.40)      -- ({2*\a}, -1.55);
		\draw[richiam] ({3*\a}, -0.10)      -- ({3*\a}, -1.55);
		\draw[richiam] ({5*\a}, -1.40)      -- ({5*\a}, -1.55);
		
		\draw[tacca] (0,     -1.55) -- (0,     -1.75);
		\draw[tacca] ({2*\a},-1.55) -- ({2*\a},-1.75);
		\draw[tacca] ({3*\a},-1.55) -- ({3*\a},-1.75);
		\draw[tacca] ({5*\a},-1.55) -- ({5*\a},-1.75);
		
		\draw[<->, thin, gray] (0,     -1.65) -- ({2*\a},-1.65);
		\node[above=0.04] at ({\a},   -1.65) {$2a$};
		\draw[<->, thin, gray] ({2*\a},-1.65) -- ({3*\a},-1.65);
		\node[above=0.04] at ({2.5*\a},-1.65) {$a$};
		\draw[<->, thin, gray] ({3*\a},-1.65) -- ({5*\a},-1.65);
		\node[above=0.04] at ({4*\a}, -1.65) {$2a$};
		
	\end{tikzpicture}
	\caption{La trave di Gerber sotto carico attivo: una forza $2F$
		verso il basso nel punto medio di $AB$, su $\mathcal{C}_1$, e
		una forza $F$ verso il basso nel punto medio di $\mathcal{C}_2$.
		Stesso sistema di \FIGref{fig:gerber_motiv}, qui nella lettura
		statica.}
	\label{fig:gerber_carichi}
\end{figure}

\paragraph{La lettura vettoriale.}
Isoliamo i due corpi, esattamente come si farebbe per un corpo singolo,
avendo cura di rendere esplicita anche l'azione mutua in $C$. Sulle
reazioni esterne adottiamo la convenzione consueta, positive se dirette
come gli assi: $(X_A,Y_A)$ in $A$ (cerniera, due componenti incognite),
$Y_B$ in $B$ (carrello, reazione verticale), $Y_D$ in $D$ (carrello,
reazione verticale). Sull'azione interna in $C$ introduciamo la coppia
incognita $(X_C,Y_C)$, definita come l'azione che $\mathcal{C}_1$
esercita su $\mathcal{C}_2$: per il principio di azione e reazione,
l'azione che $\mathcal{C}_2$ esercita su $\mathcal{C}_1$ è allora,
automaticamente, $(-X_C,-Y_C)$ --- non una nuova incognita, ma la stessa
coppia già introdotta, con il segno cambiato.

Cominciamo da $\mathcal{C}_1$, che vede il carico $2F$ in $x=a$, le reazioni
esterne $(X_A,Y_A)$ e $Y_B$, e la reazione mutua in $C$ vista da
$\mathcal{C}_1$ --- cioè, per quanto appena detto, $(-X_C,-Y_C)$: 
\begin{align}
& X_A - X_C = 0
	\;\;\Longrightarrow\;\; X_A=0, \label{eq:eqC1_Fx}\\
& Y_A + Y_B - Y_C - 2F = 0
\;\;\Longrightarrow\;\; Y_A = \tfrac{3}{4}F, \label{eq:eqC1_Fy}\\
	& Y_B\cdot 2a - 2F\cdot a - Y_C\cdot 3a = 0
	\;\;\Longrightarrow\;\; Y_B = \tfrac{7}{4}F. \label{eq:eqC1_MA}\\
	\end{align}

Passiamo a $\mathcal{C}_2$, che vede solo il carico $F$ in $x=4a$, la
reazione mutua $(X_C,Y_C)$ in $C$ e la reazione $Y_D$ in $D$:
\begin{align}
& X_C = 0, \label{eq:eqC2_Fx}\\
& Y_C + Y_D - F = 0
\;\;\Longrightarrow\;\; Y_C = \tfrac{1}{2}F, \label{eq:eqC2_Fy}\\
	& Y_D\cdot 2a - F\cdot a = 0
	\;\;\Longrightarrow\;\; Y_D = \tfrac{1}{2}F. \label{eq:eqC2_MC}\\
	\end{align}

Sei equazioni scalari, sei incognite --- tante quante i vincoli avevano
imposto in chiave cinematica nella \SECref{sec:congruenza_SCR} --- e la
soluzione è univoca:
\begin{equation}\label{eq:reazioni_gerber}
	X_A=0,\quad Y_A=\tfrac34 F,\quad Y_B=\tfrac74 F,\quad
	X_C=0,\quad Y_C=\tfrac12 F,\quad Y_D=\tfrac12 F.
\end{equation}
Come verifica, si osservi che le sole tre equazioni cardinali
dell'intero sistema, trattato come un unico blocco, sono anch'esse
soddisfatte:
\[
\underbrace{Y_A+Y_B+Y_D-2F-F}_{\;=\;0}\;=\;0,
\qquad
\underbrace{Y_B\cdot 2a + Y_D\cdot 5a - 2F\cdot a - F\cdot 4a}_{\;=\;0}
\;=\;0.
\]

\begin{oss}
	Le due equazioni globali appena verificate sono una conseguenza
	necessaria dell'equilibrio, non una condizione sufficiente a
	stabilirlo: bastano tre equazioni per l'intero sistema, contro le sei
	che abbiamo dovuto scrivere corpo per corpo. È possibile costruire
	configurazioni di forze che soddisfano le equazioni cardinali
	dell'assieme --- risultante e momento risultante nulli, come se il
	sistema fosse un unico corpo rigido --- pur non essendo affatto di
	equilibrio per le singole parti.
\end{oss}


Mettiamo ora in forma matriciale le sei equazioni di equilibrio già
risolte in \SECref{sec:gerber_equilibrio}, e confrontiamole con \eqref{eq:A_gerber_SCR}.


Raccogliamo, corpo per corpo, risultante e momento (rispetto al polo)
dei soli carichi attivi:
\begin{equation}\label{eq:f_gerber}
	\bm f = \{0\;\; -2F \;\;   -2Fa \;\; 0 \;\; -F \;\; -Fa\}^{\!\top},
\end{equation}
e le reazioni:
\begin{equation}\label{eq:r_gerber}
	\bm r = \{X_A  \;\;  Y_A \;\; Y_B \;\; X_C \;\; Y_C \;\; Y_D\}^{\!\top}.
\end{equation}


Le equazioni \eqref{eq:eqC2_Fx}--\eqref{eq:eqC1_Fy} sono le componenti
di $\bm B\bm r=-\bm f$, con
\begin{equation}\label{eq:B_gerber}
	\bm B = \begin{bmatrix}
		1 & 0 & 0 & -1 & 0  & 0 \\
		0 & 1 & 1 & 0  & -1 & 0 \\
		0 & 0 & 2a& 0  & -3a& 0 \\
		0 & 0 & 0 & 1  & 0  & 0 \\
		0 & 0 & 0 & 0  & 1  & 1 \\
		0 & 0 & 0 & 0  & 0  & 2a
	\end{bmatrix}.
\end{equation}



Osserviamo che, trasponendo \eqref{eq:A_gerber_SCR}, si ottiene esattamente
$\bm B$:
\begin{equation}\label{eq:dualismo_gerber}
	\bm B = \bm A^{\!\top}.
\end{equation}
Non è un caso: mostreremo più avanti che questo vale sempre, per qualunque
sistema e qualunque vincolo, e non solo per la Gerber.


\paragraph{Soppesare le reazioni: la via del lavoro virtuale.}
C'è un secondo modo di trovare le stesse sei reazioni, che non passa per
la soluzione simultanea di un sistema di equazioni corpo per corpo. Nel capitolo sul corpo singolo, il lavoro virtuale
serviva a \emph{caratterizzare} l'equilibrio, non a \emph{calcolare} una
reazione. L'idea però è stessa già incontrata all'apertura della statica del
corpo singolo, quando si osservava che un peso, prima ancora di posarlo
su una bilancia, lo si \emph{soppesa}: gli si imprime uno spostamento e
se ne valuta la fatica che costa. Qui facciamo lo stesso, non con un
oggetto ma con un intero sistema.

Il punto di partenza è ciò che già sappiamo dal
\CAPref{cap:statica_CR}: un corpo rigido è in equilibrio se e solo se il
lavoro virtuale delle forze su di esso applicate è nullo per \emph{ogni}
moto rigido del corpo, non solo per un moto particolare. Sappiamo, dalla
lettura vettoriale appena conclusa, che il sistema è in equilibrio: in
particolare lo sono, separatamente, $\mathcal{C}_1$ e $\mathcal{C}_2$.
Per il principio appena richiamato, applicato una volta a ciascuno, il
lavoro delle forze su $\mathcal{C}_1$ è allora nullo per ogni moto rigido
di $\mathcal{C}_1$, e lo stesso vale per $\mathcal{C}_2$ --- due
affermazioni vere \emph{indipendentemente} l'una dall'altra, per due moti
scelti anche a piacere e senza alcuna relazione fra loro.

Scegliamo ora, fra tutti i moti possibili, quelli \emph{congruenti}: che
rispettino cioè la cerniera in $C$, assegnando alle due travi lo stesso
spostamento di quel punto (pur lasciandole libere di ruotare in modo
diverso attorno ad esso). Sommando i due lavori --- entrambi nulli
separatamente, dunque nulla anche la loro somma --- succede qualcosa di
utile.

\begin{oss}
	Nella somma, il contributo della coppia di reazione mutua in $C$
	scompare, qualunque sia il suo valore. Per il principio di azione e
	reazione le due travi si scambiano in $C$ una forza uguale e
	contraria; per la congruenza appena imposta, quel punto materiale
	subisce, visto dall'una o dall'altra trave, \emph{lo stesso}
	spostamento. Il lavoro che $\mathcal{C}_1$ compie su $\mathcal{C}_2$
	attraverso quella reazione e il lavoro che $\mathcal{C}_2$ compie su
$\mathcal{C}_1$ sono dunque uguali in modulo e opposti in segno: si
elidono sempre, senza bisogno di conoscere la reazione stessa. 
\end{oss}

Nella somma restano dunque solo i carichi attivi e le reazioni
\emph{esterne} --- quelle in $A$, $B$, $D$. Se scegliamo il moto
congruente in modo che tutte le reazioni esterne tranne una restino
ferme nel loro punto di applicazione, la richiesta che il \textit{lavoro totale} --- somma 
di quello compiuto dalle forze attive e da quelle reattive --- sia nullo coinvolge una sola incognita, e la determina. Un moto congruente
di questo tipo si costruisce sopprimendo, solo concettualmente, il
vincolo che a quella reazione corrisponde: il sistema, con un vincolo in
meno, guadagna un grado di libertà, e i moti congruenti compatibili con
i vincoli rimasti sono esattamente la famiglia a un parametro che serve.
Vediamola funzionare due volte sulla Gerber.

\medskip

\emph{Reazione $Y_D$} (\FIGref{fig:gerber_mecc_D}). Sopprimiamo il
carrello in $D$. Su $\mathcal{C}_1$ restano cerniera in $A$ e carrello in
$B$: due vincoli che, da soli, bloccano completamente un corpo rigido
piano (tre condizioni per tre parametri), sicché $\mathcal{C}_1$ non ha
alcun moto congruente disponibile all'infuori di quello nullo, e in
particolare il punto $C$ resta fermo. La cerniera interna impone allora
che anche $\mathcal{C}_2$ passi per $C$ restando fermo in quel punto:
l'unico moto congruente residuo è una rotazione rigida di
$\mathcal{C}_2$ attorno a $C$, di ampiezza $\vartheta_2$ (presa positiva
se antioraria, come sempre). Essendo $C$ e $D$ sull'asse della trave, lo
spostamento verticale di un punto di $\mathcal{C}_2$ a distanza $d$ da
$C$ vale, per la formula fondamentale della cinematica rigida, $d\,
\vartheta_2$; in particolare
\[
v(D) = 2a\,\vartheta_2, \qquad v(4a) = a\,\vartheta_2,
\]
dove $v(4a)$ è lo spostamento del punto di applicazione del carico $F$.
Le reazioni in $A$ e $B$ non compiono lavoro, essendo $\mathcal{C}_1$
fermo; la reazione mutua in $C$ nemmeno, per l'\hyperref[]{osservazione}
appena fatta. Restano il carico $F$, verso il basso, e la reazione
incognita $Y_D$, verso l'alto per come l'abbiamo definita: il lavoro
totale nullo dà
\begin{equation}\label{eq:sopp_YD}
Y_D\cdot\bigl(2a\,\vartheta_2\bigr) \;-\; F\cdot\bigl(a\,
\vartheta_2\bigr) \;=\; 0
\quad\Longrightarrow\quad Y_D = \tfrac12 F,
\end{equation}
essendo $\vartheta_2$ arbitrario e non nullo. È lo stesso valore già
trovato in \eqref{eq:eqC2_MC}, per una strada del tutto diversa.

\begin{figure}[htbp]
\centering
\begin{tikzpicture}[>=Stealth, line width=0.8pt, font=\small,
nodo/.style={circle, fill=white, draw=black,
	inner sep=0pt, minimum size=5pt, line width=0.8pt},
cint/.style={circle, fill=white, draw=black,
	inner sep=0pt, minimum size=6pt, line width=0.8pt},
fantasma/.style={thin, gray, dashed},
tacca/.style={thin, gray!70}]

\def\a{1.0}
\def\tw{0.22}\def\thA{0.38}\def\thR{0.32}
\def\bw{0.16}\def\gw{0.30}\def\sr{0.08}

%% riferimento indeformato, tratteggiato
\draw[fantasma] (0,0) -- ({5*\a},0);

%% configurazione del meccanismo: C1 fermo, C2 ruotato attorno a C
\draw[line width=2.2pt] (0,0) -- ({3*\a},0);
\draw[line width=2.2pt] ({3*\a},0) -- ({5*\a},-0.6);

\node[above=2pt] at (0,0)      {$A$};
\node[above=2pt] at ({2*\a},0) {$B$};
\node[above=2pt] at ({3*\a},0) {$C$};
\node[below right=1pt] at ({5*\a},-0.6) {$D$};

\node at ({1.4*\a},0.28)  {$\mathcal{C}_1$};
\node at ({4.3*\a},-0.2)   {$\mathcal{C}_2$};

%% angolo virtuale theta_2 in C
\draw[tacca] ({3.2*\a-0.35},0) -- ({3.2*\a+0.35},0);
\draw[thin] ({3.2*\a+0.3},0) arc (0:-16.7:0.3);
\node[right=1pt] at ({3*\a+0.32},-0.4) {$\vartheta_2$};

%% cerniera in A e carrello in B, invariati
\draw (0,0) -- ++(-\tw,-\thA) -- ++(2*\tw,0) -- cycle;
\draw (-\gw,-\thA) -- (\gw,-\thA);
\foreach \x in {-0.24,-0.12,0,0.12,0.24}
\draw[thin] (\x,-\thA) -- ++(-0.09,-0.11);
\node[nodo] at (0,0) {};

\draw ({2*\a},0) -- ++(-\tw,-\thR) -- ++(2*\tw,0) -- cycle;
\draw ({2*\a-\gw},-\thR) -- ({2*\a+\gw},-\thR);
\draw ({2*\a-\bw},-\thR-0.12) circle (2.4pt);
\draw ({2*\a+\bw},-\thR-0.12) circle (2.4pt);
\draw ({2*\a-\gw},-\thR-0.24) -- ({2*\a+\gw},-\thR-0.24);
\foreach \x in {-0.24,-0.12,0,0.12,0.24}
\draw[thin] ({2*\a+\x},-\thR-0.24) -- ++(-0.09,-0.11);
\fill[white] ({2*\a-\sr},0) arc(180:360:\sr) -- cycle;
\draw ({2*\a-\sr},0) arc(180:360:\sr);

\node[cint] at ({3*\a},0) {};

%% carrello in D soppresso: fantasma nella posizione originaria
\draw[fantasma] ({5*\a},0) -- ++(-\tw,-\thR) -- ++(2*\tw,0) -- cycle;
\node[gray, above=2pt] at ({5*\a},0) {\footnotesize(soppresso)};
\node[nodo] at ({5*\a},-0.6) {};

%% reazione incognita Y_D, verso l'alto, applicata nel punto attuale
\draw[-{Latex[length=2.8mm]}, thick] ({5*\a},-1.2) -- ({5*\a},-0.65);
\node[below=2pt] at ({5*\a},-0.95) {$Y_D$};

%% carico F, verso il basso, applicato nel punto attuale
\draw[-{Latex[length=2.8mm]}, thick] ({4*\a},0.7) -- ({4*\a},0.1);
\node[right=2pt] at ({3.4*\a},0.5) {$F$};

\end{tikzpicture}
\caption{Il meccanismo usato per determinare $Y_D$: si sopprime
(concettualmente) il carrello in $D$, in grigio nella sua posizione
originaria. Cerniera in $A$ e carrello in $B$ bloccano
$\mathcal{C}_1$, che resta fermo; $\mathcal{C}_2$ è allora libero di
ruotare, di un angolo virtuale $\vartheta_2$, attorno al punto fisso
$C$. Solo $F$ e la reazione incognita $Y_D$ compiono lavoro in
questo moto.}
\label{fig:gerber_mecc_D}
\end{figure}

\medskip

\emph{Reazione $Y_B$} (\FIGref{fig:gerber_mecc_B}). Sopprimiamo invece il
carrello in $B$, lasciando intatti gli altri tre vincoli. Su
$\mathcal{C}_1$ resta la sola cerniera in $A$: $\mathcal{C}_1$ è libero
di ruotare, di un angolo virtuale $\vartheta_1$, attorno ad $A$; il punto
$C$, a distanza $3a$ da $A$, si sposta allora di $v(C) = 3a\,
\vartheta_1$. Per congruenza questo è anche lo spostamento verticale di
$C$ come punto di $\mathcal{C}_2$. Ma il carrello in $D$ non è stato
toccato, $v(D)=0$: $\mathcal{C}_2$ deve dunque accompagnare il moto con
una propria rotazione $\vartheta_2$ attorno a $C$, tale che
\[
v(D) = 3a\,\vartheta_1 + 2a\,\vartheta_2 = 0
\quad\Longrightarrow\quad
\vartheta_2 = -\tfrac32\,\vartheta_1 .
\]
Con questa relazione fra le due rotazioni, il meccanismo ha un solo
parametro libero, $\vartheta_1$ --- esattamente quanti gradi di libertà
si guadagnano sopprimendo un solo vincolo. Gli spostamenti verticali dei
punti che interessano sono allora
\[
v(B) = 2a\,\vartheta_1, \qquad
v(a) = a\,\vartheta_1, \qquad
v(4a) = 3a\,\vartheta_1 + a\,\vartheta_2 = \tfrac32 a\,\vartheta_1,
\]
dove $v(a)$ e $v(4a)$ sono gli spostamenti dei punti di applicazione,
rispettivamente, di $2F$ e di $F$. Il lavoro totale nullo dà, con la
sola incognita $Y_B$,
\begin{equation}\label{eq:sopp_YB}
Y_B\cdot\bigl(2a\,\vartheta_1\bigr)
\;-\; 2F\cdot\bigl(a\,\vartheta_1\bigr)
\;-\; F\cdot\bigl(\tfrac32 a\,\vartheta_1\bigr) = 0
\quad\Longrightarrow\quad Y_B = \tfrac74 F,
\end{equation}
di nuovo lo stesso valore già trovato in \eqref{eq:eqC1_MA}.

\begin{figure}[htbp]
\centering
\begin{tikzpicture}[>=Stealth, line width=0.8pt, font=\small,
nodo/.style={circle, fill=white, draw=black,
	inner sep=0pt, minimum size=5pt, line width=0.8pt},
cint/.style={circle, fill=white, draw=black,
	inner sep=0pt, minimum size=6pt, line width=0.8pt},
fantasma/.style={thin, gray, dashed},
tacca/.style={thin, gray!70}]

\def\a{1.0}
\def\tw{0.22}\def\thA{0.38}\def\thR{0.32}
\def\bw{0.16}\def\gw{0.30}\def\sr{0.08}

%% riferimento indeformato, tratteggiato
\draw[fantasma] (0,0) -- ({5*\a},0);

%% configurazione del meccanismo
\draw[line width=2.2pt] (0,0) -- ({3*\a},-0.9);
\draw[line width=2.2pt] ({3*\a},-0.9) -- ({5*\a},0);

\node[above=2pt] at (0,0)             {$A$};
%\node[below=3pt] at ({2*\a},-0.6)     {$B$};
\node[below=3pt] at ({3*\a},-0.9)     {$C$};
\node[above=2pt] at ({5*\a},0)        {$D$};

%\node at ({1*\a},-0.15)  {$\mathcal{C}_1$};
%\node at ({4*\a},-0.6)   {$\mathcal{C}_2$};

%% angolo virtuale theta_1 in A
\draw[tacca] (0,0) -- (0.9,0);
\draw[thin] (0.75,0) arc (0:-16.7:0.75);
\node[below right=1pt] at (0.9,.09) {$\vartheta_1$};

%% angolo virtuale theta_2 in C, contro una tacca locale orizzontale
\draw[tacca] ({3*\a-0.4},-0.9) -- ({3*\a+0.4},-0.9);
\draw[thin] ({3*\a+0.35},-0.9) arc (0:24.2:0.35);
\node[above right=1pt] at ({3*\a+0.4},-1.1) {$\vartheta_2$};

%% cerniera in A, invariata
\draw (0,0) -- ++(-\tw,-\thA) -- ++(2*\tw,0) -- cycle;
\draw (-\gw,-\thA) -- (\gw,-\thA);
\foreach \x in {-0.24,-0.12,0,0.12,0.24}
\draw[thin] (\x,-\thA) -- ++(-0.09,-0.11);
\node[nodo] at (0,0) {};

%% carrello in B soppresso: fantasma nella posizione originaria
\draw[fantasma] ({2*\a},0) -- ++(-\tw,-\thR) -- ++(2*\tw,0) -- cycle;
\node[gray, above=2pt] at ({2*\a},0) {\footnotesize(soppresso)};
%\node[nodo] at ({2*\a},-0.6) {};

%% cerniera interna, nella posizione attuale
\node[cint] at ({3*\a},-0.9) {};

%% carrello in D, invariato (coincide con la posizione di riferimento)
\draw ({5*\a},0) -- ++(-\tw,-\thR) -- ++(2*\tw,0) -- cycle;
\draw ({5*\a-\gw},-\thR) -- ({5*\a+\gw},-\thR);
\draw ({5*\a-\bw},-\thR-0.12) circle (2.4pt);
\draw ({5*\a+\bw},-\thR-0.12) circle (2.4pt);
\draw ({5*\a-\gw},-\thR-0.24) -- ({5*\a+\gw},-\thR-0.24);
\foreach \x in {-0.24,-0.12,0,0.12,0.24}
\draw[thin] ({5*\a+\x},-\thR-0.24) -- ++(-0.09,-0.11);
\node[nodo] at ({5*\a},0) {};

%% reazione incognita Y_B, applicata nel punto attuale
\draw[-{Latex[length=2.8mm]}, thick] ({2*\a},-1.2) -- ({2*\a},-0.6);
\node[below=2pt] at ({1.7*\a},-0.85) {$Y_B$};

%% carichi attivi, applicati nei punti attuali
\draw[-{Latex[length=2.8mm]}, thick] ({1*\a},0.65) -- ({1*\a},0.1);
\node[right=2pt] at ({0.2*\a},0.5) {$2F$};

\draw[-{Latex[length=2.8mm]}, thick] ({4*\a},0.6) -- ({4*\a},0.1);
\node[right=2pt] at ({4*\a},0.5) {$F$};

\end{tikzpicture}
\caption{Il meccanismo usato per determinare $Y_B$: si sopprime il
carrello in $B$. $\mathcal{C}_1$ ruota di $\vartheta_1$ attorno alla
cerniera in $A$; $\mathcal{C}_2$ accompagna il moto ruotando di
$\vartheta_2=-\tfrac32\vartheta_1$ attorno a $C$, in modo da lasciare
fermo il punto $D$, dove il carrello non è stato toccato.}
\label{fig:gerber_mecc_B}
\end{figure}

Le rimanenti reazioni --- $X_A$, $X_C$ (nulle, in assenza di carichi
orizzontali) e $Y_A$, $Y_C$ --- si trovano allo stesso modo, con un
proprio meccanismo ciascuna. In alternativa, è anche possibile sopprimere \textit{tutti}
i vincoli e studiare il più generico meccanismo che ne consegue, determinando il 
lavoro totale. 

Due strade dunque, concettualmente lontane --- l'una che bilancia forze
e momenti corpo per corpo, l'altra che soppesa un vincolo alla volta con
un meccanismo e un'equazione di lavoro --- portano, sulla stessa
struttura, esattamente alle stesse sei reazioni. 





\subsection{Il problema statico, in generale.}\index{problema statico|textbf}
Generalizziamo ora a un sistema qualunque di $n_c$ corpi e $m$ vincoli
quanto visto nel caso della trave di Gerber. Le forze attive su
ciascun corpo $\mathcal C_i$ si raccolgono, come già fatto per il
corpo singolo (\CAPref{cap:statica_CR}), nel blocco
\begin{equation}
	\bm f_i \coloneqq \bigl(F_{ix} \;\; F_{iy}  \;\; C_i\bigr)^{\!\top}\in\mathbb R^3,
\end{equation}
risultante e momento risultante dei soli carichi attivi su
$\mathcal C_i$; l'insieme delle forze attive su tutto il sistema forma
il vettore
\begin{equation}
	\bm f \coloneqq (\bm f_1\;\dots\;\bm f_{n_c})^{\!\top}\in\mathbb R^n,
	\qquad n=3n_c.
\end{equation}
Le reazioni vincolari\nomenclature[L]{$\bm{r}$}{vettore delle reazioni vincolari (intensità scalari $R_k$)} costituiscono il vettore
$\bm{r} \in \mathbb{R}^{m}$:
\begin{equation}
	\bm{r} \coloneqq
	\bigl(R_1\ R_2 \ \dots\ R_m\bigr)^{\!\top},
\end{equation}
dove ogni $R_k$ è l'intensità scalare della reazione associata
alla $k$-esima equazione di vincolo.
\medskip

Per ricavare il contributo di $R_k$ all'equilibrio del sistema
seguiamo, come nel problema cinematico, due passi: prima la
condizione vettoriale --- la forza applicata in un punto --- poi la
condizione di rigidità --- il trasporto di quella forza al polo
scelto per il corpo.

La reazione del $k$-esimo vincolo esterno (carrello di asse
$\ab$) agisce sul corpo $\mathcal{C}_i$ come una forza concentrata
$R_k\,\ab$ applicata nel punto $A$: questa è la condizione
vettoriale, che a questo stadio non dice nulla sull'equilibrio del
corpo nel suo insieme, ma solo sulla forza e sul punto in cui è
applicata. Per la rigidità di $\mathcal{C}_i$, l'effetto di questa
forza sull'equilibrio del corpo è interamente descritto dal suo
risultante $R_k\,\ab$ e dal suo momento rispetto a un polo
qualsiasi --- in particolare rispetto al polo $O_i$ già scelto per
la descrizione cinematica di $\mathcal{C}_i$. Questo trasporto è
lecito proprio perché $\mathcal{C}_i$ è rigido, ed è l'analogo
statico del trasporto dello spostamento dal polo $O_i$ al punto $A$
usato nella formula fondamentale del moto rigido. Il momento di
$R_k\,\ab$ rispetto a $O_i$ vale
\begin{equation*}
	\bigl(\overrightarrow{O_i A} \times R_k\,\ab\bigr)
	\cdot\ab_z
	= R_k\,\bigl(x_i(A)\,a_y - y_i(A)\,a_x\bigr).
\end{equation*}
Raccogliendo le tre componenti relative a $\mathcal{C}_i$
nel vettore
\begin{equation}\label{eq:aki_def}
	\bm{b}_k^{(i)} \coloneqq
	\begin{Bmatrix}
		a_x \\[3pt] a_y \\[3pt]
		x_i(A)\,a_y - y_i(A)\,a_x
	\end{Bmatrix}
	\in \mathbb{R}^3,
\end{equation}
il contributo del vincolo $k$ al vettore delle forze
generalizzate di $\mathcal{C}_i$ assume la forma compatta
\begin{equation}\label{eq:contrib_rea_vett}
	\bm{f}_i^{(k)} \coloneqq
	\begin{Bmatrix}
		R_k\,a_x \\[3pt]
		R_k\,a_y \\[3pt]
		R_k\,\bigl(x_i(A)\,a_y - y_i(A)\,a_x\bigr)
	\end{Bmatrix}
	= R_k\,\bm{b}_k^{(i)}.
\end{equation}
Il vettore $\bm{b}_k^{(i)}$ è stato costruito qui in modo del
tutto indipendente, a partire dalla sola definizione di momento
di una forza. Ma, per calcolo diretto, $\bm{b}_k^{(i)}$ coincide
componente per componente con il blocco $i$-esimo non nullo del
vettore $\bm{a}_k^{(i)}$ definito nella~\eqref{eq:ak_esterno} per
il problema cinematico: la stessa combinazione
$x_i(A)\,a_y - y_i(A)\,a_x$ vi compariva infatti come coefficiente
della rotazione $\vartheta_i$ nella proiezione su $\ab$ dello spostamento di $A$. Da qui in avanti scriviamo dunque semplicemente
$\bm{a}_k^{(i)}$ anche per il blocco statico, intendendo che si
tratta, per calcolo, dello stesso vettore già introdotto in
cinematica.
\medskip

Per un vincolo interno tra $\mathcal{C}_i$ e $\mathcal{C}_j$ lo
stesso ragionamento si applica a entrambi i corpi: la reazione agisce come una forza
$R_k\,\ab$ applicata in $A$ su $\mathcal{C}_j$ e come una forza
$-R_k\,\ab$ applicata nello stesso punto su $\mathcal{C}_i$.
Trasportando ciascuna delle due al proprio polo --- $O_j$ e $O_i$
rispettivamente --- con lo stesso argomento di rigidità appena
usato, si ottengono i contributi $+R_k\,\bm{a}_k^{(j)}$ e
$-R_k\,\bm{a}_k^{(i)}$, che riproducono esattamente la struttura
a due blocchi di segno opposto già incontrata
nella~\eqref{eq:ak_interno}.
\bigskip

Possiamo quindi ora formulare il problema statico in termini
algebrici. Dato un sistema vincolato, e assegnate le forze
attive applicate ai corpi --- descritte dal vettore delle
forze generalizzate $\bm{f} \in \mathbb{R}^n$ --- si
chiede di determinare, se esistono, le reazioni vincolari
$\bm{r} \in \mathbb{R}^m$ che, insieme alle forze attive,
soddisfano simultaneamente:
\begin{itemize}[noitemsep]
	\item le \emph{equazioni di equilibrio} di ciascun
	corpo $\mathcal{C}_i$, inteso come solido rigido
	soggetto alle forze attive e alle reazioni dei vincoli
	ad esso connessi;
	\item le \emph{condizioni di reazione}, che impongono
	che le reazioni siano orientate secondo i rispettivi
	versori di vincolo $\ab$.
\end{itemize}
Un sistema di reazioni che soddisfa entrambe le condizioni
si dice \emph{equilibrato}.
Assemblando i
contributi di tutte le $m$ reazioni, la condizione di
equilibrio del sistema esprime l'annullarsi della somma
del vettore delle forze generalizzate equivalenti alle
reazioni vincolari con quello delle forze attive:
\begin{equation}\label{eq:eq_globale}
	\bm{B}\,\bm{r} + \bm{f} = \bm{0},
\end{equation}
dove $\bm{B} \in \mathbb{R}^{n \times m}$ è la
\emph{matrice di equilibrio}\index{matrice di equilibrio|textbf}\index{matrice di equilibrio|seealso{matrice dei vincoli}}\nomenclature[L]{$\bm{B}$}{matrice di equilibrio del sistema} del sistema vincolato.
Il \emph{problema statico del sistema vincolato} è rappresentato da
un sistema lineare di $n = 3 n_c$ equazioni in $m$
incognite.

Riassumendo quanto appena osservato caso per caso --- per i
vincoli esterni nella~\eqref{eq:contrib_rea_vett}, per quelli
interni con lo stesso argomento applicato ai due corpi --- ogni
colonna $k$-esima di $\bm{B}$ coincide con la riga $k$-esima di
$\bm{A}$. In altre parole,
\begin{equation}
	\bm{B} = \bm{A}^{\top}.
\end{equation}
La matrice di equilibrio $\bm{B}$ è dunque l'aggiunta di
$\bm{A}$, e il problema statico è duale di quello cinematico.
Non è una coincidenza: è la conseguenza diretta del fatto che
la stessa geometria di vincolo --- la posizione del punto $A$
rispetto al polo $O_i$ e la direzione $\ab$ --- entra, tramite
lo stesso prodotto vettoriale $(A-O_i)\times\ab$, sia nella
componente di spostamento che il vincolo consente o proibisce
(lettura cinematica), sia nella componente di forza che la
reazione trasmette (lettura statica): la medesima geometria,
letta in un verso, descrive uno spostamento consentito; letta
nel verso opposto, descrive una forza trasmessa.

Lo schema del \CAPref{app:applicazionilin_dualita} (Figura~\ref{fig:duale}),
che rappresentava la dualità in chiave puramente
algebrica, si arricchisce ora di questo contenuto meccanico
preciso. Nella Figura~\ref{fig:duale_mec} le righe
corrispondono alle $m$ equazioni di vincolo e le colonne
alle $n$ componenti del vettore di configurazione
$\bm{u}$: ogni coefficiente $a_{ku_j}$ compare una volta
sola nella tabella, ma va letto in due modi distinti a
seconda della direzione di lettura. Letto \emph{per
	righe}, il coefficiente $a_{ku_j}$ è il peso della
componente $u_j$ nell'equazione di vincolo $k$-esima:
la riga intera è il vettore $\bm{a}_k^\top$ e il termine
noto è il cedimento $s_k$. Letto \emph{per colonne},
lo stesso $a_{ku_j}$ è il contributo della reazione $R_k$
all'equazione di equilibrio associata a $u_j$: la colonna
intera è la riga $j$-esima di $\bm{A}^\top$ e il termine
noto è $-X_j$.
\begin{figure}[htbp]
	\centering
	\renewcommand{\arraystretch}{1.3}
	\small
	\begin{tabular}{c|ccc|c|ccc|c}
		& $u_1$ & $v_1$ & $\vartheta_1$
		& $\cdots$ & $u_j$ & $v_j$ & $\vartheta_j$ & \\
		\hline
		$R_1$ & $a_{1u_1}$ & $a_{1v_1}$ & $a_{1\vartheta_1}$
		& $\cdots$ & $a_{1u_j}$ & $a_{1v_j}$ & $a_{1\vartheta_j}$ & $s_1$ \\
		$\vdots$ & $\vdots$ & $\vdots$ & $\vdots$
		& $\ddots$ & $\vdots$ & $\vdots$ & $\vdots$ & $\vdots$ \\
		$R_k$ & $a_{ku_1}$ & $a_{kv_1}$ & $a_{k\vartheta_1}$
		& $\cdots$ & $a_{ku_j}$ & $a_{kv_j}$ & $a_{k\vartheta_j}$ & $s_k$ \\
		$\vdots$ & $\vdots$ & $\vdots$ & $\vdots$
		& $\ddots$ & $\vdots$ & $\vdots$ & $\vdots$ & $\vdots$ \\
		$R_m$ & $a_{mu_1}$ & $a_{mv_1}$ & $a_{m\vartheta_1}$
		& $\cdots$ & $a_{mu_j}$ & $a_{mv_j}$ & $a_{m\vartheta_j}$ & $s_m$ \\
		\hline
		& $-X_1$ & $-Y_1$ & $-\mu(O_1)$
		& $\cdots$ & $-X_j$ & $-Y_j$ & $-\mu(O_j)$ &
	\end{tabular}
	\caption{Rilettura meccanica dello schema di dualità
		del \CAPref{app:applicazionilin_dualita}: per righe il problema cinematico
		$\bm{A}\,\bm{u}=\bm{s}$, per colonne il problema
		statico $\bm{B}\,\bm{r}=-\bm{f}$.}
	\label{fig:duale_mec}
\end{figure}
La Figura~\ref{fig:duale_mec} rende dunque visibile, in forma
schematica, la rilettura meccanica dello schema di dualità del
\CAPref{app:applicazionilin_dualita}: le grandezze astratte
$x_j$, $y_i$, $b_i$, $c_j$ sono qui sostituite dalle loro
controparti fisiche $u_j$, $R_k$, $s_k$, $-X_j$.
\begin{oss}
	Il significato del caso $a_{ku_j} = 0$ --- e
	analogamente per $a_{kv_j} = 0$ e
	$a_{k\vartheta_j} = 0$ --- è duplice e perfettamente
	simmetrico. In chiave \emph{cinematica}, $a_{ku_j} = 0$
	significa che la componente $u_j$ non compare
	nell'equazione di vincolo $k$-esima: il vincolo $k$ è
	\emph{insensibile} a quello spostamento, ovvero $u_j$
	è libero rispetto a quel vincolo. In chiave
	\emph{statica}, lo stesso $a_{ku_j} = 0$ significa che
	la reazione $R_k$ non contribuisce all'equazione di
	equilibrio associata a $u_j$: il vincolo $k$ non
	trasmette forza in quella direzione. Le due letture
	sono inseparabili: un vincolo che non sente uno
	spostamento non può trasmettergli una forza, e
	viceversa.
\end{oss}


% I risultati del \CAPref{app:applicazionilin_dualita}  si applicano
%direttamente: la soluzione esiste se e solo se
%$\bm{f} \in \Ic(\bm{B})$; se esiste, l'insieme delle
%soluzioni è un sottospazio affine di dimensione
%$g_i \coloneqq \nl(\bm{B}) = m - \rg(\bm{B})$,
%detto \emph{grado di iperstaticità}\index{iperstaticita@iperstaticità!grado di|textbf}\nomenclature[L]{$g_i$}{grado di iperstaticità del sistema ($\dim\Nc(\bm{A}^\top)$)} del sistema; la parte
%omogenea (con $\bm{f} = \bm{0}$) costituisce il
%$\Nc(\bm{B})$: lo spazio degli \emph{stati di
%	autoequilibrio}, cioè delle reazioni che equilibrano il
%sistema in assenza di forze attive.  Per ispezione diretta,
%confrontando i coefficienti della $k$-esima colonna di
%$\bm{B}$ --- calcolati nella~\eqref{eq:contrib_rea_vett}
%--- con i coefficienti della $k$-esima riga di $\bm{A}$
%--- dati dalla~\eqref{eq:ak_esterno} ---  si riconosce
%che le due coincidono: questo \emph{anticipa per ispezione}
%il dualismo $\bm{B} = \bm{A}^{\top}$, che nel \CAPref{cap:TLV}
%sarà ricondotto alla struttura bilineare del lavoro virtuale.



\paragraph{Il principio dei lavori virtuali.}\index{principio dei lavori virtuali!per un sistema di corpi rigidi|textbf}
Il problema statico appena introdotto ammette una riformulazione
equivalente in termini di lavoro virtuale, che generalizza al
sistema il principio dei lavori virtuali già stabilito per il
singolo corpo rigido (\CAPref{cap:statica_CR}): un corpo rigido
soggetto a un sistema di forze --- attive e reattive insieme ---
è equilibrato se e solo se il lavoro virtuale di tali forze è
nullo per ogni atto di moto rigido virtuale del corpo, comunque
scelto, senza alcun vincolo di congruenza.
\medskip

Applichiamo questo principio a ciascun corpo $\mathcal C_i$ del
sistema. Il risultante generalizzato delle forze --- attive e
reattive --- su $\mathcal C_i$ è, per costruzione, il blocco
$i$-esimo $(\bm B\bm r+\bm f)_i\in\mathbb R^3$ del vettore
assemblato nella~\eqref{eq:eq_globale}; il lavoro virtuale delle
forze su $\mathcal C_i$, per un atto di moto rigido virtuale
$\bm u_i\in\mathbb R^3$ del corpo, è dunque
\begin{equation}
	\Wc_i[\bm u_i]\coloneqq(\bm B\bm r+\bm f)_i\cdot \bm u_i.
\end{equation}
Per il principio dei lavori virtuali del corpo singolo,
\begin{equation}
	\mathcal C_i \text{ è equilibrato} \iff
	\Wc_i[\bm u_i] = 0
	\quad \forall\,\bm u_i\in\mathbb R^3.
\end{equation}
Per definizione (si veda il paragrafo precedente), il sistema è
equilibrato se e solo se ciascun corpo $\mathcal C_i$,
$i=1,\dots,n_c$, è individualmente equilibrato; dunque
\begin{equation}\label{eq:plv_percorpo}
	\text{sistema equilibrato} \iff
	\Wc_i[\bm u_i]=0 \quad \forall\, i=1,\dots,n_c,\
	\forall\,\bm u_i\in\mathbb R^3.
\end{equation}


\medskip

Resta da tradurre questa condizione, formulata corpo per corpo,
in un'unica condizione sull'intero sistema. Poniamo
$\bm u=(\bm u_1,\dots,\bm u_{n_c})\in\mathbb R^n$ --- un atto di
moto virtuale del sistema, cioè un atto di moto rigido virtuale
indipendente per ciascun corpo, non soggetto ad alcuna
condizione di congruenza --- e definiamo il lavoro virtuale
totale
\begin{equation}
	\Wc_{\rm tot}[\bm u] \coloneqq \sum_{i=1}^{n_c} \Wc_i[\bm u_i].
\end{equation}

($\Rightarrow$) Se $\Wc_i[\bm u_i]=0$ per ogni $i$ e ogni
$\bm u_i$, allora in particolare $\Wc_{\rm tot}[\bm u]=
\sum_i \Wc_i[\bm u_i]=0$ per ogni $\bm u\in\mathbb R^n$.

($\Leftarrow$) Viceversa, supponiamo $\Wc_{\rm tot}[\bm u]=0$
per ogni $\bm u\in\mathbb R^n$. Fissati un indice $i$ e un
$\bm u_i\in\mathbb R^3$ arbitrari, scegliamo $\bm u\in\mathbb
R^n$ nullo su tutti i blocchi tranne l'$i$-esimo, posto uguale
a $\bm u_i$: è una scelta lecita, perché $\bm u$ varia
liberamente su tutto $\mathbb R^n$. Poiché ciascun
$\Wc_j[\cdot]$ è lineare e $\Wc_j[\bm 0]=0$,
\[
0 = \Wc_{\rm tot}[\bm u] = \Wc_i[\bm u_i] + \sum_{j\neq i}\Wc_j[\bm 0]
= \Wc_i[\bm u_i].
\]
Essendo $i$ e $\bm u_i$ arbitrari, ciò mostra $\Wc_i[\bm u_i]=0$
per ogni $i$ e ogni $\bm u_i$.

Le due implicazioni mostrano che la~\eqref{eq:plv_percorpo}
equivale alla condizione unica $\Wc_{\rm tot}[\bm u]=0$ per ogni
$\bm u\in\mathbb R^n$. Possiamo dunque affermare che 

\bigskip

\emph{Un sistema di corpi rigidi è equilibrato se e solo se il lavoro
virtuale totale delle forze attive e reattive è nullo per ogni
spostamento rigido virtuale del sistema:
\begin{equation}\label{eq:plv_SCR}
	\bm B\bm r+\bm f=\bm 0 \iff \Wc_{\rm tot}[\bm u]=0
	\quad \forall\,\bm u\in\mathbb R^n.
\end{equation}
}
%\medskip
%
%\begin{oss}
%	Vale, come conseguenza immediata delle definizioni date,
%	l'identità
%	\begin{equation}
%		\Wc_{\rm tot}[\bm u] = (\bm B\bm r+\bm f)^{\!\top} \bm u
%		\qquad \forall\,\bm r,\bm f,\bm u,
%	\end{equation}
%	valida qualunque sia $\bm r$ --- non solo per le reazioni
%	equilibrate --- poiché $\Wc_i[\bm u_i]\coloneqq(\bm B\bm
%	r+\bm f)_i\cdot\bm u_i$ per definizione. Da questa forma il
%	principio dei lavori virtuali segue anche per via puramente
%	algebrica: un vettore $\bm g\in\mathbb R^n$ è nullo se e solo
%	se $\bm g^{\top}\bm u=0$ per ogni $\bm u\in\mathbb R^n$ (basta
%	scegliere $\bm u$ uguale, in successione, ai vettori della
%	base canonica); applicato a $\bm g=\bm B\bm r+\bm f$, questo
%	fatto elementare riproduce esattamente
%	la~\eqref{eq:plv_SCR}.
%\end{oss}

Notiamo infine che, a differenza del problema cinematico ---
dove ci si limita agli atti di moto congruenti, cioè tali che
$\bm A\bm u=\bm 0$ --- qui $\bm u$ varia su \emph{tutto}
$\mathbb R^n$, senza alcuna restrizione di congruenza: è
proprio questa libertà completa a rendere il principio un test
necessario e sufficiente per l'equilibrio. Sarà proprio questa
libertà, ristretta ora agli atti di moto congruenti, il punto di
partenza per confrontare equilibrio, congruenza e lavoro
virtuale nella \emph{triade meccanica} della prossima sezione.

\subsection{Equilibrio, congruenza e lavoro virtuale.}\label{sec:triade_SCR}\index{triade meccanica|textbf}
Il dualismo $\bm B=\bm A^\top$ permette di collegare tre
nozioni: l'\emph{equilibrio} di una coppia $(\bm r,\bm f)$, la \emph{congruenza} di una coppia
$(\bm u,\bm s)$ e il \emph{lavoro
	virtuale} reciproco fra un sistema di forze e uno di
spostamenti. Vedremo che equilibrio e congruenza, applicati a
due sistemi indipendenti, forzano sempre un'identità di lavoro
virtuale; il viceversa vale anch'esso, ma va enunciato con più
cura, perché nasconde un'insidia sui quantificatori che è bene
rendere esplicita.
\begin{oss}
	Dato $\bm u\in\mathbb R^n$, è sempre possibile trovare
	$\bm s$ tale che $(\bm u,\bm s)$ sia congruente: basta
	prendere $\bm s\coloneqq\bm A\bm u$. Non è vero il
	viceversa: dato $\bm s$, non sempre esiste $\bm u$ tale che
	$\bm A\bm u=\bm s$: occorre infatti che $\bm s\in\operatorname{Im}\bm
	A$. Analogamente, dato $\bm r\in\mathbb R^m$, è sempre
	possibile trovare $\bm f$ tale che $(\bm r,\bm f)$ sia
	equilibrata: basta prendere $\bm f\coloneqq-\bm B\bm r$.
	Non è vero il viceversa: dato $\bm f$, non sempre esiste
	$\bm r$ tale che $\bm B\bm r+\bm f=\bm 0$.
\end{oss}


\paragraph{L'Identità del lavoro virtuale tra stati indipendenti.}
Siano $(\bm r,\bm f)$ un sistema equilibrato e $(\bm u,\bm s)$
un sistema congruente, \textit{indipendenti tra loro}:
\[
\bm A\bm u=\bm s
\ \ (\text{congruenza}),
\qquad
\bm B\bm r+\bm f=\bm 0
\ \ (\text{equilibrio}).
\]
Allora
\begin{equation}\label{eq:reciprocita_SCR}
	\bm r^{\top}\bm s + \bm f^{\top}\bm u = 0.
\end{equation}
\begin{proof}
	Per \eqref{eq:plv_SCR}, essendo $(\bm r,\bm f)$ equilibrato,
	$\Wc_{\rm tot}[\bm u]=0$ per ogni $\bm u\in\mathbb R^n$, in
	particolare per il nostro $\bm u$. D'altra parte, per
	definizione,
	\[
	\Wc_{\rm tot}[\bm u]=(\bm B\bm r+\bm f)^\top\bm u
	=\bm r^\top\bm A\bm u+\bm f^\top\bm u
	=\bm r^\top\bm s+\bm f^\top\bm u,
	\]
	avendo usato $\bm B=\bm A^\top$ e $\bm s=\bm A\bm u$. Dunque
	$\bm r^\top\bm s+\bm f^\top\bm u=0$.
\end{proof}
\medskip
Schematicamente,
\[
\left.
\begin{aligned}
	(\bm u,\bm s) &\text{ congruente}\\
	(\bm r,\bm f) &\text{ equilibrato}
\end{aligned}
\right\}
\quad\Longrightarrow\quad
\bm r^\top\bm s+\bm f^\top\bm u=0.
\]
Nessun legame è richiesto fra i due sistemi: possono riferirsi
a due situazioni fisiche del tutto scorrelate e non è da pensare che gli spostamenti $\bm{u}$ siano in qualche modo conseguenti
alle forze esterne $\bm{f}$.

\paragraph{Congruenza e lavoro virtuale implicano l'equilibrio.}
L'identità appena dimostrata è una condizione \emph{necessaria}:
se $(\bm r,\bm f)$ è equilibrato, $\bm r^\top\bm s+\bm f^\top\bm
u=0$ vale per ogni sistema congruente $(\bm u,\bm s)$, scelto
indipendentemente da $(\bm r,\bm f)$. È naturale chiedersi se
questa condizione sia anche \emph{sufficiente}: se il lavoro
virtuale è nullo per ogni sistema congruente, possiamo concludere
che $(\bm r,\bm f)$ è equilibrato? La risposta è sì, ma vale la
pena capire con precisione perché, dato che il ``per ogni sistema
congruente'' nasconde un'insidia sui quantificatori. Per
l'osservazione iniziale, per ogni $\bm u\in\mathbb R^n$ esiste
sempre $\bm s\coloneqq\bm A\bm u$ che rende $(\bm u,\bm s)$
congruente: non c'è alcun $\bm u$ escluso. L'ipotesi equivale
dunque, semplicemente, a richiedere
\[
\bm r^\top\bm s+\bm f^\top\bm u=0
\qquad\forall\,\bm u\in\mathbb R^n\ \big(\text{posto }\bm s\coloneqq\bm A\bm u\big),
\]
cioè, per lo stesso calcolo appena fatto,
$\Wc_{\rm tot}[\bm u]=0$ per ogni $\bm u\in\mathbb R^n$. Ma
questa è precisamente l'ipotesi del principio dei lavori
virtuali \eqref{eq:plv_SCR}, che dà allora direttamente
$\bm B\bm r+\bm f=\bm 0$: il risultato è già acquisito, non
serve una nuova dimostrazione.


\paragraph{Equilibrio e lavoro virtuale implicano la congruenza.}
Chiediamoci ora se valga il viceversa, scambiando i ruoli di
$\bm u$ e $\bm r$: se $\bm r^\top\bm s+\bm f^\top\bm u=0$ vale
per ogni sistema equilibrato $(\bm r,\bm f)$, indipendentemente
scelto, possiamo concludere che $(\bm u,\bm s)$ è congruente?
Per l'osservazione iniziale, per ogni $\bm r\in\mathbb R^m$
esiste sempre $\bm f\coloneqq-\bm B\bm r$ che rende
$(\bm r,\bm f)$ equilibrato: l'ipotesi equivale dunque a
richiederla per ogni $\bm r\in\mathbb R^m$. Questa volta il
principio dei lavori virtuali non basta --- riguarda solo il
lato $\bm u$, non il lato $\bm r$ --- e la dimostrazione va
condotta autonomamente, in modo speculare, in $\mathbb R^m$:
\bigskip
\emph{
	Siano $\bm u\in\mathbb R^n$ e $\bm s\in\mathbb R^m$. Se
	\begin{equation}
		\bm r^\top\bm s+\bm f^\top\bm u=0
		\qquad\forall\,\bm r\in\mathbb R^m\ \big(\text{posto }\bm f\coloneqq-\bm A^\top\bm r\big),
	\end{equation}
	allora $\bm A\bm u=\bm s$.
}
\begin{proof}
	L'ipotesi si scrive $\bm r^\top(\bm s-\bm A\bm u)=0$ per
	ogni $\bm r\in\mathbb R^m$; vista la quantificazione, segue immediatamente che $\bm s-\bm A\bm u=\bm 0$.
\end{proof}

\bigskip

Le proposizioni appena dimostrate stabiliscono equivalenze fra
nozioni, ma non garantiscono da sole l'esistenza di una
soluzione per un dato assegnato, né dicono come calcolarla:
riguardano sistemi che già esistono entrambi. Le applicazioni
che seguono rispondono a queste due domande, e condividono
tutte la stessa origine: si applica l'identità
\eqref{eq:reciprocita_SCR} a una scelta particolare, non più
arbitraria, di uno dei due sistemi indipendenti. Si dividono
naturalmente in due famiglie speculari: restringendo il sistema
congruente $(\bm u,\bm s)$ si ottengono un criterio di esistenza
per l'equilibrio e una formula per il calcolo esplicito di una
reazione; restringendo, specularmente, il sistema equilibrato
$(\bm r,\bm f)$, si ottengono un criterio di esistenza per la
congruenza e una formula per il calcolo esplicito di uno
spostamento.

\paragraph{Compatibilità statica: l'esistenza dell'equilibrio.}
	\label{cor:compat_stat}
\emph{	Assegnato $\bm f\in\mathbb R^n$, il problema statico
	$\bm B\bm r+\bm f=\bm 0$ ammette soluzione se e solo se
	\begin{equation}\label{eq:compat_stat}
		\bm f^\top\bm u = 0\qquad\forall\,\bm u\in\Nc(\bm A),
	\end{equation}
	cioè le forze attive compiono lavoro virtuale nullo su
	ogni meccanismo del sistema.}
	
	\medskip

\begin{proof}
	($\Rightarrow$) Se esiste $\bm r$ con $\bm B\bm r+\bm f=\bm
	0$, l'identità~\eqref{eq:reciprocita_SCR}, applicata con
	$\bm s=\bm 0$ (cioè con $\bm u\in\Nc(\bm A)$), dà
	$\bm f^\top\bm u=0$ per ogni $\bm u\in\Nc(\bm A)$.
	
	($\Leftarrow$) Supponiamo $\bm f^\top\bm u=0$ per ogni
	$\bm u\in\Nc(\bm A)$, cioè $\bm f\in(\Nc(\bm A))^\perp$. Per
	il fatto di algebra lineare
	$(\Nc(\bm A))^\perp=\Ic(\bm A^\top)$ (si veda
	\CAPref{app:applicazionilin_dualita}), segue $\bm
	f\in\Ic(\bm A^\top)=\Ic\bm B$:
	esiste $\bm w\in\mathbb R^m$ con $\bm f=\bm B\bm w$; basta
	porre $\bm r\coloneqq-\bm w$.
\end{proof}
\begin{oss}
	Le due implicazioni del corollario hanno origine diversa.
	La necessità segue direttamente dall'identità
	\eqref{eq:reciprocita_SCR}, e vale sempre, senza altre
	ipotesi. La sufficienza richiede invece un fatto ulteriore,
	che non discende dal lavoro virtuale in sé, ma dalla
	struttura dello spazio vettoriale, che impone
$(\Nc(\bm A))^\perp=\Ic(\bm A^\top)$. Notiamo
	infine che la conclusione è un'esistenza: non dice quale
	sia $\bm r$, né se sia l'unico.
\end{oss}
\paragraph{Il teorema degli spostamenti virtuali (Müller-Breslau).}
\index{teorema!degli spostamenti virtuali|textbf}\index{TSV|see{teorema degli spostamenti virtuali}}\index{principio di Müller-Breslau|see{teorema degli spostamenti virtuali}}
È lo stesso procedimento già seguito, un vincolo alla volta, nella
\SECref{sec:gerber_equilibrio}: lì abbiamo sospeso concettualmente
un vincolo, imposto il lavoro totale nullo sul meccanismo
residuo, e ricavato così $Y_D$ (\eqref{eq:sopp_YD}) e $Y_B$
(\eqref{eq:sopp_YB}), senza risolvere il sistema di equazioni di
equilibrio corpo per corpo. Un'altra applicazione dell'identità
\eqref{eq:reciprocita_SCR} permette di formalizzare questo
procedimento per un sistema qualunque, trasformando il problema
statico in uno cinematico. Supponiamo che $\bm A$ sia suriettiva,
$\operatorname{rango}\bm A=m$ (nel linguaggio della
classificazione che vedremo, un sistema isostatico o labile:
assenza di autotensioni), cosicché per ogni $k$ esiste un atto
di moto che realizza un cedimento unitario al solo vincolo $k$;
e che $\bm f$ soddisfi la~\eqref{eq:compat_stat}, cosicché il
problema statico ammette soluzione $\bm r$. Vale dunque il seguente teorema.


\medskip

\emph{Sia $\bm e_k\in\mathbb R^m$ il $k$-esimo versore della base
	canonica, e sia $\bm u^*_k$ una soluzione del
	\emph{problema virtuale} $\bm A\bm u^*_k=\bm e_k$. La
	$k$-esima reazione si ottiene da
	\begin{equation}\label{eq:mueller_breslau}
		R_k = -\bm f^{\top}\bm u^*_k.
	\end{equation}
}

\medskip 

\begin{proof}
	Applichiamo l'identità \eqref{eq:reciprocita_SCR} alla
	coppia reale equilibrata $(\bm r,\bm f)$ e alla coppia
	virtuale congruente $(\bm u^*_k,\bm e_k)$:
	\[
	\bm r^\top\bm e_k+\bm f^\top\bm u^*_k=0.
	\]
	Il primo termine è $R_k$; riordinando si ottiene la tesi.
\end{proof}
\begin{oss}
	Se il sistema ammette autotensioni --- reazioni
	$\bm d\in\ker\bm B$ non nulle, capaci di autoequilibrarsi
	anche in assenza di carichi esterni: il fenomeno già
	annunciato nella \SECref{sec:motivazione_SCR}, e che chiameremo
	\emph{iperstaticità} quando classificheremo i sistemi ---
	la soluzione $\bm r$ del problema statico non è unica:
	eppure la~\eqref{eq:mueller_breslau} determina $R_k$ senza
	ambiguità, perché il secondo membro non dipende da $\bm r$.
	La ragione è che ogni autotensione $\bm d\in\ker\bm B$ ha
	automaticamente componente $k$-esima nulla: da $\bm
	A^\top\bm d=\bm 0$,
	\[
	d_k=\bm d^\top\bm e_k=\bm d^\top(\bm A\bm u_k^*)
	=(\bm A^\top\bm d)^\top\bm u_k^*=0.
	\]
	Il corollario trasforma così il calcolo di una singola
	reazione in un problema cinematico: si impone un cedimento
	unitario al vincolo $k$ e si calcola il lavoro delle forze
	attive nella configurazione che ne risulta. È il contenuto
	del \emph{teorema degli spostamenti virtuali} (TSV), noto
	nella tradizione come \emph{principio di Müller-Breslau}.
\end{oss}
\begin{oss}[Forma operativa del TSV]
	La scelta del cedimento virtuale $\bm s^*=\bm e_k$ usata
	nella dimostrazione è la più diretta dal punto di vista
	matriciale, ma non la più comoda nei conti. Assegnando il
	cedimento in verso \emph{opposto} al verso positivo della
	reazione --- cioè prendendo $\bm u^{**}_k$ soluzione di
	$\bm A\bm u^{**}_k=-\bm e_k$ --- l'identità
	\eqref{eq:reciprocita_SCR} dà direttamente
	\begin{equation}\label{eq:mueller_breslau_op}
		R_k=+\bm f^\top\bm u^{**}_k,
	\end{equation}
	senza cambio di segno. In questa forma $R_k$ coincide con
	il lavoro virtuale delle forze attive nello spostamento del
	meccanismo, ed è la scrittura che corrisponde al
	procedimento classico: \emph{cedere} il vincolo $k$-esimo
	di una quantità unitaria nel verso opposto alla reazione
	assunta positiva, e calcolare il lavoro delle forze attive
	nel moto rigido che ne risulta. Le due forme sono
	equivalenti; useremo sistematicamente
	la~\eqref{eq:mueller_breslau_op} negli esempi.
\end{oss}

Passiamo ora alla famiglia speculare, ottenuta restringendo il
sistema equilibrato $(\bm r,\bm f)$ anziché quello congruente.

\paragraph{Compatibilità cinematica: l'esistenza della congruenza.}
\index{compatibilita@compatibilità!cinematica|textbf}
	\label{cor:compat_cin}
	
	
\emph{	Assegnato $\bm s\in\mathbb R^m$, il problema cinematico
	$\bm A\bm u=\bm s$ ammette soluzione se e solo se
	\begin{equation}\label{eq:compat_cin}
		\bm s^\top\bm r = 0\qquad\forall\,\bm r\in\Nc(\bm A^\top),
	\end{equation}
	cioè i cedimenti compiono lavoro virtuale nullo su ogni
	stato di autoequilibrio del sistema.}
	
	\medskip

\begin{proof}
	($\Rightarrow$) Se esiste $\bm u$ con $\bm A\bm u=\bm s$,
	l'identità~\eqref{eq:reciprocita_SCR}, applicata con
	$\bm f=\bm 0$ (cioè con $\bm r\in\Nc(\bm A^\top)$, essendo
	l'equilibrio $\bm A^\top\bm r+\bm f=\bm 0$), dà
	$\bm r^\top\bm s=0$ per ogni $\bm r\in\Nc(\bm A^\top)$.
	
	($\Leftarrow$) Supponiamo $\bm s^\top\bm r=0$ per ogni
	$\bm r\in\Nc(\bm A^\top)$, cioè
	$\bm s\in(\Nc(\bm A^\top))^\perp$. Per il fatto di algebra
	lineare $(\Nc(\bm A^\top))^\perp=\Ic(\bm A)$
	(si veda \CAPref{app:applicazionilin_dualita}), segue
	$\bm s\in\Ic(\bm A)$: esiste $\bm u\in\mathbb
	R^n$ con $\bm A\bm u=\bm s$.
\end{proof}
\begin{oss}
	Come per la compatibilità statica, le due implicazioni
	hanno origine diversa: la necessità segue direttamente
	dall'identità \eqref{eq:reciprocita_SCR}, la sufficienza da
	un fatto di struttura dello spazio vettoriale
	($(\Nc(\bm A^\top))^\perp=\Ic(\bm A)$) --- il
	duale esatto di quello usato per la compatibilità statica.
	Anche qui la conclusione è un'esistenza: non dice quale sia
	$\bm u$, né se sia l'unico.
\end{oss}
\paragraph{Il teorema delle forze virtuali.}	\label{cor:TFV}
Specularmente al TSV, un'altra applicazione dell'identità
\eqref{eq:reciprocita_SCR} permette di calcolare esplicitamente
una componente dello spostamento, trasformando il problema
cinematico in uno statico. Supponiamo che $\bm A$ sia
iniettiva, $\operatorname{rango}\bm A=n$ (nel linguaggio della
classificazione che vedremo, assenza di meccanismi), cosicché
per ogni $j$ esiste un sistema di forze virtuali equilibrato
con risultante generalizzato pari al solo versore $\bm e_j$; e
che $\bm s$ soddisfi la~\eqref{eq:compat_cin}, cosicché il
problema cinematico ammette soluzione $\bm u$.

\index{teorema!delle forze virtuali|textbf}\index{TFV|see{teorema delle forze virtuali}}

\medskip
\emph{	Sia $\bm e_j\in\mathbb R^n$ il $j$-esimo versore della base
	canonica, e sia $\bm r^*_j$ una soluzione del
	\emph{problema virtuale} $\bm B\bm r^*_j+\bm e_j=\bm 0$. La
	$j$-esima componente di $\bm u$ si ottiene da
	\begin{equation}\label{eq:carico_unitario}
		u_j = -\bm s^{\top}\bm r^*_j.
	\end{equation}
}

\medskip

\begin{proof}
	Applichiamo l'identità \eqref{eq:reciprocita_SCR} alla
	coppia reale congruente $(\bm u,\bm s)$ e alla coppia
	virtuale equilibrata $(\bm r_j^*,\bm e_j)$:
	\[
	\bm r_j^{*\top}\bm s+\bm e_j^\top\bm u=0.
	\]
	Il secondo termine è $u_j$; riordinando si ottiene la tesi.
\end{proof}
\begin{oss}
	Il corollario è il duale esatto del TSV: il calcolo di una
	componente di spostamento si riconduce a un problema
	statico virtuale, applicando una forza generalizzata
	unitaria nella direzione corrispondente e calcolando il
	lavoro delle reazioni virtuali nei cedimenti reali
	assegnati. È il contenuto del \emph{teorema delle forze
		virtuali} (TFV).
\end{oss}
\begin{oss}[Spostamento generico, relativo e rotazione]
	Il corollario fornisce direttamente le componenti $u_j$
	del vettore di configurazione, associate ai poli $O_i$. La
	sua portata è però più generale, perché la scelta del
	sistema di forze virtuali $\bm f^*$ è arbitraria: quanto
	visto vale, senza alcuna modifica nella dimostrazione, per
	il calcolo di \emph{qualunque} grandezza cinematica che sia
	un funzionale lineare di $\bm u$, cioè della forma
	$\bm c^\top\bm u$ per un qualche $\bm c\in\mathbb R^n$ ---
	basta ripetere la dimostrazione con $\bm f^*=\bm c$ al
	posto di $\bm e_j$, ottenendo
	$\bm c^\top\bm u=-\bm s^\top\bm r^*$, dove $\bm r^*$ risolve
	$\bm B\bm r^*+\bm c=\bm 0$.
	
	\emph{Spostamento assoluto.} Per calcolare lo spostamento
	$\eta$ di un punto $P$ di $\mathcal C_i$ secondo una
	direzione $\ab$, si costruisce il problema virtuale con una
	forza fittizia unitaria applicata in $P$ nella direzione
	$\ab$: il vettore delle forze generalizzate corrispondente
	è esattamente il blocco già introdotto per quel punto e
	quella direzione (\eqref{eq:aki_def}), preso come $\bm c$.
	
	\emph{Spostamento relativo.} Per lo spostamento relativo fra
	due punti $P,Q$ di corpi distinti, secondo la congiungente,
	si applica una coppia di forze fittizie unitarie e opposte,
	una in $P$ e una in $Q$: $\bm c$ è la sovrapposizione dei
	due contributi.
	
	\emph{Rotazione relativa.} Per la rotazione relativa fra due
	corpi $\mathcal C_i$, $\mathcal C_j$, si applica una coppia
	di momenti fittizi unitari e opposti, uno su ciascun corpo.
	
	In tutti i casi la quantità cercata è il funzionale lineare
	duale del sistema di forze fittizie scelto: costruire il
	problema virtuale giusto è tutto ciò che serve.
\end{oss}



\paragraph{Uno sguardo d'insieme.}
Abbiamo usato l'espressione ``lavoro virtuale'' per otto
risultati diversi, tutti imparentati ma logicamente distinti:
vale la pena, prima di proseguire, mettere ordine.

\begin{center}
	\renewcommand{\arraystretch}{1.5}
	\resizebox{\textwidth}{!}{%
		\begin{tabular}{p{3.4cm}p{4.6cm}p{5.4cm}}
			\toprule
			Risultato & Sistemi coinvolti & Cosa afferma \\
			\midrule
			Principio dei lavori virtuali (PLV), \eqref{eq:plv_SCR}
			& un solo sistema $(\bm r,\bm f)$, testato contro
			\emph{ogni} $\bm u\in\mathbb R^n$, senza vincoli di
			congruenza
			& $\bm B\bm r+\bm f=\bm 0 \iff \Wc_{\rm tot}[\bm u]=0\
			\forall\bm u$ \\
			Identità del lavoro virtuale tra stati indipendenti,
			\eqref{eq:reciprocita_SCR}
			& due sistemi indipendenti: $(\bm r,\bm f)$ equilibrato,
			$(\bm u,\bm s)$ congruente
			& $\bm r^\top\bm s+\bm f^\top\bm u=0$ (conseguenza
			immediata del PLV) \\
			``Congruenza e lavoro virtuale implicano l'equilibrio''
			& come sopra, testato per \emph{ogni} $\bm u\in\mathbb
			R^n$
			& è di nuovo il PLV: nessun fatto nuovo \\
			``Equilibrio e lavoro virtuale implicano la congruenza''
			& come sopra, testato per \emph{ogni} $\bm r\in\mathbb
			R^m$
			& $\bm A\bm u=\bm s$  \\
			Compatibilità statica
			& l'identità ristretta a $\bm u\in\Nc(\bm A)$
			& esistenza di un $\bm r$ equilibrato, se $\bm
			f\perp\Nc(\bm A)$ \\
			Teorema degli spostamenti virtuali (TSV/Müller-Breslau)
			& l'identità con la scelta $\bm s=\bm e_k$
			& formula esplicita per la reazione $R_k$ \\
			Compatibilità cinematica
			& l'identità ristretta a $\bm r\in\Nc(\bm A^\top)$
			& esistenza di un $\bm u$ congruente, se $\bm
			s\perp\Nc(\bm A^\top)$ \\
			Teorema delle forze virtuali (TFV)
			& l'identità con la scelta $\bm f=\bm e_j$
			& formula esplicita per lo spostamento $u_j$ \\
			\bottomrule
		\end{tabular}%
	}
\end{center}

Un solo fatto algebrico genera l'intera famiglia: il dualismo
$\bm B=\bm A^\top$, di cui il PLV, dimostrato una volta per
tutte in \eqref{eq:plv_SCR}, è la traduzione nel linguaggio del
lavoro virtuale. È questo dualismo a rendere possibile il
passaggio fra il mondo statico e quello cinematico: l'identità
\eqref{eq:reciprocita_SCR} ne è la forma bilineare, valida per
una coppia qualunque di sistemi --- uno equilibrato, l'altro
congruente --- anche privi di ogni relazione fisica reciproca.
È proprio questa indifferenza alla relazione fra i due sistemi
a renderla utilizzabile come sonda: scegliendo di volta in
volta il sistema virtuale più conveniente, la stessa identità
restituisce ogni volta un'informazione diversa sull'altro
sistema, quello ``reale''.


Compatibilità statica e compatibilità cinematica sono risultati
di \emph{esistenza}: la prima stabilisce quando un carico $\bm
f$ ammette reazioni equilibrate --- se e solo se non compie
lavoro sui meccanismi del sistema, cioè su $\Nc(\bm A)$ --- e la
seconda, specularmente, quando un cedimento imposto $\bm s$
ammette un moto congruente --- se e solo se non compie lavoro
sulle autotensioni, cioè su $\Nc(\bm A^\top)$. TSV e TFV sono
invece risultati di \emph{calcolo}: scegliendo come sistema
virtuale non un intero sottospazio ma un singolo versore,
l'identità isola una singola incognita --- una reazione $R_k$,
uno spostamento $u_j$ --- e la restituisce come un unico
prodotto scalare, al prezzo di dover costruire una singola
configurazione virtuale ad hoc (un cedimento unitario,
un'autotensione unitaria) invece di risolvere per intero il
sistema di equazioni di equilibrio o di congruenza.

Meccanismi e autotensioni, e la loro dimensione, sono
precisamente l'oggetto della classificazione a cui ci dedichiamo
nel prossimo paragrafo.


%----------------------------------------------------------------------------------------------------------------------------
\section{I quattro sottospazi e la classificazione}\label{sec:quattro}
%----------------------------------------------------------------------------------------------------------------------------

I problemi cinematico e statico sono duali: la stessa matrice $\bm{A}$
governa entrambi, e i quattro sottospazi fondamentali\index{sottospazi fondamentali} di $\bm{A}$ e
$\bm{A}^\top$ determinano simultaneamente l'esistenza e la struttura
delle soluzioni di entrambi. Il teorema rango-nullità garantisce che
ciascuno dei due spazi si decompone come somma diretta ortogonale:
\begin{equation}
	\mathbb{R}^{n} = \Nc(\bm{A}) \oplus \Ic(\bm{A}^{\top}), \qquad
	\mathbb{R}^{m} = \Nc(\bm{A}^{\top}) \oplus \Ic(\bm{A}).
\end{equation}
Ogni vettore di configurazione $\bm{u}$ si decompone univocamente in
una parte ``meccanismo'' $\bm{u}_n \in \Nc(\bm{A})$ e una parte
``compatibile'' $\bm{u}_c \in \Ic(\bm{A}^\top)$; analogamente ogni
cedimento $\bm{s}$ si decompone in una parte ammissibile
$\bm{s}_c \in \Ic(\bm{A})$ e una parte incompatibile
$\bm{s}_n \in \Nc(\bm{A}^\top)$.


Le domande a cui questi quattro sottospazi rispondono sono già
emerse, in casi particolari, più volte nel corso del capitolo:
quando il problema cinematico ammette soluzione, e quante ne
ammette; quando il problema statico ammette soluzione, e quante.
Nei paragrafi che seguono li introduciamo uno a uno, ne raccogliamo il significato
meccanico in una sintesi tabellare,
ricaviamo da essi una classificazione completa dei sistemi di
corpi rigidi, e infine descriviamo
esplicitamente, caso per caso, la struttura delle soluzioni dei
due problemi.


\subsection{Struttura duale dei due problemi}\label{sec:dualita_cin_stat}

I problemi cinematico~\eqref{eq:prob_cin} e statico~\eqref{eq:prob_stat}
sono governati dalla stessa matrice $\bm{A}$, e collegati
dall'identità del lavoro virtuale tra stati indipendenti
\eqref{eq:reciprocita_SCR}, già dimostrata nella \SECref{sec:triade_SCR}.
%(i due addendi $\bm r^\top\bm s$ e $\bm f^\top\bm u$ hanno entrambi le
%dimensioni di un lavoro $[\text{N\,m}]$, pur appartenendo a spazi
%diversi: $\bm u,\bm f\in\mathbb R^n$, mentre $\bm s,\bm r\in\mathbb
%R^m$).

I quattro sottospazi coinvolti assumono significati meccanici precisi:
\begin{itemize}[nosep]
	\item $\Nc(\bm{A})$: \textbf{meccanismi} --- spostamenti $\bm{u}_n$ che
	non producono cedimenti ($\bm{A}\bm{u}_n = \bm{0}$);
	\item $\Ic(\bm{A}^\top)$: \textbf{forze equilibrabili} --- forze
	generalizzate $\bm{f}_c$ per cui esiste una reazione equilibrante,
	cioè un $\bm{r}$ tale che $\bm{A}^\top\bm{r}=-\bm{f}_c$;
	\item $\Nc(\bm{A}^\top)$: \textbf{stati di autoequilibrio}\index{stato di autoequilibrio|textbf}
	(dette anche \emph{autotensioni}) --- reazioni $\bm{r}_n$ la cui
	forza generalizzata equivalente è nulla ($\bm{A}^\top\bm{r}_n = \bm{0}$);
	\item $\Ic(\bm{A})$: \textbf{cedimenti ammissibili} --- cedimenti
	$\bm{s}_c$ effettivamente realizzabili da una configurazione del
	sistema, cioè tali che $\bm{s}_c=\bm{A}\bm{u}$ per qualche
	$\bm{u}\in\mathbb R^n$.
\end{itemize}

Per il teorema dei complementi ortogonali (\CAPref{app:applicazionilin_dualita}),
$\Ic(\bm A^\top)=\big(\Nc(\bm A)\big)^\perp$ e
$\Ic(\bm A)=\big(\Nc(\bm A^\top)\big)^\perp$: le condizioni di
Compatibilità statica~\eqref{eq:compat_stat} e Compatibilità
cinematica~\eqref{eq:compat_cin} --- $\bm f\perp\Nc(\bm A)$ e
$\bm s\perp\Nc(\bm A^\top)$, già dimostrate --- si riscrivono dunque
come
\begin{equation}\label{eq:compat_via_im}
	\bm f\in\Ic(\bm A^\top)
	\qquad\text{e}\qquad
	\bm s\in\Ic(\bm A)
\end{equation}
rispettivamente: il problema statico $\bm A^\top\bm r=-\bm f$ ammette
soluzione se e solo se $\bm f\in\Ic(\bm A^\top)$; il problema
cinematico $\bm A\bm u=\bm s$ ammette soluzione se e solo se
$\bm s\in\Ic(\bm A)$. È questa riformulazione, in termini di
appartenenza a un sottospazio anziché di ortogonalità a un altro, a
rendere immediato il passo successivo: contare le dimensioni.



La \FIGref{fig:problemi_duali_SCR} illustra questa struttura: i cedimenti
ammissibili $\bm{s}_c$ e le forze equilibrabili $\bm{f}_c$ vivono nelle immagini
dei rispettivi operatori, mentre le componenti incompatibili $\bm{s}_n$ e
$\bm{f}_n$ --- che devono annullarsi per garantire l'esistenza della soluzione
--- appartengono ai nuclei degli operatori duali.

%---------------------------------------------------------------------------------------------------------------------
\begin{figure}[htbp]
	\centering
	\begin{tikzpicture}[>=Stealth, line width=0.8pt, scale=0.8]
		
		% --- Ellisse sinistra: spazio R^n ---
		\draw[thick] (-3.5, 0) ellipse [x radius=2.2, y radius=3.0];
		\node at (-3.5, 3.4) {$\mathbb{R}^n$};
		
		% Curva divisoria sinistra
		\draw[thick] (-5.7, 0)
		.. controls (-4.5, 0.5) and (-2.5, -0.5) ..
		(-1.3, 0);
		
		% Regione superiore — N(A)
		\node          at (-3.5,  2.3) {$\Nc(\bm{A})$};
		\node[font=\small\itshape] at (-3.5,  1.65) {meccanismi};
		\fill (-3.5,  1.0) circle (1.5pt) node[right=2pt] {$\bm{u}_n$};
		\node[gray!60] at (-3.5,  0.35) {$\bm{f}_n = \bm{0}$};
		
		% Regione inferiore — Im(A^T)
		\fill (-3.5, -0.35) circle (1.5pt) node[right=2pt] {$\bm{u}_c$};
		\fill (-3.5, -1.0)  circle (1.5pt) node[right=2pt] {$\bm{f}_c$};
		\node[font=\small\itshape] at (-3.5, -1.65) {forze equilibrabili};
		\node          at (-3.5, -2.3) {$\Ic(\bm{A}^\top)$};
		
		% --- Ellisse destra: spazio R^m ---
		\draw[thick] (3.5, 0) ellipse [x radius=2.2, y radius=3.0];
		\node at (3.5, 3.4) {$\mathbb{R}^m$};
		
		% Curva divisoria destra
		\draw[thick] (1.3, 0)
		.. controls (2.5, 0.5) and (4.5, -0.5) ..
		(5.7, 0);
		
		% Regione superiore — Im(A)
		\node          at (3.5,  2.3) {$\Ic(\bm{A})$};
		\node[font=\small\itshape] at (3.5,  1.65) {ced.\ ammissibili};
		\fill (3.5,  1.0) circle (1.5pt) node[right=2pt] {$\bm{s}_c$};
		\fill (3.5,  0.35) circle (1.5pt) node[right=2pt] {$\bm{r}_c$};
		
		% Regione inferiore — N(A^T)
		\node[gray!60] at (3.5, -0.35) {$\bm{s}_n = \bm{0}$};
		\fill (3.5, -1.0)  circle (1.5pt) node[right=2pt] {$\bm{r}_n$};
		\node[font=\small\itshape] at (3.5, -1.65) {autoequilibrio};
		\node          at (3.5, -2.3) {$\Nc(\bm{A}^\top)$};
		
		% --- Freccia problema cinematico ---
		\draw[->, thick] (-1.0, 1.5) to[bend left=20]
		node[midway, above=4pt] {$\bm{A}\bm{u}_c = \bm{s}_c$}
		(1.0, 1.5);
		
		% --- Freccia problema statico ---
		\draw[->, thick] (1.0, -1.5) to[bend left=20]
		node[midway, below=4pt] {$\bm{A}^\top\bm{r}_c = -\bm{f}_c$}
		(-1.0, -1.5);
		
	\end{tikzpicture}
	\caption{Struttura duale dei problemi cinematico $\bm{A}\bm{u} = \bm{s}$
		e statico $\bm{A}^\top\bm{r} = -\bm{f}$. Il problema cinematico ammette
		soluzione se e solo se $\bm{s}_c \in \Ic(\bm{A})$, cioè $\bm{s}_n = \bm{0}$:
		il cedimento deve essere cinematicamente ammissibile. Il problema statico
		ammette soluzione se e solo se $\bm{f}_c \in \Ic(\bm{A}^\top)$, cioè
		$\bm{f}_n = \bm{0}$: la forza attiva deve essere equilibrabile dai vincoli.
		Le soluzioni generali sono $\bm{u} = \bm{u}_c + \bm{u}_n$ con
		$\bm{u}_n \in \Nc(\bm{A})$ (meccanismo) arbitrario, e
		$\bm{r} = \bm{r}_c + \bm{r}_n$ con $\bm{r}_n \in \Nc(\bm{A}^\top)$
		(stato di autoequilibrio) arbitrario.}
	\label{fig:problemi_duali_SCR}
\end{figure}


\begin{oss}
	La condizione $\bm{f}_n = \bm{0}$ --- la forza attiva deve essere
	equilibrabile --- è la condizione di equilibrio globale del sistema:
	essa richiede che la forza risultante e il momento risultante di tutte
	le forze attive siano bilanciabili dalle reazioni dei vincoli. Nei
	sistemi labili ($\Nc(\bm{A}) \neq \{\bm{0}\}$) questa condizione non è
	automaticamente soddisfatta e costituisce il criterio di compatibilità
	del problema statico. Dualmente, la condizione $\bm{s}_n = \bm{0}$ ---
	il cedimento deve essere cinematicamente ammissibile --- è automaticamente
	soddisfatta nei sistemi isostatici e labili (dove ogni configurazione è
	raggiungibile), ma impone restrizioni nei sistemi iperstatici
	($\Nc(\bm{A}^\top) \neq \{\bm{0}\}$).
\end{oss}

\medskip

La \TABref{tab:quattro_sottospazi} raccoglie i quattro sottospazi
fondamentali di $\bm{A}$  con la loro definizione e il
corrispondente significato meccanico.

\begin{table}[htbp]
	\centering
	\renewcommand{\arraystretch}{2.0}
	\small
	\begin{adjustbox}{max width=\textwidth}
		\begin{tabular}{
				>{\bfseries}p{3.4cm}
				p{1.5cm}
				p{4.2cm}
				p{4.3cm}}
			\toprule
			Sottospazio
			& Spazio
			& Definizione
			& Significato meccanico\\
			\midrule
			$\Nc(\bm{A})$
			& $\mathbb{R}^{n}$
			& $\bm{A}\,\bm{u} = \bm{0}$
			& \emph{Meccanismi}: spostamenti che non producono
			cedimenti\\
			$\Ic(\bm{A}^{\top})$
			& $\mathbb{R}^{n}$
			& $\Span\{\text{colonne di }\bm{A}^{\top}\}$
			& \emph{Forze equilibrabili}: $\bm{f}_c$ per cui esiste
			$\bm{r}$ con $\bm{A}^\top\bm{r}=-\bm{f}_c$\\
			\midrule
			$\Nc(\bm{A}^{\top})$
			& $\mathbb{R}^{m}$
			& $\bm{A}^{\top}\bm{r} = \bm{0}$
			& \emph{Stati di autoequilibrio} (autotensioni): reazioni
			la cui forza generalizzata equivalente è nulla\\
			$\Ic(\bm{A})$
			& $\mathbb{R}^{m}$
			& $\Span\{\text{colonne di }\bm{A}\}$
			& \emph{Cedimenti ammissibili}: $\bm{s}_c$ per cui esiste
			$\bm{u}$ con $\bm{s}_c=\bm{A}\bm{u}$\\
			\bottomrule
		\end{tabular}
	\end{adjustbox}
	\caption{Sintesi dei quattro sottospazi fondamentali di $\bm{A}$
		e del loro significato meccanico.
		I due sottospazi di $\mathbb{R}^{n}$ sono ortogonali
		($\Nc(\bm{A}) \perp \Ic(\bm{A}^{\top})$);
		idem per i due di $\mathbb{R}^{m}$
		($\Nc(\bm{A}^\top) \perp \Ic(\bm{A})$).}
	\label{tab:quattro_sottospazi}
\end{table}

\subsection{La formula fondamentale di classificazione}
\label{sec:classificazione}

Poniamo:
\begin{align}
	r &\coloneqq \rg(\bm{A}), \label{eq:rango}\\
	g_\ell &\coloneqq \nl(\bm{A}) = n - r
	\quad \text{(grado di labilità)},\label{eq:labilita}\\
	g_i &\coloneqq \nl(\bm{A}^{\top}) = m - r
	\quad \text{(grado di iperstaticità)}.\label{eq:iperstaticita}
\end{align}
Il teorema rango-nullità applicato a $\bm{A}$ e a
$\bm{A}^{\top}$ (che hanno lo stesso rango) dà direttamente
la \emph{formula fondamentale di classificazione}\index{classificazione dei sistemi di corpi rigidi!formula fondamentale}:
\begin{equation}\label{eq:formula_class}
g_\ell - g_i 
		= n - m \eqqcolon \nu,
\end{equation}
dove $\nu = n - m$\nomenclature[G]{$\nu$}{numero nominale di gradi di libertà ($n-m$)} è il \emph{numero nominale di gradi
	di libertà}, calcolabile a priori dal solo conteggio dei
corpi e dei vincoli scalari.

La~\eqref{eq:formula_class} mostra che $g_\ell$ e
$g_i$ non sono indipendenti: la loro differenza
è fissata dalla topologia del sistema ($n_c$ e $m$),
indipendentemente dalla geometria. La geometria può però
far sì che entrambi siano maggiori del minimo imposto da
$\nu$: si parla allora di \emph{degenerazione geometrica}\index{degenerazione geometrica|textbf}. La
classificazione che ne consegue è raccolta nella
\TABref{tab:classificazione_SCR}.




\begin{table}[htbp]
	\centering
	\renewcommand{\arraystretch}{1.8}
	\small
	\begin{tabular}{cccc p{5.5cm}}
		\toprule
		$\nu = n - m$ & $\rg(\bm{A})$
		& $g_\ell$ & $g_i$ & Tipo\\
		\midrule
		$= 0$ & $= m = n$ & $0$ & $0$
		& \textbf{Isostatico}: soluzione unica
		a entrambi i problemi\\
		$> 0$ & $= m$ & $\nu$ & $0$
		& \textbf{Labile} di grado $g_\ell = \nu$:
		il problema cinematico ha $\infty^{g_\ell}$
		soluzioni; il problema statico ammette
		soluzione solo se $\bm{f} \in \Ic(\bm{A}^{\top})$\\
		$< 0$ & $= n$ & $0$ & $|\nu|$
		& \textbf{Iperstatico} di grado $g_i = |\nu|$: il problema
		cinematico ammette soluzione solo se $\bm{s} \in \Ic(\bm{A})$
		(in tal caso unica); il problema statico ammette sempre
		soluzione, con $\infty^{g_i}$ possibilità\\
		\multirow{2}{*}{qualsiasi}
		& \multirow{2}{*}{$< \min(m, n)$}
		& \multirow{2}{*}{$> 0$}
		& \multirow{2}{*}{$> 0$}
		& \textbf{Labile-iperstatico}
		(degenerazione geometrica): entrambe le
		problematiche coesistono;
		$g_\ell - g_i = \nu$ resta valida\\
		\bottomrule
	\end{tabular}
	\caption{Classificazione dei sistemi di corpi rigidi piani in base	al rango di $\bm{A}$. Il caso labile-iperstatico si manifesta quando la geometria dei vincoli è degenere (ad esempio, vincoli 	tutti paralleli o tutti concorrenti in un punto). Ogni volta che 	$g_\ell > 0$, il problema statico è sovradeterminato e ammette	soluzione solo se $\bm{U}_q^{\top}\bm{f}=\bm{0}$, dove le		colonne di $\bm{U}_q$ formano una base di $\Nc(\bm{A})$; ogni	volta che $g_i > 0$, il problema cinematico è	sovradeterminato e ammette soluzione solo se $\bm{R}_\chi^{\top}\bm{s}=\bm{0}$, dove le colonne di		$\bm{R}_\chi$ formano una base di $\Nc(\bm{A}^\top)$. Entrambe	le matrici sono costruite esplicitamente nel	\PARref{sec:struttura_soluzione}.}
	\label{tab:classificazione_SCR}
\end{table}

\begin{oss}[Limiti del conteggio]
	La formula $\nu = n - m$ fornisce una condizione
	\emph{necessaria} ma non \emph{sufficiente} per
	l'isostaticità. Un sistema con $\nu = 0$ ma con vincoli
	geometricamente degeneri ha $g_\ell = g_i > 0$:
	è simultaneamente labile e iperstatico. Per stabilire
	il tipo effettivo occorre calcolare $\rg(\bm{A})$ o valutare in qualche modo
	$g_\ell$ o $g_i $.
	
\end{oss}


\subsection{Struttura della soluzione}
\label{sec:struttura_soluzione}

La classificazione della sezione precedente non riguarda solo
l'esistenza della soluzione ma anche la sua struttura. Per
ciascun caso si descrivono entrambi i problemi, distinguendo
la condizione di compatibilità, la soluzione particolare e
la parte omogenea.

\paragraph{Sistema isostatico ($g_\ell = 0$,
	$g_i = 0$).}
La matrice $\bm{A}$\index{isostaticita@isostaticità} è quadrata e invertibile. Entrambi i
problemi hanno soluzione unica per ogni dato:
\begin{equation}\label{eq:sol_iso}
	\bm{u} = \bm{A}^{-1}\bm{s},
	\qquad
	\bm{r} = -\bm{A}^{-\top}\bm{f}.
\end{equation}
Non vi è alcuna condizione di compatibilità.

\paragraph{Sistema labile ($g_\ell > 0$,
	$g_i = 0$).}
Sia $\bm{U}_q \in \mathbb{R}^{n \times g_\ell}$\index{labilita@labilità} la matrice
le cui colonne formano una base di $\Nc(\Abb)$:
\begin{equation}\label{eq:ker_A_lab}
	\bm{A}\,\bm{U}_q = \bm{0}.
\end{equation}

\noindent\emph{Problema cinematico} (sottodeterminato):
ammette sempre soluzione per ogni $\bm{s}$. La soluzione
generale è:
\begin{equation}\label{eq:sol_lab_cin}
	\bm{u} = \bm{u}_0 + \bm{U}_q\,\bm{q},
	\qquad \bm{q} \in \mathbb{R}^{g_\ell},
\end{equation}
dove $\bm{u}_0$ è una soluzione particolare di
$\bm{A}\,\bm{u}_0 = \bm{s}$ e $\bm{q}$ è il vettore
delle \emph{coordinate lagrangiane o di meccanismo}\index{coordinate lagrangiane}\nomenclature[L]{$\bm{q}$}{coordinate lagrangiane (di meccanismo)}, libero.

\noindent\emph{Problema statico} (sovradeterminato):
ammette soluzione se e solo se:
\begin{equation}\label{eq:compat_stat_lab}
	\bm{U}_q^{\top}\bm{f} = \bm{0},
\end{equation}
ovvero le forze attive non compiono lavoro in alcun
meccanismo. Se la~\eqref{eq:compat_stat_lab} è soddisfatta,
la soluzione è unica ($g_i = 0$):
\begin{equation}\label{eq:sol_lab_stat}
	\bm{r} = \bm{r}_0.
\end{equation}

\paragraph{Sistema iperstatico ($g_\ell = 0$,
	$g_i > 0$).}
Sia $\bm{R}_\chi \in \mathbb{R}^{m \times g_i}$\index{iperstaticita@iperstaticità}
la matrice le cui colonne formano una base di
$\Nc(\Abb^{\top})$:
\begin{equation}\label{eq:ker_AT_iper}
	\bm{A}^{\top}\bm{R}_\chi = \bm{0}.
\end{equation}

\noindent\emph{Problema cinematico} (sovradeterminato):
ammette soluzione se e solo se:
\begin{equation}\label{eq:compat_cin_iper}
	\bm{R}_\chi^{\top}\bm{s} = \bm{0},
\end{equation}
ovvero i cedimenti non compiono lavoro in alcuno stato
di autoequilibrio. Se la~\eqref{eq:compat_cin_iper} è
soddisfatta, la soluzione è unica ($g_\ell = 0$):
\begin{equation}\label{eq:sol_iper_cin}
	\bm{u} = \bm{u}_0.
\end{equation}

\noindent\emph{Problema statico} (sottodeterminato):
ammette sempre soluzione per ogni $\bm{f}$. La soluzione
generale è:
\begin{equation}\label{eq:sol_iper_stat}
	\bm{r} = \bm{r}_0 + \bm{R}_\chi\,\bm{\chi},
	\qquad \bm{\chi} \in \mathbb{R}^{g_i},
\end{equation}
dove $\bm{r}_0$ è una soluzione particolare di
$\bm{A}^{\top}\bm{r}_0 = -\bm{f}$ e $\bm{\chi}$ è il
vettore delle \emph{incognite iperstatiche}\index{vettore delle incognite iperstatiche}\nomenclature[G]{$\bm{\chi}$}{vettore delle incognite iperstatiche}, libero.

\paragraph{Sistema labile-iperstatico ($g_\ell > 0$,
	$g_i > 0$).}
Entrambe le matrici $\bm{U}_q$ e $\bm{R}_\chi$ sono non
banali.

\noindent\emph{Problema cinematico}: ammette soluzione
se e solo se $\bm{R}_\chi^{\top}\bm{s} = \bm{0}$;
in tal caso la soluzione generale è:
\begin{equation}\label{eq:sol_li_cin}
	\bm{u} = \bm{u}_0 + \bm{U}_q\,\bm{q},
	\qquad \bm{q} \in \mathbb{R}^{g_\ell}.
\end{equation}

\noindent\emph{Problema statico}: ammette soluzione
se e solo se $\bm{U}_q^{\top}\bm{f} = \bm{0}$;
in tal caso la soluzione generale è:
\begin{equation}\label{eq:sol_li_stat}
	\bm{r} = \bm{r}_0 + \bm{R}_\chi\,\bm{\chi},
	\qquad \bm{\chi} \in \mathbb{R}^{g_i}.
\end{equation}


\begin{oss}
	Le condizioni~\eqref{eq:compat_stat_lab} e~\eqref{eq:compat_cin_iper}
	non sono altro che Compatibilità statica~\eqref{eq:compat_stat} e
	Compatibilità cinematica~\eqref{eq:compat_cin}, già dimostrate nella
	\SECref{sec:triade_SCR}, riscritte in forma matriciale tramite una
	base del nucleo: $\bm f^\top\bm u=0$ per ogni $\bm u\in\Nc(\bm A)$
	equivale a $\bm U_q^\top\bm f=\bm 0$ perché le colonne di $\bm U_q$
	generano $\Nc(\bm A)$; analogamente $\bm s^\top\bm r=0$ per ogni
	$\bm r\in\Nc(\bm A^\top)$ equivale a $\bm R_\chi^\top\bm s=\bm 0$.
	Per la stessa ragione, esse coincidono anche con la riformulazione
	in termini di immagine vista in~\eqref{eq:compat_via_im}:
	$\bm U_q^\top\bm f=\bm 0 \iff \bm f\in\Ic(\bm A^\top)$ e
	$\bm R_\chi^\top\bm s=\bm 0 \iff \bm s\in\Ic(\bm A)$. Sono tre
	scritture della stessa condizione, utili in contesti diversi: la
	prima per il calcolo esplicito (nota una base del nucleo), la
	seconda per il significato fisico (lavoro virtuale nullo su
	meccanismi o autotensioni), la terza per il conteggio dimensionale
	usato nella classificazione.
\end{oss}










%%-------------------------------------------------------------------------------------------------------------------
\section{Esempi applicativi}\label{sec:esempi}
%%-------------------------------------------------------------------------------------------------------------------


\begin{esempio}[Trave semplicemente appoggiata]
	\label{ex:trave_semplice}
	Si consideri una trave rigida di lunghezza $\ell$,
	vincolata da una cerniera in $A = (0,0)$ e da un
	carrello verticale in $B = (\ell, 0)$. Si sceglie
	il polo $O \equiv A$.
	
	\noindent\textit{Problema cinematico.} Dato il
	cedimento verticale $\bar{v}_B > 0$ al carrello in
	$B$ (tutti gli altri vincoli perfetti), determinare
	la configurazione della trave.
	
	\noindent\textit{Problema statico.} Data una forza
	concentrata $F$ (verso il basso) applicata nel punto
	$x = a$ ($0 < a < \ell$), determinare le reazioni
	vincolari.
	
	I dati dei due problemi sono illustrati nella
	Figura~\ref{fig:trave_semplice}.
	
	\begin{figure}[ht]
		\centering
		\begin{tikzpicture}[>=Stealth, line width=0.8pt, font=\small,
			nodo/.style={circle, fill=white, draw=black,
				inner sep=0pt, minimum size=5pt, line width=0.8pt}, tacca/.style={thin, gray!70}]
			
			
			\def\L{5.0}
			\def\tw{0.22}
			\def\thA{0.38}
			\def\thR{0.32}
			\def\bw{0.16}
			\def\gw{0.30}
			\def\dx{0.6}    % offset orizzontale del cedimento
			
			% ============================================================
			% PANNELLO (a) — dato cinematico
			% ============================================================
			\begin{scope}[xshift=0cm]
				
				%% trave
				\draw[line width=2.2pt] (0,0) -- (\L,0);
				
				%% etichette nodi
				\node[above=2pt] at (0,0)  {$A$};
				\node[above=2pt] at (\L,0) {$B$};
				
				%% cerniera in A
				\draw (0,0) -- ++(-\tw,-\thA) -- ++(2*\tw,0) -- cycle;
				\draw (-\gw,-\thA) -- (\gw,-\thA);
				\foreach \x in {-0.24,-0.12,0,0.12,0.24}
				\draw[thin] (\x,-\thA) -- ++(-0.09,-0.11);
				\node[nodo] at (0,0) {};
				
				%% carrello terminale in B
				\draw (\L,0) -- ++(-\tw,-\thR) -- ++(2*\tw,0) -- cycle;
				\draw (\L-\gw,-\thR) -- (\L+\gw,-\thR);
				\draw (\L-\bw,-\thR-0.12) circle (2.4pt);
				\draw (\L+\bw,-\thR-0.12) circle (2.4pt);
				\draw (\L-\gw,-\thR-0.24) -- (\L+\gw,-\thR-0.24);
				\foreach \x in {-0.24,-0.12,0,0.12,0.24}
				\draw[thin] (\L+\x,-\thR-0.24) -- ++(-0.09,-0.11);
				\node[nodo] at (\L,0) {};
				
				%% cedimento di B (a destra: tacca orizzontale + freccia verticale)
				\draw[thin] (\L,0) -- (\L+\dx,0);
				\draw[->, thick] (\L+\dx,0) -- (\L+\dx,-0.50)
				node[right=3pt, pos=0.5] {$\bar{v}_B$};
				
				%% quota
				\draw[<->, thin] (0,-0.90) -- (\L,-0.90)
				node[midway, above=2pt] {$\ell$};
				
				\draw[tacca] (0, -1.00) -- (0, -0.80);
				\draw[tacca] (\L,-1.00) -- (\L,-0.80);
				
				
				
				\node at (\L/2, -1.40) {(a)};
				
			\end{scope}
			
			% ============================================================
			% PANNELLO (b) — dato statico
			% ============================================================
			\begin{scope}[xshift=7.5cm]
				
				\def\xa{3.0}
				
				%% trave
				\draw[line width=2.2pt] (0,0) -- (\L,0);
				
				%% etichette nodi
				\node[above=2pt] at (0,0)  {$A$};
				\node[above=2pt] at (\L,0) {$B$};
				
				%% cerniera in A
				\draw (0,0) -- ++(-\tw,-\thA) -- ++(2*\tw,0) -- cycle;
				\draw (-\gw,-\thA) -- (\gw,-\thA);
				\foreach \x in {-0.24,-0.12,0,0.12,0.24}
				\draw[thin] (\x,-\thA) -- ++(-0.09,-0.11);
				\node[nodo] at (0,0) {};
				
				%% carrello terminale in B
				\draw (\L,0) -- ++(-\tw,-\thR) -- ++(2*\tw,0) -- cycle;
				\draw (\L-\gw,-\thR) -- (\L+\gw,-\thR);
				\draw (\L-\bw,-\thR-0.12) circle (2.4pt);
				\draw (\L+\bw,-\thR-0.12) circle (2.4pt);
				\draw (\L-\gw,-\thR-0.24) -- (\L+\gw,-\thR-0.24);
				\foreach \x in {-0.24,-0.12,0,0.12,0.24}
				\draw[thin] (\L+\x,-\thR-0.24) -- ++(-0.09,-0.11);
				\node[nodo] at (\L,0) {};
				
				%% forza F
				\draw[thin] (\xa,0.05) -- (\xa,-0.05);
				\draw[->, thick] (\xa,0.85) -- (\xa,0.10)
				node[right=3pt, pos=0] {$F$};
				
				%% quotatura: parziali (a, ell-a)
				\draw[<->, thin] (0,-0.90) -- (\xa,-0.90)
				node[midway, above=2pt] {$a$};
				\draw[<->, thin] (\xa,-0.90) -- (\L,-0.90)
				node[midway, above=2pt] {$\ell - a$};
				
				\draw[tacca] (0,  -1.00) -- (0,  -0.80);
				\draw[tacca] (\xa,-1.00) -- (\xa,-0.80);
				\draw[tacca] (\L, -1.00) -- (\L, -0.80);
				
				\node at (\L/2, -1.40) {(b)};
				
			\end{scope}
			
		\end{tikzpicture}
		\caption{Trave semplicemente appoggiata
			(\ESEref{ex:trave_semplice}): dato cinematico (a) e dato
			statico (b). In entrambi i casi la cerniera è in $A=(0,0)$
			e il carrello verticale in $B=(\ell,0)$, con polo
			$O\equiv A$.}
		\label{fig:trave_semplice}
	\end{figure}
	
	\begin{svolgimento}
		
		\noindent\textbf{Vettore di configurazione e matrice
			dei vincoli.}
		Con $O = A$, $n_c = 1$, $n = 3$:
		\begin{equation*}
			\bm{u} =
			\bigl\{u(O),\; v(O),\; \vartheta
			\bigr\}^\top \in \mathbb{R}^3.
		\end{equation*}
		I vincoli scalari sono $m = 3$: componente $x$ e
		componente $y$ della cerniera in $A$, e componente
		verticale del carrello in $B$. La matrice dei vincoli
		è:
		\begin{equation*}\label{eq:A_es1}
			\bm{A} =
			\begin{bmatrix}
				1 & 0 & 0 \\
				0 & 1 & 0 \\
				0 & 1 & \ell
			\end{bmatrix}.
		\end{equation*}
		
		\noindent\textbf{Classificazione.}
		$\nu = n - m = 0$ e $\det\bm{A} = \ell \neq 0$:
		il sistema è \textbf{isostatico}.
		
		\smallskip
		\noindent\textbf{Problema cinematico:
			$\bm{A}\,\bm{u} = \bm{s}$.}
		Il vettore dei cedimenti è:
		\begin{equation*}
			\bm{s} =
			\bigl\{0,\; 0,\; -\bar{v}_B
			\bigr\}^\top,
		\end{equation*}
		dove il segno negativo indica che il cedimento è
		verso il basso mentre l'asse del carrello è $\ab_y$.
		Il sistema $\bm{A}\,\bm{u} = \bm{s}$ dà:
		\begin{equation*}
			u(O) = 0, \qquad v(O) = 0, \qquad
			\vartheta = -\frac{\bar{v}_B}{\ell}.
		\end{equation*}
		La trave ruota rigidamente attorno ad $A$ di ampiezza
		$\bar{v}_B/\ell$ in senso orario.
		
		\smallskip
		\noindent\textbf{Problema statico:
			$\bm{A}^{\top}\bm{r} = -\bm{f}$.}
		Le incognite sono le reazioni vincolari:
		\begin{equation*}
			\bm{r} =
			\bigl\{X_A,\; Y_A,\; Y_B
			\bigr\}^\top.
		\end{equation*}
		Il vettore delle forze generalizzate riferite al polo
		$O \equiv A$ è:
		\begin{equation*}
			\bm{f} =
			\bigl\{0,\; -F,\; -Fa
			\bigr\}^\top,
		\end{equation*}
		dove $-F$ è la risultante verticale e $-Fa$ è il
		momento rispetto ad $A$. La matrice del problema
		statico è:
		\begin{equation*}\label{eq:B_es1}
			\bm{A}^{\top} =
			\begin{bmatrix}
				1 & 0 & 0 \\
				0 & 1 & 1 \\
				0 & 0 & \ell
			\end{bmatrix}.
		\end{equation*}
		Il sistema $\bm{A}^{\top}\bm{r} = -\bm{f}$ si
		risolve per sostituzione all'indietro:
		\begin{equation*}
			Y_B = \frac{Fa}{\ell}, \qquad
			Y_A = F\left(1 - \frac{a}{\ell}\right), \qquad
			X_A = 0.
		\end{equation*}
		
		\noindent\textbf{Dualismo.}
		Il coefficiente $\ell$ che compare nella terza riga
		di $\bm{A}$ --- il braccio del carrello in $B$
		rispetto al polo $O$ nella prestazione cinematica
		--- è lo stesso che compare nella terza riga di
		$\bm{A}^\top$ --- il braccio della reazione $Y_B$
		nel momento di equilibrio rispetto ad $A$. La stessa
		geometria di vincolo governa entrambi i problemi.
	\end{svolgimento}
\end{esempio}

%\begin{esempio}[Trave di Gerber]
%	\label{ex:gerber}
%	Si consideri un sistema di due travi rigide
%	$\mathcal{C}_1$ e $\mathcal{C}_2$ collegate da una
%	cerniera interna in $C = (3a, 0)$. I dati geometrici
%	sono:
%	\begin{itemize}[noitemsep]
%		\item $\mathcal{C}_1$: polo $O_1 \equiv A$,
%		cerniera al suolo in $A$, carrello verticale in
%		$B$, cerniera interna in $C$;
%		\item $\mathcal{C}_2$: polo $O_2 \equiv C$,
%		carrello verticale in $D $.
%	\end{itemize}
%	
%	\noindent\textit{Problema cinematico.} Sono assegnati
%	un cedimento verticale $2\bar{v}$ verso il basso al
%	carrello in $B$ e un cedimento verticale $\bar{v}$
%	verso il basso al carrello in $D$. Determinare la
%	configurazione del sistema.
%	
%	\noindent\textit{Problema statico.} Sono applicate
%	una forza $2F$ verso il basso nel punto medio del
%	tratto $AB$ (a distanza $a$ da $A$) su $\mathcal{C}_1$
%	e una forza $F$ verso il basso nel punto medio di
%	$\mathcal{C}_2$ (a distanza $a$ da $C$) su
%	$\mathcal{C}_2$. Determinare le reazioni vincolari.
%	
%	I dati dei due problemi sono illustrati nella
%	Figura~\ref{fig:gerber}.
%	
%	\begin{figure}[ht]
%		\centering
%		\begin{tikzpicture}[>=Stealth, line width=0.8pt, font=\small,
%			nodo/.style={circle, fill=white, draw=black,
%				inner sep=0pt, minimum size=5pt, line width=0.8pt},
%			cint/.style={circle, fill=white, draw=black,
%				inner sep=0pt, minimum size=6pt, line width=0.8pt},
%			richiam/.style={thin, gray, dashed},
%			tacca/.style={thin, gray!70}]
%			
%			
%			\def\a{1.0}     % unità di lunghezza (= a)
%			\def\tw{0.22}   % semi-larghezza triangolo
%			\def\thA{0.38}  % altezza triangolo cerniera
%			\def\thR{0.32}  % altezza triangolo carrello
%			\def\sr{0.08}   % raggio mezzaluna (carrello intermedio)
%			\def\bw{0.16}   % semi-distanza rotelle
%			\def\gw{0.30}   % semi-larghezza barra di terra
%			
%			% ============================================================
%			% PANNELLO (a) — dato cinematico
%			% ============================================================
%			\begin{scope}[xshift=0cm]
%				
%				%% trave (suddivisa in C1 e C2 dalla cerniera in C)
%				\draw[line width=2.2pt] (0,0) -- ({3*\a},0);
%				\draw[line width=2.2pt] ({3*\a},0) -- ({5*\a},0);
%				
%				%% etichette nodi
%				\node[above=2pt] at (0,0)       {$A$};
%				\node[above=2pt] at ({2*\a},0)  {$B$};
%				\node[above=2pt] at ({3*\a},0)  {$C$};
%				\node[above=2pt] at ({5*\a},0)  {$D$};
%				
%				%% etichette corpi
%				\node at ({\a},  -0.30) {$\mathcal{C}_1$};
%				\node at ({4*\a},-0.30) {$\mathcal{C}_2$};
%				
%				%% cerniera in A
%				\draw (0,0) -- ++(-\tw,-\thA) -- ++(2*\tw,0) -- cycle;
%				\draw (-\gw,-\thA) -- (\gw,-\thA);
%				\foreach \x in {-0.24,-0.12,0,0.12,0.24}
%				\draw[thin] (\x,-\thA) -- ++(-0.09,-0.11);
%				\node[nodo] at (0,0) {};
%				
%				%% carrello intermedio in B
%				\draw ({2*\a},0) -- ++(-\tw,-\thR) -- ++(2*\tw,0) -- cycle;
%				\draw ({2*\a-\gw},-\thR) -- ({2*\a+\gw},-\thR);
%				\draw ({2*\a-\bw},-\thR-0.12) circle (2.4pt);
%				\draw ({2*\a+\bw},-\thR-0.12) circle (2.4pt);
%				\draw ({2*\a-\gw},-\thR-0.24) -- ({2*\a+\gw},-\thR-0.24);
%				\foreach \x in {-0.24,-0.12,0,0.12,0.24}
%				\draw[thin] ({2*\a+\x},-\thR-0.24) -- ++(-0.09,-0.11);
%				\fill[white] ({2*\a-\sr},0) arc(180:360:\sr) -- cycle;
%				\draw ({2*\a-\sr},0) arc(180:360:\sr);
%				
%				%% cerniera interna in C
%				\node[cint] at ({3*\a},0) {};
%				
%				%% carrello terminale in D
%				\draw ({5*\a},0) -- ++(-\tw,-\thR) -- ++(2*\tw,0) -- cycle;
%				\draw ({5*\a-\gw},-\thR) -- ({5*\a+\gw},-\thR);
%				\draw ({5*\a-\bw},-\thR-0.12) circle (2.4pt);
%				\draw ({5*\a+\bw},-\thR-0.12) circle (2.4pt);
%				\draw ({5*\a-\gw},-\thR-0.24) -- ({5*\a+\gw},-\thR-0.24);
%				\foreach \x in {-0.24,-0.12,0,0.12,0.24}
%				\draw[thin] ({5*\a+\x},-\thR-0.24) -- ++(-0.09,-0.11);
%				\node[nodo] at ({5*\a},0) {};
%				
%				%% cedimenti
%				\draw[->, thick] ({2*\a},-0.75) -- ({2*\a},-1.30)
%				node[right=3pt, pos=0.5] {$2\bar{v}$};
%				\draw[->, thick] ({5*\a},-0.75) -- ({5*\a},-1.30)
%				node[right=3pt, pos=0.5] {$\bar{v}$};
%				
%				%% quotatura
%				\draw[richiam] (0,      -\thA-0.15) -- (0,      -1.55);
%				\draw[richiam] ({2*\a}, -1.40)      -- ({2*\a}, -1.55);
%				\draw[richiam] ({3*\a}, -0.10)      -- ({3*\a}, -1.55);
%				\draw[richiam] ({5*\a}, -1.40)      -- ({5*\a}, -1.55);
%				
%				\draw[tacca] (0,     -1.55) -- (0,     -1.75);
%				\draw[tacca] ({2*\a},-1.55) -- ({2*\a},-1.75);
%				\draw[tacca] ({3*\a},-1.55) -- ({3*\a},-1.75);
%				\draw[tacca] ({5*\a},-1.55) -- ({5*\a},-1.75);
%				
%				\draw[<->, thin, gray] (0,     -1.65) -- ({2*\a},-1.65);
%				\node[above=0.04] at ({\a},   -1.65) {$2a$};
%				\draw[<->, thin, gray] ({2*\a},-1.65) -- ({3*\a},-1.65);
%				\node[above=0.04] at ({2.5*\a},-1.65) {$a$};
%				\draw[<->, thin, gray] ({3*\a},-1.65) -- ({5*\a},-1.65);
%				\node[above=0.04] at ({4*\a}, -1.65) {$2a$};
%				
%				\node at ({2.5*\a}, -2.05) {(a)};
%				
%			\end{scope}
%			
%			% ============================================================
%			% PANNELLO (b) — dato statico
%			% ============================================================
%			\begin{scope}[xshift=7cm]
%				
%				\draw[line width=2.2pt] (0,0) -- ({3*\a},0);
%				\draw[line width=2.2pt] ({3*\a},0) -- ({5*\a},0);
%				
%				\node[above=2pt] at (0,0)       {$A$};
%				\node[above=2pt] at ({2*\a},0)  {$B$};
%				\node[above=2pt] at ({3*\a},0)  {$C$};
%				\node[above=2pt] at ({5*\a},0)  {$D$};
%				
%				\node at ({\a},  -0.30) {$\mathcal{C}_1$};
%				\node at ({4*\a},-0.30) {$\mathcal{C}_2$};
%				
%				%% cerniera in A
%				\draw (0,0) -- ++(-\tw,-\thA) -- ++(2*\tw,0) -- cycle;
%				\draw (-\gw,-\thA) -- (\gw,-\thA);
%				\foreach \x in {-0.24,-0.12,0,0.12,0.24}
%				\draw[thin] (\x,-\thA) -- ++(-0.09,-0.11);
%				\node[nodo] at (0,0) {};
%				
%				%% carrello intermedio in B
%				\draw ({2*\a},0) -- ++(-\tw,-\thR) -- ++(2*\tw,0) -- cycle;
%				\draw ({2*\a-\gw},-\thR) -- ({2*\a+\gw},-\thR);
%				\draw ({2*\a-\bw},-\thR-0.12) circle (2.4pt);
%				\draw ({2*\a+\bw},-\thR-0.12) circle (2.4pt);
%				\draw ({2*\a-\gw},-\thR-0.24) -- ({2*\a+\gw},-\thR-0.24);
%				\foreach \x in {-0.24,-0.12,0,0.12,0.24}
%				\draw[thin] ({2*\a+\x},-\thR-0.24) -- ++(-0.09,-0.11);
%				\fill[white] ({2*\a-\sr},0) arc(180:360:\sr) -- cycle;
%				\draw ({2*\a-\sr},0) arc(180:360:\sr);
%				
%				\node[cint] at ({3*\a},0) {};
%				
%				%% carrello terminale in D
%				\draw ({5*\a},0) -- ++(-\tw,-\thR) -- ++(2*\tw,0) -- cycle;
%				\draw ({5*\a-\gw},-\thR) -- ({5*\a+\gw},-\thR);
%				\draw ({5*\a-\bw},-\thR-0.12) circle (2.4pt);
%				\draw ({5*\a+\bw},-\thR-0.12) circle (2.4pt);
%				\draw ({5*\a-\gw},-\thR-0.24) -- ({5*\a+\gw},-\thR-0.24);
%				\foreach \x in {-0.24,-0.12,0,0.12,0.24}
%				\draw[thin] ({5*\a+\x},-\thR-0.24) -- ++(-0.09,-0.11);
%				\node[nodo] at ({5*\a},0) {};
%				
%				%% forze
%				\draw[thin] ({\a},0.05) -- ({\a},-0.05);
%				\draw[->, thick] ({\a},0.70) -- ({\a},0.10)
%				node[right=3pt, pos=0] {$2F$};
%				\draw[thin] ({4*\a},0.05) -- ({4*\a},-0.05);
%				\draw[->, thick] ({4*\a},0.70) -- ({4*\a},0.10)
%				node[right=3pt, pos=0] {$F$};
%				
%				%% quotatura
%				\draw[richiam] (0,      -\thA-0.15) -- (0,      -1.10);
%				\draw[richiam] ({2*\a}, -\thR-0.40) -- ({2*\a}, -1.10);
%				\draw[richiam] ({3*\a}, -0.10)      -- ({3*\a}, -1.10);
%				\draw[richiam] ({5*\a}, -\thR-0.40) -- ({5*\a}, -1.10);
%				
%				\draw[tacca] (0,     -1.10) -- (0,     -1.30);
%				\draw[tacca] ({2*\a},-1.10) -- ({2*\a},-1.30);
%				\draw[tacca] ({3*\a},-1.10) -- ({3*\a},-1.30);
%				\draw[tacca] ({5*\a},-1.10) -- ({5*\a},-1.30);
%				
%				\draw[<->, thin, gray] (0,     -1.20) -- ({2*\a},-1.20);
%				\node[above=0.04] at ({\a},   -1.20) {$2a$};
%				\draw[<->, thin, gray] ({2*\a},-1.20) -- ({3*\a},-1.20);
%				\node[above=0.04] at ({2.5*\a},-1.20) {$a$};
%				\draw[<->, thin, gray] ({3*\a},-1.20) -- ({5*\a},-1.20);
%				\node[above=0.04] at ({4*\a}, -1.20) {$2a$};
%				
%				\node at ({2.5*\a}, -2.05) {(b)};
%				
%			\end{scope}
%			
%		\end{tikzpicture}
%		\caption{Trave di Gerber (\ESEref{ex:gerber}): dato cinematico
%			(a) e dato statico (b). Cerniera in $A$, carrello (intermedio)
%			in $B$, cerniera interna in $C$, carrello (terminale) in $D$;
%			poli $O_1\equiv A$ e $O_2\equiv C$.}
%		\label{fig:gerber}
%	\end{figure}
%	
%	\begin{svolgimento}
%		
%		\noindent\textbf{Vettore di configurazione e matrice
%			dei vincoli.}
%		Con $n_c = 2$, $n = 6$:
%		\begin{equation*}
%			\bm{u} =
%			\bigl\{u(O_1),\; v(O_1),\; \vartheta_1,\;
%			u(O_2),\; v(O_2),\; \vartheta_2
%			\bigr\}^\top \in \mathbb{R}^6.
%		\end{equation*}
%		I vincoli scalari sono $m = 6$. La matrice dei
%		vincoli e il vettore delle reazioni sono:
%		\begin{equation*}\label{eq:A_gerber}
%			\bm{A} =
%			\begin{bmatrix}
%				1 &  0 &   0 & 0 & 0 &  0 \\
%				0 &  1 &   0 & 0 & 0 &  0 \\
%				0 &  1 &  2a & 0 & 0 &  0 \\
%				-1 &  0 &   0 & 1 & 0 &  0 \\
%				0 & -1 & -3a & 0 & 1 &  0 \\
%				0 &  0 &   0 & 0 & 1 & 2a
%			\end{bmatrix},
%			\qquad
%			\bm{r} =
%			\begin{Bmatrix}
%				X_A \\ Y_A \\ Y_B \\ X_C \\ Y_C \\ Y_D
%			\end{Bmatrix}.
%		\end{equation*}
%		
%		\noindent\textbf{Classificazione.}
%		$\nu = n - m = 0$ e $\rg(\bm{A}) = 6$ per ogni
%		$a \neq 0$: il sistema è \textbf{isostatico}.
%		
%		\smallskip
%		\noindent\textbf{Problema cinematico:
%			$\bm{A}\,\bm{u} = \bm{s}$.}
%		Il vettore dei cedimenti è:
%		\begin{equation*}
%			\bm{s} =
%			\bigl\{0,\; 0,\; -2\bar{v},\;
%			0,\; 0,\; -\bar{v}
%			\bigr\}^\top.
%		\end{equation*}
%		La soluzione si ottiene per sostituzione in avanti:
%		\begin{align*}
%			\text{riga (i):}&\quad u(O_1) = 0,\\
%			\text{riga (ii):}&\quad v(O_1) = 0,\\
%			\text{riga (iii):}&\quad
%			v(O_1) + 2a\,\vartheta_1 = -2\bar{v}
%			\;\Longrightarrow\;
%			\vartheta_1 = -\frac{\bar{v}}{a},\\
%			\text{riga (iv):}&\quad
%			u(O_2) = u(O_1) = 0,\\
%			\text{riga (v):}&\quad
%			v(O_2) = v(O_1) + 3a\,\vartheta_1
%			= -3\bar{v},\\
%			\text{riga (vi):}&\quad
%			v(O_2) + 2a\,\vartheta_2 = -\bar{v}
%			\;\Longrightarrow\;
%			\vartheta_2 = \frac{-\bar{v} + 3\bar{v}}{2a}
%			= \frac{\bar{v}}{a}.
%		\end{align*}
%		La configurazione del sistema è:
%		\begin{equation*}
%			u(O_1) = 0, \quad
%			v(O_1) = 0, \quad
%			\vartheta_1 = -\frac{\bar{v}}{a},
%			\quad
%			u(O_2) = 0, \quad
%			v(O_2) = -3\bar{v}, \quad
%			\vartheta_2 = \frac{\bar{v}}{a}.
%		\end{equation*}
%		$\mathcal{C}_1$ ruota in senso orario attorno ad $A$;
%		$\mathcal{C}_2$ subisce una traslazione verticale di
%		$-3\bar{v}$ e una rotazione antioraria di $\bar{v}/a$.
%		Si noti che il cedimento in $D$ è:
%		\begin{equation*}
%			v(O_2) + 2a\,\vartheta_2 =
%			-3\bar{v} + 2\bar{v} = -\bar{v} \;\checkmark
%		\end{equation*}
%		e quello in $B$:
%		\begin{equation*}
%			v(O_1) + 2a\,\vartheta_1 =
%			0 - 2\bar{v} = -2\bar{v}. \;\checkmark
%		\end{equation*}
%		
%		\smallskip
%		\noindent\textbf{Problema statico:
%			$\bm{A}^{\top}\bm{r} = -\bm{f}$.}
%		Il vettore delle forze generalizzate è:
%		\begin{equation*}
%			\bm{f} =
%			\bigl\{0,\; -2F,\; -2Fa,\;
%			0,\; -F,\; -Fa
%			\bigr\}^\top,
%		\end{equation*}
%		dove $-2F$ e $-2Fa$ sono rispettivamente la
%		risultante verticale e il momento rispetto ad $A$
%		delle forze su $\mathcal{C}_1$; analogamente $-F$ e
%		$-Fa$ per le forze su $\mathcal{C}_2$ rispetto al
%		polo $C$. La matrice del problema statico è:
%		\begin{equation*}\label{eq:B_gerber}
%			\bm{A}^{\top} =
%			\begin{bmatrix}
%				1 &  0 &  0  & -1 &  0  & 0  \\
%				0 &  1 &  1  &  0 & -1  & 0  \\
%				0 &  0 & 2a  &  0 & -3a & 0  \\
%				0 &  0 &  0  &  1 &  0  & 0  \\
%				0 &  0 &  0  &  0 &  1  & 1  \\
%				0 &  0 &  0  &  0 &  0  & 2a
%			\end{bmatrix}.
%		\end{equation*}
%		Il sistema $\bm{A}^{\top}\bm{r} = -\bm{f}$ si
%		risolve per sostituzione all'indietro:
%		\begin{align*}
%			\text{riga (vi):}&\quad
%			2a\,Y_D = Fa
%			\;\Longrightarrow\; Y_D = \frac{F}{2},\\
%			\text{riga (v):}&\quad
%			Y_C + Y_D = F
%			\;\Longrightarrow\; Y_C = \frac{F}{2},\\
%			\text{riga (iv):}&\quad X_C = 0,\\
%			\text{riga (iii):}&\quad
%			2a\,Y_B - 3a\,Y_C = 2Fa
%			\;\Longrightarrow\;
%			Y_B = \frac{2Fa + 3Fa/2}{2a} = \frac{7F}{4},\\
%			\text{riga (ii):}&\quad
%			Y_A + Y_B - Y_C = 2F
%			\;\Longrightarrow\;
%			Y_A = 2F + Y_C - Y_B =
%			2F + \frac{F}{2} - \frac{7F}{4} = \frac{3F}{4},\\
%			\text{riga (i):}&\quad X_A = 0.
%		\end{align*}
%		
%		\noindent\textbf{Dualismo.}
%		La struttura triangolare a blocchi di $\bm{A}^\top$
%		riflette la topologia del sistema: le righe (iv)--(vi)
%		coinvolgono le sole reazioni di $\mathcal{C}_2$ e si
%		risolvono per prime; le righe (i)--(iii) coinvolgono
%		le reazioni di $\mathcal{C}_1$ e la reazione interna
%		$Y_C$ già nota. L'ordine di risoluzione del problema
%		statico --- da $\mathcal{C}_2$ a $\mathcal{C}_1$ ---
%		è l'inverso di quello del problema cinematico, dove
%		$\mathcal{C}_1$ determina $\mathcal{C}_2$.
%		
%		\noindent\textbf{Verifica.}
%		\textit{Corpo $\mathcal{C}_2$}:
%		\begin{equation*}
%			\sum_{(2)} Y = Y_C + Y_D - F =
%			\frac{F}{2} + \frac{F}{2} - F = 0,\;\checkmark
%			\qquad
%			\sum\mu_C = Y_D \cdot 2a - F \cdot a =
%			0.\;\checkmark
%		\end{equation*}
%		\textit{Corpo $\mathcal{C}_1$}:
%		\begin{equation*}
%			\sum_{(1)}  Y= Y_A + Y_B - Y_C - 2F =
%			\frac{3F}{4} + \frac{7F}{4} -
%			\frac{F}{2} - 2F = 0,\;\checkmark
%		\end{equation*}
%		\begin{equation*}
%			\sum_{(1)}\mu_A = Y_B \cdot 2a - Y_C \cdot 3a
%			- 2F \cdot a =
%			\frac{7F}{4} \cdot 2a -
%			\frac{F}{2} \cdot 3a - 2Fa =
%			\frac{7Fa}{2} - \frac{3Fa}{2} - 2Fa =
%			0.\;\checkmark
%		\end{equation*}
%	\end{svolgimento}
%\end{esempio}
\begin{esempio}[Sistema labile: compatibilità statica
	--- trave di Gerber senza carrello in $B$]
	\label{ex:gerber_labile}
	Si consideri il sistema della trave di Gerber di
	\FIGref{fig:gerber_motiv} (già risolto nelle Sezioni
	\ref{sec:congruenza_SCR}-\ref{sec:gerber_equilibrio}),
	dal quale si rimuove il carrello in $B$, mantenendo tutti gli
	altri vincoli. Allora $m = 5$ e $\nu = 6 - 5 = 1$: il sistema è
	labile. Sono applicati gli stessi carichi già visti
	in~\eqref{eq:f_gerber}: una forza $2F$ verso il
	basso nel punto medio di $AB$ e una forza $F$ verso
	il basso nel punto medio di $\mathcal{C}_2$.
	Determinare il meccanismo e verificare la
	compatibilità statica.
	
	\begin{svolgimento}
		
		\noindent\textbf{Meccanismo.}
		La matrice $\bm{A}$ è la sottomatrice delle righe
		(i), (ii), (iv), (v), (vi) di~\eqref{eq:A_gerber_SCR},
		ottenuta sopprimendo la riga (iii) corrispondente al
		carrello in $B$. Risolvendo $\bm{A}\,\bm{u} = \bm{0}$:
		\begin{align*}
			\text{riga (i):}&\quad u(O_1) = 0,\\
			\text{riga (ii):}&\quad v(O_1) = 0,\\
			\text{riga (iv):}&\quad u(O_2) = 0,\\
			\text{riga (v):}&\quad v(O_2) = 3a\,q,\\
			\text{riga (vi):}&\quad
			\vartheta_2 = -\frac{3a\,q}{2a}
			= -\frac{3}{2}\,q,
		\end{align*}
		con $q \in \mathbb{R}$ coordinata lagrangiana libera.
		La soluzione generale è $\bm{u} = q\,\bm{U}_q$,
		dove, normalizzando a $q = 1$, la base del meccanismo
		è:
		\begin{equation*}
			\bm{U}_q =
			\bigl\{0,\; 0,\; 1,\; 0,\; 3a,\;
			-\tfrac{3}{2}
			\bigr\}^\top.
		\end{equation*}
		Il sistema è \textbf{labile di grado} $g_\ell = 1$.
		
		\noindent\textbf{Verifica del meccanismo.}
		Lo spostamento verticale di $D$ nel meccanismo è:
		\begin{equation*}
			v(D) = v(O_2) + 2a\,\vartheta_2
			= 3a\,q - 2a\cdot\frac{3}{2}\,q
			= 3aq - 3aq = 0. \;\checkmark
		\end{equation*}
		Il punto $D$ rimane fisso, coerentemente con il
		fatto che il carrello in $D$ è un vincolo attivo
		che impedisce lo spostamento verticale di $D$.
		
		Fisicamente: $\mathcal{C}_1$ ruota attorno ad $A$
		di ampiezza $q$; la cerniera interna in $C$
		costringe $\mathcal{C}_2$ a seguire, con $C$ che si
		alza di $3a\,q$ e $D$ che rimane fisso. Il meccanismo
		è una rotazione coordinata dei due corpi, vincolata
		dal carrello in $D$.
		
		
		
		\noindent\textbf{Compatibilità statica.}
		Il problema statico $\bm{A}^\top\bm{r} = -\bm{f}$
		ammette soluzione se e solo se
		$\bm{U}_q^\top\bm{f} = 0$. Il vettore delle forze
		generalizzate è ancora quello di~\eqref{eq:f_gerber}:
		\begin{equation*}
			\bm{f} =
			\bigl\{0,\; -2F,\; -2Fa,\;
			0,\; -F,\; -Fa
			\bigr\}^\top.
		\end{equation*}
		La condizione di compatibilità dà:
		\begin{equation*}
			\bm{U}_q^\top\bm{f} =
			1\cdot(-2Fa) + 3a\cdot(-F) +
			\left(-\frac{3}{2}\right)\cdot(-Fa)
			= -2Fa - 3Fa + \frac{3Fa}{2}
			= -\frac{7Fa}{2} \neq 0.
		\end{equation*}
		La condizione non è soddisfatta: il sistema labile
		non è in equilibrio sotto questo carico. Fisicamente,
		senza il carrello in $B$, non esiste alcuna
		combinazione di reazioni vincolari — nelle sole
		cerniere in $A$, $C$ e nel carrello in $D$ ---
		capace di equilibrare il momento risultante delle
		forze attive nel meccanismo: la reazione in $B$
		era l'unica in grado di fornire il contributo
		mancante.
		
		\vspace{6pt}
		\noindent\textbf{Carico compatibile.}
		La condizione $\bm{U}_q^\top\bm{f} = 0$ può
		essere soddisfatta per opportune scelte del
		carico. Ad esempio, applicando su $\mathcal{C}_1$
		una coppia $\mu$ in aggiunta alle forze assegnate,
		la condizione diventa:
		\begin{equation*}
			-2Fa + \mu - 3Fa + \frac{3Fa}{2} = 0
			\;\Longrightarrow\;
			\mu = \frac{7Fa}{2}.
		\end{equation*}
		Solo con questa coppia aggiuntiva il sistema
		labile è in equilibrio.
	\end{svolgimento}
\end{esempio}


\tR{Da verificare}
\begin{esempio}[Trave su tre appoggi: autotensione e
	compatibilità cinematica]
	\label{ex:trave_tre_appoggi}
	Si consideri una trave rigida vincolata da una cerniera in
	$A=(0,0)$ e da due carrelli verticali, in $B=(a,0)$ e in
	$C=(2a,0)$. Si scelga il polo $O\equiv A$.
	
	\begin{figure}[ht]
		\centering
		\begin{tikzpicture}[>=Stealth, line width=0.8pt, font=\small,
			nodo/.style={circle, fill=white, draw=black,
				inner sep=0pt, minimum size=5pt, line width=0.8pt},
			tacca/.style={thin, gray!70}]
			
			\def\a{2.0}     % unità di lunghezza (= a)
			\def\tw{0.22}
			\def\thA{0.38}
			\def\thR{0.32}
			\def\bw{0.16}
			\def\gw{0.30}
			\def\sr{0.08}
			
			%% trave
			\draw[line width=2.2pt] (0,0) -- ({2*\a},0);
			
			%% etichette nodi
			\node[above=2pt] at (0,0)      {$A$};
			\node[above=2pt] at (\a,0)     {$B$};
			\node[above=2pt] at ({2*\a},0) {$C$};
			
			%% cerniera in A
			\draw (0,0) -- ++(-\tw,-\thA) -- ++(2*\tw,0) -- cycle;
			\draw (-\gw,-\thA) -- (\gw,-\thA);
			\foreach \x in {-0.24,-0.12,0,0.12,0.24}
			\draw[thin] (\x,-\thA) -- ++(-0.09,-0.11);
			\node[nodo] at (0,0) {};
			
			%% carrello intermedio in B
			\draw (\a,0) -- ++(-\tw,-\thR) -- ++(2*\tw,0) -- cycle;
			\draw (\a-\gw,-\thR) -- (\a+\gw,-\thR);
			\draw (\a-\bw,-\thR-0.12) circle (2.4pt);
			\draw (\a+\bw,-\thR-0.12) circle (2.4pt);
			\draw (\a-\gw,-\thR-0.24) -- (\a+\gw,-\thR-0.24);
			\foreach \x in {-0.24,-0.12,0,0.12,0.24}
			\draw[thin] (\a+\x,-\thR-0.24) -- ++(-0.09,-0.11);
			\fill[white] (\a-\sr,0) arc(180:360:\sr) -- cycle;
			\draw (\a-\sr,0) arc(180:360:\sr);
			
			%% carrello terminale in C
			\draw ({2*\a},0) -- ++(-\tw,-\thR) -- ++(2*\tw,0) -- cycle;
			\draw ({2*\a-\gw},-\thR) -- ({2*\a+\gw},-\thR);
			\draw ({2*\a-\bw},-\thR-0.12) circle (2.4pt);
			\draw ({2*\a+\bw},-\thR-0.12) circle (2.4pt);
			\draw ({2*\a-\gw},-\thR-0.24) -- ({2*\a+\gw},-\thR-0.24);
			\foreach \x in {-0.24,-0.12,0,0.12,0.24}
			\draw[thin] ({2*\a+\x},-\thR-0.24) -- ++(-0.09,-0.11);
			\node[nodo] at ({2*\a},0) {};
			
			%% quotatura
			\draw[tacca] (0,-0.90) -- (0,-1.10);
			\draw[tacca] (\a,-0.90) -- (\a,-1.10);
			\draw[tacca] ({2*\a},-0.90) -- ({2*\a},-1.10);
			
			\draw[<->, thin, gray] (0,-1.00) -- (\a,-1.00);
			\node[above=0.04] at ({0.5*\a},-1.00) {$a$};
			\draw[<->, thin, gray] (\a,-1.00) -- ({2*\a},-1.00);
			\node[above=0.04] at ({1.5*\a},-1.00) {$a$};
			
		\end{tikzpicture}
		\caption{Trave su tre appoggi (\ESEref{ex:trave_tre_appoggi}):
			cerniera in $A=(0,0)$, carrello verticale intermedio in
			$B=(a,0)$, carrello verticale terminale in $C=(2a,0)$;
			polo $O\equiv A$. Il vincolo ridondante --- uno dei tre
			è in eccesso rispetto al minimo necessario per un unico
			corpo rigido --- rende il sistema iperstatico di grado 1.}
		\label{fig:trave_tre_appoggi}
	\end{figure}
	
	\begin{svolgimento}
		
		\noindent\textbf{Vettore di configurazione e matrice
			dei vincoli.} Con $n_c=1$, $n=3$:
		$\bm u=\{u(O),\,v(O),\,\vartheta\}^\top$. I vincoli
		scalari sono $m=4$: le due componenti della cerniera in
		$A$ e le componenti verticali dei due carrelli. La
		matrice dei vincoli e il vettore delle reazioni sono:
		\begin{equation*}
			\bm A =
			\begin{bmatrix}
				1 & 0 & 0\\
				0 & 1 & 0\\
				0 & 1 & a\\
				0 & 1 & 2a
			\end{bmatrix},
			\qquad
			\bm r =
			\begin{Bmatrix}
				X_A\\ Y_A\\ Y_B\\ Y_C
			\end{Bmatrix}.
		\end{equation*}
		
		\noindent\textbf{Classificazione.}
		$\nu=n-m=3-4=-1$. Le prime tre righe di $\bm A$ hanno
		determinante $a\neq0$, quindi $\rg(\bm A)=3=n$: il
		sistema è \textbf{iperstatico di grado} $g_i=|\nu|=1$
		(e $g_\ell=0$: nessun meccanismo).
		
		\noindent\textbf{Autotensione.}
		Sia $\bm R_\chi\in\mathbb R^{4\times 1}$ una base di
		$\ker\bm A^\top$. Risolvendo $\bm A^\top\bm r_n=\bm 0$:
		\begin{align*}
			\text{(colonna $u(O)$):}&\quad X_A=0,\\
			\text{(colonna $\vartheta$):}&\quad
			a\,Y_B+2a\,Y_C=0
			\;\Longrightarrow\; Y_B=-2Y_C,\\
			\text{(colonna $v(O)$):}&\quad
			Y_A+Y_B+Y_C=0
			\;\Longrightarrow\; Y_A=Y_C.
		\end{align*}
		Ponendo $Y_C=\chi$ e normalizzando a $\chi=1$:
		\begin{equation*}
			\bm R_\chi = \{0,\;1,\;-2,\;1\}^\top.
		\end{equation*}
		
		\noindent\textbf{Verifica dell'autotensione.}
		Risultante verticale: $Y_A+Y_B+Y_C=1-2+1=0$.
		\;$\checkmark$\quad Momento rispetto ad $A$:
		$Y_B\cdot a+Y_C\cdot 2a=-2a+2a=0$. \;$\checkmark$
		È dunque un sistema di reazioni in equilibrio pur in
		assenza di ogni carico esterno --- un'autotensione:
		possibile proprio perché il vincolo in $C$ è ridondante
		rispetto ai primi tre.
		
		\noindent\textbf{Compatibilità cinematica.}
		Il problema cinematico $\bm A\bm u=\bm s$ ammette
		soluzione se e solo se $\bm R_\chi^\top\bm s=0$.
		Imponiamo cedimenti $\bar s_B$ in $B$ e $\bar s_C$ in
		$C$ (cerniera in $A$ perfetta):
		$\bm s=\{0,\,0,\,\bar s_B,\,\bar s_C\}^\top$. La
		condizione di compatibilità dà:
		\begin{equation*}
			\bm R_\chi^\top\bm s =
			-2\,\bar s_B+\bar s_C = 0
			\quad\Longrightarrow\quad
			\bar s_C = 2\,\bar s_B.
		\end{equation*}
		Il risultato si verifica per ispezione diretta: con la
		cerniera perfetta, $v(O)=0$, e ogni configurazione
		rigida ha $v(x)=x\vartheta$; dunque
		$v(a)=a\vartheta=\bar s_B$ e $v(2a)=2a\vartheta=2\bar
		s_B$, che deve coincidere con $\bar s_C$ --- esattamente
		la condizione trovata.
		
		\vspace{6pt}
		\noindent\textbf{Cedimento incompatibile.}
		Se i due appoggi cedono \emph{indipendentemente} ---
		per esempio $\bar s_B=\delta$ e $\bar s_C=\delta$
		(stesso cedimento in entrambi) --- la condizione non è
		soddisfatta ($\delta\neq 2\delta$ per $\delta\neq0$):
		nessuna configurazione rigida della trave è compatibile
		con questi due dati simultaneamente. Fisicamente, è
		proprio l'iperstaticità --- il vincolo ridondante in
		$C$ --- a rendere il sistema sensibile a cedimenti
		differenziali: una trave isostatica (con un solo
		carrello) si adatterebbe a qualunque cedimento
		imposto, ruotando semplicemente attorno agli altri due
		vincoli; qui, il terzo vincolo lo impedisce, a meno che
		i due cedimenti non siano già nel rapporto giusto.
		
	\end{svolgimento}
\end{esempio}

\begin{esempio}[Trave semplicemente appoggiata: TSV e TFV]
	\label{ex:trave_semplice_virt}
	Si riprende la trave semplicemente appoggiata
	dell'\ESEref{ex:trave_semplice}: cerniera in $A$,
	carrello in $B=(\ell,0)$, sistema isostatico con
	cedimento $\bar v_B$ al carrello $B$ e carico $F$
	applicato in $x=a$. Calcolare la reazione $Y_B$ tramite
	il teorema degli spostamenti virtuali (Müller-Breslau) e
	lo spostamento $v(a)$ tramite il teorema delle forze
	virtuali.
	
	\begin{svolgimento}
		
		\noindent\textbf{Teorema degli spostamenti virtuali.}
		Si assegna un cedimento unitario al vincolo 3 (il
		carrello in $B$): $\bm s^*=\bm e_3=\{0\;\,0\;\,1\}^\top$.
		Il problema virtuale cinematico $\bm A\bm u^*=\bm s^*$ dà:
		\begin{equation*}
			u^*=0,\qquad v^*=0,\qquad \vartheta^*=\frac{1}{\ell}.
		\end{equation*}
		Con il carico reale $\bm f=\{0\;\,{-F}\;\,{-Fa}\}^\top$:
		\begin{equation*}
			\bm f^\top\bm u^* =
			0\cdot 0+(-F)\cdot 0+(-Fa)\cdot\frac{1}{\ell}
			= -\frac{Fa}{\ell}.
		\end{equation*}
		Per~\eqref{eq:mueller_breslau}:
		\begin{equation*}
			Y_B=-\bm f^\top\bm u^*=\frac{Fa}{\ell},
		\end{equation*}
		in accordo con il valore già trovato
		nell'\ESEref{ex:trave_semplice}.
		
		\noindent\textbf{Teorema delle forze virtuali.}
		Si vuole calcolare lo spostamento verticale $v(a)$ del
		punto di applicazione del carico $F$. Per la formula del
		moto rigido, $v(a)=v+a\vartheta$: è un funzionale lineare
		di $\bm u$ con coefficienti $\bm c=\{0\;\,1\;\,a\}^\top$.
		Per la generalizzazione a funzionali lineari vista sopra,
		si risolve il problema virtuale statico
		$\bm A^\top\bm r^*=-\bm c$:
		\begin{equation*}
			X_A^*=0,\qquad
			\ell\,Y_B^*=-a \;\Longrightarrow\; Y_B^*=-\frac{a}{\ell},
			\qquad
			Y_A^*+Y_B^*=-1 \;\Longrightarrow\;
			Y_A^*=\frac{a}{\ell}-1.
		\end{equation*}
		Con il cedimento reale $\bm s=\{0\;\,0\;\,{-\bar v_B}\}^\top$:
		\begin{equation*}
			v(a) = -\bm s^\top\bm r^* =
			-\Bigl[(-\bar v_B)\cdot\Bigl(-\frac{a}{\ell}\Bigr)\Bigr]
			= -\frac{a}{\ell}\,\bar v_B.
		\end{equation*}
		Il segno è coerente con il calcolo diretto già svolto
		nell'\ESEref{ex:trave_semplice}: se $B$ scende
		($\bar v_B>0$), l'intera trave ruota in senso orario
		attorno ad $A$, e ogni punto fra $A$ e $B$ --- in
		particolare quello in $x=a$ --- scende con essa.
		
		\begin{oss}
			I due calcoli mostrano l'uso speculare di TSV e TFV:
			per $Y_B$ si risolve un problema cinematico virtuale
			con cedimento unitario al vincolo cercato; per $v(a)$
			si risolve un problema statico virtuale con un carico
			generalizzato pari ai coefficienti del funzionale
			cercato.
		\end{oss}
		
	\end{svolgimento}
\end{esempio}


\begin{esempio}[Trave di Gerber: teorema degli spostamenti virtuali]
	\label{ex:gerber_TSV}
	Si riprende la trave di Gerber di \FIGref{fig:gerber_motiv},
	con il carico già visto in~\eqref{eq:f_gerber} (una forza
	$2F$ sul punto medio di $AB$, una forza $F$ sul punto medio
	di $\mathcal{C}_2$). Calcolare la reazione $Y_B$ tramite il
	teorema degli spostamenti virtuali, senza risolvere
	l'intero sistema $\bm A^\top\bm r=-\bm f$.
	
	\begin{svolgimento}
		
		Si assegna un cedimento unitario al vincolo 3 (il
		carrello in $B$): $\bm s^*=\bm e_3$. Il problema
		virtuale $\bm A\bm u^*=\bm e_3$, con la stessa matrice
		$\bm A$ di~\eqref{eq:A_gerber_SCR}, si risolve
		sequenzialmente:
		\begin{align*}
			&\text{righe (i)--(ii):}& u_1^*&=v_1^*=0;\\
			&\text{riga (iii):}& \vartheta_1^*&=\frac{1}{2a};\\
			&\text{riga (iv):}& u_2^*&=0;\\
			&\text{riga (v):}& v_2^*&=3a\,\vartheta_1^*=\frac{3}{2};\\
			&\text{riga (vi):}& \vartheta_2^*&=
			-\frac{v_2^*}{2a}=-\frac{3}{4a}.
		\end{align*}
		Con il carico reale~\eqref{eq:f_gerber}:
		\begin{equation*}
			\bm f^\top\bm u^* =
			(-2Fa)\cdot\frac{1}{2a} + (-F)\cdot\frac{3}{2}
			+ (-Fa)\cdot\left(-\frac{3}{4a}\right)
			= -F - \frac{3F}{2} + \frac{3F}{4}
			= -\frac{7F}{4}.
		\end{equation*}
		Per~\eqref{eq:mueller_breslau}:
		\begin{equation*}
			Y_B = -\bm f^\top\bm u^* = \frac{7F}{4},
		\end{equation*}
		esattamente il valore già trovato risolvendo il sistema
		completo in~\eqref{eq:reazioni_gerber} --- qui ottenuto
		senza calcolare le altre cinque reazioni.
		
	\end{svolgimento}
\end{esempio}

%-------------------------------------------------------------------------------
\section{Esercizi proposti}
%-------------------------------------------------------------------------------

\paragraph*{Convenzioni di nomenclatura.}
\textbf{Reazioni.} Le componenti del vettore $\bm{r}$ si
denotano con simboli mnemonici legati al punto di applicazione:
$H_X, V_X$ per le componenti orizzontale e verticale di una
cerniera in $X$; $M_X$ per il momento di un incastro in $X$;
$R_X$ per il modulo della reazione di un carrello in $X$ lungo
il proprio asse, sostituito da $V_X$ (o $H_X$) se l'asse del
carrello è verticale (od orizzontale).
\textbf{Cedimenti.} Le componenti del vettore $\bm{s}$ si
scrivono $\bar{u}_X, \bar{v}_X$ per i cedimenti traslazionali e
$\bar{\theta}$ per la rotazione di un incastro; $\bar{s}_X$ per
il cedimento di un vincolo orientato (carrello inclinato), con
significato di proiezione dello spostamento sul proprio asse.



\begin{esercizio}[Trave incastrata: mensola]
	\label{ex:C05_mensola}
	
	Una trave rigida $AB$ di lunghezza $\ell$ è incastrata in $B$,
	con $A=(0,0)$ e $B=(\ell,0)$. Si scelga il polo $O\equiv A$.
	(a) Scrivere la matrice $\bm{A}$ del problema cinematico.
	(b) Verificare che il sistema è isostatico.
	(c) Risolvere il problema cinematico assumendo come unico
	cedimento la rotazione $\bar{\theta}$ dell'incastro;
	identificare il centro di rotazione assoluto del corpo.
	(d) Risolvere il problema statico per una forza $F$ verticale
	verso il basso applicata in $A$, e verificare le equazioni
	cardinali di equilibrio.
	
	\risultato (a) Le tre equazioni di vincolo dell'incastro in
	$B$, espresse rispetto al polo $A$, sono $u_A=\bar{u}_B$,
	$v_A+\theta\ell=\bar{v}_B$, $\theta=\bar{\theta}$, dove la
	seconda è la formula del moto rigido $v(B)=v_A+\theta\,(x_B-x_A)$
	scritta come equazione di vincolo. Da qui
	\[
	\bm{A}=
	\begin{pmatrix}
		1 & 0 & 0\\
		0 & 1 & \ell\\
		0 & 0 & 1
	\end{pmatrix}.
	\]
	(b) $\det(\bm{A})=1\neq 0$, quindi $\rg(\bm{A})=3$ e
	$g_\ell=g_{\mathscr{i}}=0$: il sistema è isostatico.
	(c) Da $\bm{s}=(0\;0\;\bar{\theta})^\top$ si ottiene
	\[
	u_A=0,\qquad v_A=-\bar{\theta}\,\ell,\qquad \theta=\bar{\theta}:
	\]
	il polo $A$ trasla verticalmente di $-\bar{\theta}\,\ell$ e
	il corpo ruota di $\bar{\theta}$. Il CR assoluto coincide
	con $B$: dalla formula del moto rigido
	$u(B)=u_A=0$ e
	$v(B)=v_A+\theta\,\ell=-\bar{\theta}\,\ell+\bar{\theta}\,\ell=0$.
	(d) $\bm{f}=(0\;-F\;0)^\top$. Risolvendo
	$\bm{A}^\top\bm{r}=-\bm{f}$ si ottengono le tre componenti
	dell'incastro $H_B=0$, $V_B=F$, $M_B=-F\ell$. Risultante
	verticale: $V_B-F=0$. Momento rispetto ad $A$:
	$M_B+V_B\,\ell=-F\ell+F\ell=0$. Equilibrio verificato.
\end{esercizio}

\begin{figure}[ht]
	\centering
	\begin{tikzpicture}[>=Stealth, line width=0.8pt, font=\small,
		nodo/.style={circle, fill=white, draw=black,
			inner sep=0pt, minimum size=5pt, line width=0.8pt},
		tacca/.style={thin, gray!70}]
		
		\def\L{5.0}
		\def\tw{0.22}\def\thA{0.38}\def\thR{0.32}
		\def\bw{0.16}\def\gw{0.30}
		
		%% trave
		\draw[line width=2.2pt] (0,0) -- (\L,0);
		\node[above=2pt] at (0,0)  {$A$};
		\node[above=2pt] at (\L,0) {$B$};
		\node[above=2pt] at (\L/2,0) {$C$};
		
		%% cerniera in A
		\draw (0,0) -- ++(-\tw,-\thA) -- ++(2*\tw,0) -- cycle;
		\draw (-\gw,-\thA) -- (\gw,-\thA);
		\foreach \x in {-0.24,-0.12,0,0.12,0.24}
		\draw[thin] (\x,-\thA) -- ++(-0.09,-0.11);
		\node[nodo] at (0,0) {};
		
		%% carrello terminale in B, ruotato di -45° (asse a 45° dalla verticale)
		\begin{scope}[rotate around={-45:(\L,0)}]
			\draw (\L,0) -- ++(-\tw,-\thR) -- ++(2*\tw,0) -- cycle;
			\draw (\L-\gw,-\thR) -- (\L+\gw,-\thR);
			\draw (\L-\bw,-\thR-0.12) circle (2.4pt);
			\draw (\L+\bw,-\thR-0.12) circle (2.4pt);
			\draw (\L-\gw,-\thR-0.24) -- (\L+\gw,-\thR-0.24);
			\foreach \x in {-0.24,-0.12,0,0.12,0.24}
			\draw[thin] (\L+\x,-\thR-0.24) -- ++(-0.09,-0.11);
		\end{scope}
		\node[nodo] at (\L,0) {};
		
		%% asse del carrello (versore a)
		\draw[->, thick, gray!50] (\L,0) -- ++(0.50, 0.50)
		node[right=2pt, black] {$\ab$};
		
		%% forza F a 45° in C (verso il basso a destra)
		\draw[thin] (\L/2,0.05) -- (\L/2,-0.05);
		\draw[->, thick] (\L/2-0.55, 0.55) -- (\L/2-0.06, 0.06)
		node[above left=2pt, pos=0] {$F$};
		
		%% quota
		\draw[tacca] (0, -1.05) -- (0, -0.85);
		\draw[tacca] (\L,-1.05) -- (\L,-0.85);
		\draw[<->, thin] (0,-0.95) -- (\L,-0.95)
		node[midway, above=2pt] {$\ell$};
		
	\end{tikzpicture}
	\caption{Trave appoggiata con carrello inclinato
		(\ESEref{ex:C05_trave_inclinata_stat} e
		\ESEref{ex:C05_trave_inclinata_cin}): cerniera in $A$,
		carrello in $B$ con asse $\ab$ a $45^\circ$ rispetto alla
		verticale, polo $O\equiv A$. Nel dato statico è applicata
		in $C=(\ell/2,0)$ una forza $F$ inclinata di $45^\circ$
		sull'orizzontale.}
	\label{fig:C05_trave_inclinata}
\end{figure}

\begin{esercizio}[Trave appoggiata con carrello inclinato — dato statico]
	\label{ex:C05_trave_inclinata_stat}
	
	Si consideri la trave di \FIGref{fig:C05_trave_inclinata}:
	cerniera in $A=(0,0)$, carrello in $B=(\ell,0)$ di asse
	$\ab=\tfrac{\sqrt{2}}{2}(\ab_x+\ab_y)$, polo $O\equiv A$.
	È applicata in $C=(\ell/2,0)$ una forza inclinata di
	$45^\circ$ sull'orizzontale, di componenti
	$\fb=\tfrac{\sqrt{2}}{2}F\,\ab_x-\tfrac{\sqrt{2}}{2}F\,\ab_y$.
	(a) Scrivere $\bm{A}$. (b) Classificare il sistema.
	(c) Calcolare le reazioni vincolari risolvendo
	$\bm{A}^\top\bm{r}=-\bm{f}$.
	(d) Verificare le equazioni cardinali di equilibrio.
	
	\risultato (a) Le due equazioni della cerniera in $A$ sono
	$u_A=\bar{u}_A$ e $v_A=\bar{v}_A$. Per il carrello, la
	condizione $\ub(B)\cdot\ab=\bar{s}_B$ si esplicita applicando
	la formula del moto rigido $\ub(B)=\ub(A)+\bm{\upvartheta}\times
	\overrightarrow{AB}$ e calcolando il prodotto misto come
	in~\eqref{eq:vincolo_est_scal}, con $\overrightarrow{AB}=\ell\,\ab_x$:
	\[
	\tfrac{\sqrt{2}}{2}\,\bigl(u_A+v_A+\theta\,\ell\bigr)
	=\bar{s}_B.
	\]
	Quindi
	\[
	\bm{A}=
	\begin{pmatrix}
		1 & 0 & 0\\
		0 & 1 & 0\\
		\sqrt{2}/2 & \sqrt{2}/2 & \sqrt{2}\,\ell/2
	\end{pmatrix}.
	\]
	(b) $\det(\bm{A})=\sqrt{2}\,\ell/2\neq 0$, $\rg(\bm{A})=3$:
	sistema isostatico ($g_\ell=g_{\mathscr{i}}=0$).
	(c) Le forze generalizzate sono
	\[
	\bm{f}=\bigl(\tfrac{\sqrt{2}}{2}F\;\;
	-\tfrac{\sqrt{2}}{2}F\;\;
	-\tfrac{\sqrt{2}}{4}F\ell\bigr)^\top.
	\]
	Risolvendo $\bm{A}^\top\bm{r}=-\bm{f}$ a partire dalla terza
	equazione si ottengono le reazioni: $R_B=F/2$ (reazione del
	carrello, lungo $\ab$), $H_A=-3\sqrt{2}\,F/4$,
	$V_A=\sqrt{2}\,F/4$ (componenti orizzontale e verticale della
	cerniera).
	(d) Risultante orizzontale:
	\[
	H_A+R_B\,a_x+f_x
	=-\tfrac{3\sqrt{2}}{4}F+\tfrac{F}{2}\cdot\tfrac{\sqrt{2}}{2}
	+\tfrac{\sqrt{2}}{2}F=0.
	\]
	Risultante verticale:
	\[
	V_A+R_B\,a_y+f_y
	=\tfrac{\sqrt{2}}{4}F+\tfrac{F}{2}\cdot\tfrac{\sqrt{2}}{2}
	-\tfrac{\sqrt{2}}{2}F=0.
	\]
	Momento rispetto ad $A$:
	\[
	R_B\,(x_B\,a_y - y_B\,a_x) + \mu_f(A)
	=\tfrac{F}{2}\cdot\tfrac{\sqrt{2}\,\ell}{2}
	-\tfrac{\sqrt{2}}{4}F\ell=0.
	\]
	Equilibrio verificato.
\end{esercizio}

\begin{esercizio}[Trave appoggiata con carrello inclinato — dato cinematico]
	\label{ex:C05_trave_inclinata_cin}
	
	Si consideri lo stesso sistema dell'\ESEref{ex:C05_trave_inclinata_stat}
	in assenza di forze attive ($\bm{f}=\bm{0}$), con un cedimento
	$\bar{s}_B$ \emph{lungo l'asse del carrello}.
	(a) Risolvere $\bm{A}\bm{u}=\bm{s}$ per ottenere il vettore di
	configurazione $\bm{u}$. (b) Calcolare lo spostamento $\ub(B)$
	tramite la formula del moto rigido e mostrare che la sua
	componente verticale non coincide con $\bar{s}_B$.
	(c) Identificare il centro di rotazione assoluto.
	(d) Verificare il bilancio dei lavori virtuali
	$\bm{u}^\top\bm{f}+\bm{s}^\top\bm{r}=0$ con il sistema
	$\bm{f}$, $\bm{r}$ dell'\ESEref{ex:C05_trave_inclinata_stat}.
	
	\risultato (a) Il vettore di cedimenti
	$\bm{s}=(\bar{u}_A\;\bar{v}_A\;\bar{s}_B)^\top$ ha qui
	$\bar{u}_A=\bar{v}_A=0$. Da $\bm{A}\bm{u}=\bm{s}$ con la forma di
	$\bm{A}$ dell'\ESEref{ex:C05_trave_inclinata_stat} si ottiene
	\[
	u_A=0,\qquad v_A=0,\qquad
	\theta=\frac{\sqrt{2}\,\bar{s}_B}{\ell}.
	\]
	(b) Applicando la formula del moto rigido
	$\ub(B)=\ub(A)+\bm{\upvartheta}\times\overrightarrow{AB}$, con
	$\ub(A)=\bm{0}$ e $\overrightarrow{AB}=\ell\,\ab_x$:
	\[
	\ub(B)=\theta\ell\,\ab_y=\sqrt{2}\,\bar{s}_B\,\ab_y.
	\]
	Il punto $B$ si sposta dunque solo verticalmente, di modulo
	$\sqrt{2}\,\bar{s}_B$. Si verifica così la condizione di
	vincolo $\ub(B)\cdot\ab=\bar{s}_B$ già scritta
	nell'\ESEref{ex:C05_trave_inclinata_stat}:
	$\sqrt{2}\,\bar{s}_B\cdot\tfrac{\sqrt{2}}{2}=\bar{s}_B$. Il dato
	$\bar{s}_B$ è dunque la \emph{proiezione} di $\ub(B)$ sull'asse
	del carrello, non il suo modulo $\sqrt{2}\,\bar{s}_B$ né la sua
	componente verticale (che qui coincide col modulo).
	(c) Poiché $\ub(A)=\bm{0}$, il polo $A$ è fisso e coincide con
	il CR assoluto del corpo.
	(d) Con $\bm{f}$ e $\bm{r}$ dell'esercizio precedente, e
	ricordando che la terza componente di $\bm{f}$ è
	$\mu_f(O)=-\tfrac{\sqrt{2}}{4}F\ell$ con $O\equiv A$:
	\[
	\bm{u}^\top\bm{f}
	= \theta\cdot\mu_f(O)
	= \frac{\sqrt{2}\,\bar{s}_B}{\ell}\cdot
	\Bigl(-\tfrac{\sqrt{2}}{4}F\ell\Bigr)
	= -\frac{\bar{s}_B\,F}{2},
	\]
	\[
	\bm{s}^\top\bm{r} = \bar{s}_B\cdot R_B = \frac{\bar{s}_B\,F}{2}.
	\]
	La somma è zero, come previsto da~\eqref{eq:reciprocita_SCR}.
\end{esercizio}

\begin{figure}[ht]
	\centering
	\begin{tikzpicture}[>=Stealth, line width=0.8pt, font=\small,
		nodo/.style={circle, fill=white, draw=black,
			inner sep=0pt, minimum size=5pt, line width=0.8pt},
		tacca/.style={thin, gray!70}]
		
		\def\L{6.0}
		\def\tw{0.22}\def\thR{0.32}
		\def\bw{0.16}\def\gw{0.30}\def\sr{0.08}
		
		%% trave
		\draw[line width=2.2pt] (0,0) -- (\L,0);
		\node[above=2pt] at (0,0)    {$A$};
		\node[above=2pt] at (\L/2,0) {$C$};
		\node[above=2pt] at (\L,0)   {$B$};
		
		%% carrello terminale in A (perno bianco)
		\draw (0,0) -- ++(-\tw,-\thR) -- ++(2*\tw,0) -- cycle;
		\draw (-\gw,-\thR) -- (\gw,-\thR);
		\draw (-\bw,-\thR-0.12) circle (2.4pt);
		\draw (\bw,-\thR-0.12) circle (2.4pt);
		\draw (-\gw,-\thR-0.24) -- (\gw,-\thR-0.24);
		\foreach \x in {-0.24,-0.12,0,0.12,0.24}
		\draw[thin] (\x,-\thR-0.24) -- ++(-0.09,-0.11);
		\node[nodo] at (0,0) {};
		
		%% carrello intermedio in C (mezzaluna)
		\draw (\L/2,0) -- ++(-\tw,-\thR) -- ++(2*\tw,0) -- cycle;
		\draw (\L/2-\gw,-\thR) -- (\L/2+\gw,-\thR);
		\draw (\L/2-\bw,-\thR-0.12) circle (2.4pt);
		\draw (\L/2+\bw,-\thR-0.12) circle (2.4pt);
		\draw (\L/2-\gw,-\thR-0.24) -- (\L/2+\gw,-\thR-0.24);
		\foreach \x in {-0.24,-0.12,0,0.12,0.24}
		\draw[thin] (\L/2+\x,-\thR-0.24) -- ++(-0.09,-0.11);
		\fill[white] (\L/2-\sr,0) arc(180:360:\sr) -- cycle;
		\draw (\L/2-\sr,0) arc(180:360:\sr);
		
		%% carrello terminale in B (perno bianco)
		\draw (\L,0) -- ++(-\tw,-\thR) -- ++(2*\tw,0) -- cycle;
		\draw (\L-\gw,-\thR) -- (\L+\gw,-\thR);
		\draw (\L-\bw,-\thR-0.12) circle (2.4pt);
		\draw (\L+\bw,-\thR-0.12) circle (2.4pt);
		\draw (\L-\gw,-\thR-0.24) -- (\L+\gw,-\thR-0.24);
		\foreach \x in {-0.24,-0.12,0,0.12,0.24}
		\draw[thin] (\L+\x,-\thR-0.24) -- ++(-0.09,-0.11);
		\node[nodo] at (\L,0) {};
		
		%% quotatura
		\draw[tacca] (0,    -1.05) -- (0,    -0.85);
		\draw[tacca] (\L/2, -1.05) -- (\L/2, -0.85);
		\draw[tacca] (\L,   -1.05) -- (\L,   -0.85);
		\draw[<->, thin] (0,-0.95)    -- (\L/2,-0.95)
		node[midway, above=2pt] {$\ell$};
		\draw[<->, thin] (\L/2,-0.95) -- (\L,-0.95)
		node[midway, above=2pt] {$\ell$};
		
	\end{tikzpicture}
	\caption{Trave singola su tre carrelli verticali paralleli
		(\ESEref{ex:C05_tre_carrelli}): un caso di degenerazione
		geometrica in cui $\nu=0$ ma $g_\ell=g_{\mathscr{i}}=1$.}
	\label{fig:C05_tre_carrelli}
\end{figure}

\begin{esercizio}[Tre carrelli paralleli — sistema labile-iperstatico]
	\label{ex:C05_tre_carrelli}
	
	Una trave rigida singola di lunghezza $2\ell$ è vincolata da
	tre carrelli verticali in $A=(0,0)$, $C=(\ell,0)$,
	$B=(2\ell,0)$ (\FIGref{fig:C05_tre_carrelli}). Polo $O\equiv A$.
	(a) Scrivere $\bm{A}$ e calcolarne il rango.
	(b) Verificare che $\nu=0$ ma il sistema non è isostatico:
	classificarlo. (c) Determinare una base di $\Nc(\bm{A})$ e
	descrivere il meccanismo. (d) Determinare una base di
	$\Nc(\bm{A}^\top)$ e dare significato meccanico allo stato
	di autoequilibrio.
	
	\risultato (a) Le tre equazioni di vincolo $v(A)=v_A=\bar{v}_A$,
	$v(C)=v_A+\theta\,\ell=\bar{v}_C$,
	$v(B)=v_A+2\theta\,\ell=\bar{v}_B$ danno
	\[
	\bm{A}=
	\begin{pmatrix}
		0 & 1 & 0\\
		0 & 1 & \ell\\
		0 & 1 & 2\ell
	\end{pmatrix}.
	\]
	La prima colonna è nulla, quindi $\rg(\bm{A})\le 2$; il minore
	$2\times 2$ in alto a destra ha determinante $\ell\neq 0$,
	dunque $\rg(\bm{A})=2$.
	(b) $\nu=3n_c-m=3-3=0$, ma $g_\ell=n-r=1$ e
	$g_{\mathscr{i}}=m-r=1$: sistema \emph{labile-iperstatico}.
	La degenerazione è geometrica: i tre vincoli sono paralleli
	e non sufficienti a immobilizzare la trave nel piano
	(cfr.~\TABref{tab:classificazione_SCR}, riga
	"labile-iperstatico").
	(c) $\Nc(\bm{A})$ è generato da $\bm{u}_n=(1\;0\;0)^\top$,
	cioè da una traslazione orizzontale uniforme della trave (CR
	all'infinito sulla verticale): tutti i punti del corpo
	traslano della stessa quantità lungo $\ab_x$, e i tre carrelli
	verticali — non sentendo lo spostamento orizzontale — non si
	oppongono al moto.
	(d) $\Nc(\bm{A}^\top)$ è generato da
	$\bm{r}_n=\{1\;-2\;1\}^\top$: gli stati di autoequilibrio sono
	tutti e soli quelli della forma
	\[
	V_A=\lambda,\qquad V_C=-2\lambda,\qquad V_B=\lambda,
	\qquad \lambda\in\mathbb{R},
	\]
	ovvero terne di reazioni verticali con risultante
	$V_A+V_C+V_B=0$ e momento rispetto ad $A$
	$V_C\,\ell+V_B\,2\ell=0$.  Per uno stato di forze attive
	$\bm{f}$, la condizione di risolubilità statica è
	$\bm{U}_q^\top\bm{f}=X=0$: il sistema non sopporta carichi
	orizzontali, dualmente al fatto che il meccanismo $\bm{u}_n$
	del punto~(c) non produce spostamento differenziale fra i tre
	carrelli (gli spazi $\Nc(\bm{A})$ e $\Ic(\bm{A}^\top)$ sono
	ortogonali nel senso del lavoro virtuale).
\end{esercizio}

\begin{figure}[ht]
	\centering
	\begin{tikzpicture}[>=Stealth, line width=0.8pt, font=\small,
		nodo/.style={circle, fill=white, draw=black,
			inner sep=0pt, minimum size=5pt, line width=0.8pt},
		tacca/.style={thin, gray!70}]
		
		\def\L{6.0}
		\def\tw{0.22}\def\thA{0.38}\def\thR{0.32}
		\def\bw{0.16}\def\gw{0.30}\def\sr{0.08}
		
		\draw[line width=2.2pt] (0,0) -- (\L,0);
		\node[above=2pt] at (0,0)    {$A$};
		\node[above=2pt] at (\L/2,0) {$C$};
		\node[above=2pt] at (\L,0)   {$B$};
		
		%% cerniera in A
		\draw (0,0) -- ++(-\tw,-\thA) -- ++(2*\tw,0) -- cycle;
		\draw (-\gw,-\thA) -- (\gw,-\thA);
		\foreach \x in {-0.24,-0.12,0,0.12,0.24}
		\draw[thin] (\x,-\thA) -- ++(-0.09,-0.11);
		\node[nodo] at (0,0) {};
		
		%% carrello intermedio in C (mezzaluna)
		\draw (\L/2,0) -- ++(-\tw,-\thR) -- ++(2*\tw,0) -- cycle;
		\draw (\L/2-\gw,-\thR) -- (\L/2+\gw,-\thR);
		\draw (\L/2-\bw,-\thR-0.12) circle (2.4pt);
		\draw (\L/2+\bw,-\thR-0.12) circle (2.4pt);
		\draw (\L/2-\gw,-\thR-0.24) -- (\L/2+\gw,-\thR-0.24);
		\foreach \x in {-0.24,-0.12,0,0.12,0.24}
		\draw[thin] (\L/2+\x,-\thR-0.24) -- ++(-0.09,-0.11);
		\fill[white] (\L/2-\sr,0) arc(180:360:\sr) -- cycle;
		\draw (\L/2-\sr,0) arc(180:360:\sr);
		
		%% carrello terminale in B (perno)
		\draw (\L,0) -- ++(-\tw,-\thR) -- ++(2*\tw,0) -- cycle;
		\draw (\L-\gw,-\thR) -- (\L+\gw,-\thR);
		\draw (\L-\bw,-\thR-0.12) circle (2.4pt);
		\draw (\L+\bw,-\thR-0.12) circle (2.4pt);
		\draw (\L-\gw,-\thR-0.24) -- (\L+\gw,-\thR-0.24);
		\foreach \x in {-0.24,-0.12,0,0.12,0.24}
		\draw[thin] (\L+\x,-\thR-0.24) -- ++(-0.09,-0.11);
		\node[nodo] at (\L,0) {};
		
		\draw[tacca] (0,    -1.05) -- (0,    -0.85);
		\draw[tacca] (\L/2, -1.05) -- (\L/2, -0.85);
		\draw[tacca] (\L,   -1.05) -- (\L,   -0.85);
		\draw[<->, thin] (0,-0.95)    -- (\L/2,-0.95)
		node[midway, above=2pt] {$\ell$};
		\draw[<->, thin] (\L/2,-0.95) -- (\L,-0.95)
		node[midway, above=2pt] {$\ell$};
		
	\end{tikzpicture}
	\caption{Trave continua su tre appoggi
		(\ESEref{ex:C05_trave_continua}): cerniera in $A$, carrelli
		verticali in $C$ e $B$. Sistema iperstatico di grado uno.}
	\label{fig:C05_trave_continua}
\end{figure}

\begin{esercizio}[Trave continua su tre appoggi — sistema iperstatico]
	\label{ex:C05_trave_continua}
	
	Una trave rigida singola di lunghezza $2\ell$ è vincolata da una
	cerniera in $A=(0,0)$ e da due carrelli verticali in $C=(\ell,0)$
	e $B=(2\ell,0)$ (\FIGref{fig:C05_trave_continua}). Polo $O\equiv A$.
	(a) Scrivere $\bm{A}$ e calcolarne il rango.
	(b) Classificare il sistema. (c) Determinare la base
	$\bm{R}_\chi$ di $\Nc(\bm{A}^\top)$ e darne l'interpretazione
	meccanica. (d) Scrivere la condizione di compatibilità per il
	problema cinematico in funzione di un assegnato vettore di
	cedimenti $\bm{s}=(\bar{u}_A\;\bar{v}_A\;\bar{v}_C\;\bar{v}_B)^\top$.
	
	\risultato (a) Le quattro equazioni di vincolo
	$u_A=\bar{u}_A$, $v_A=\bar{v}_A$, $v_A+\theta\ell=\bar{v}_C$,
	$v_A+2\theta\ell=\bar{v}_B$ danno
	\[
	\bm{A}=
	\begin{pmatrix}
		1 & 0 & 0\\
		0 & 1 & 0\\
		0 & 1 & \ell\\
		0 & 1 & 2\ell
	\end{pmatrix}
	\in\mathbb{R}^{4\times 3}.
	\]
	Il minore $3\times 3$ delle prime tre righe ha determinante
	$\ell\neq 0$, quindi $\rg(\bm{A})=3$.
	(b) $\nu=n-m=-1$; $g_\ell=n-r=0$, $g_{\mathscr{i}}=m-r=1$:
	sistema iperstatico di grado uno.
	(c) Le quattro reazioni del sistema sono $H_A$, $V_A$ (cerniera
	in $A$), $V_C$ (carrello in $C$), $V_B$ (carrello in $B$). Il
	sistema $\bm{A}^\top\bm{r}=\bm{0}$ dà
	\[
	H_A = 0,
	\qquad V_A + V_C + V_B = 0,
	\qquad V_C\,\ell + V_B\,2\ell = 0.
	\]
	Le soluzioni sono tutte e sole della forma
	\[
	H_A = 0,
	\quad V_A = \lambda,
	\quad V_C = -2\lambda,
	\quad V_B = \lambda,
	\qquad \lambda\in\mathbb{R},
	\]
	e $\Nc(\bm{A}^\top)$ è generato da
	$\bm{R}_\chi=\{0\;1\;-2\;1\}^\top$. Lo stesso rapporto $1:-2:1$
	appare come meccanismo nel sistema
	dell'\ESEref{ex:C05_tre_carrelli} e qui come stato di
	autoequilibrio: è la stessa relazione algebrica fra le tre quote
	$A$, $C$, $B$ --- letta come spostamenti verticali in un caso,
	come reazioni nell'altro --- che esprime in modo concreto il
	dualismo $\Nc(\bm{A}) \leftrightarrow \Nc(\bm{A}^\top)$.
	(d) La condizione $\bm{R}_\chi^\top\bm{s}=0$ si scrive
	$\bar{v}_A - 2\,\bar{v}_C + \bar{v}_B = 0$: i tre cedimenti
	verticali non possono essere assegnati indipendentemente, ma
	devono soddisfare la relazione lineare che esprime la rigidità
	della trave (differenza seconda discreta nulla: i tre punti
	$A$, $C$, $B$ restano allineati anche dopo il cedimento).
\end{esercizio}